\chapter{Conway--Norton for the automorphism group of the moonshine module}\label{ch:20}

\providecommand{\chtwentylean}[1]{\texttt{\directlua{ local s = "\luaescapestring{\detokenize{#1}}" local first = true for c in s:gmatch(utf8.charpattern) do if (not first) and c:match("^[A-Z]$") then tex.sprint("\string\\allowbreak{}") end tex.sprint(-2, c) if c == "_" or c == "." then tex.sprint("\string\\allowbreak{}") end first = false end}}}
\providecommand{\chtwentyfile}[1]{\texttt{\directlua{ local s = "\luaescapestring{\detokenize{#1}}" local first = true for c in s:gmatch(utf8.charpattern) do if (not first) and c:match("^[A-Z]$") then tex.sprint("\string\\allowbreak{}") end tex.sprint(-2, c) first = false end}}}
\providecommand{\chtwentyVL}{V_\Leech^{+}}
\providecommand{\chtwentyTp}{T^{+}}
\providecommand{\chtwentyvac}{\mathbf 1}
\providecommand{\chtwentym}{\mathfrak m}
\providecommand{\chtwentyg}{\mathfrak g}
\providecommand{\chtwentyLf}{\Leech_4}
\chaptermark{Conway--Norton for Aut of the moonshine module}
\begingroup\setlength{\emergencystretch}{3em}\sloppy

This chapter documents the release of the library named in the front matter; every Lean quotation is checked by
\texttt{scripts/textbook/sync.py} to occur verbatim (up to whitespace and comments) in the source files of that
release.  (The axiom audit of the \texttt{Aut*} files cited in the status boxes is
\leanpath{VOA/Moonshine/MoonshineLRLiftAudit.lean}.)  This chapter supplies the last two inputs of the library's
Conway--Norton theorem and assembles it.

\paragraph{What exactly is proved.}  The final theorem is the \emph{qualitative genus-zero form} of the Conway--Norton
conjecture, for the full automorphism group of the library's $\Vnat$: for every $g\in\Aut(\Vnat)$ the McKay--Thompson
series $T_g$ is the normalised Hauptmodul of \emph{some} genus-zero group $\Gamma_g$ containing some $\Gamma_0(N)$ and
commensurable with $\mathrm{SL}_2(\Z)$.  It does \emph{not} identify $\Gamma_g$ with the group in Conway and Norton's
table, does not require $\Gamma_g$ to lie in the normaliser of $\Gamma_0(N)$, and does not relate $N$ to the order of $g$;
so it is weaker than the table-level conjecture (Chapter~\ref{ch:16}, \S\ref{ch16:sec:fidelity}).  It is a statement about
$\Aut(\Vnat)$, whatever that group is; its link with the Monster (AxiomMath's group $\mathrm{AutO}$) is the subject of
Chapters~\ref{ch:21}--\ref{ch:22}, and the identification $\Aut(\Vnat)\cong\mathrm{AutO}$ there is conditional on the
named hypothesis \lean{GmatEqAutO}.  The machine evidence is targeted and umbrella builds with \texttt{\#print axioms}
reports (listed in the status boxes); the book's authors did not re-run them.

\paragraph{Where the chapter sits.}  Chapter~\ref{ch:16} stated the
conjecture as the Lean proposition \lean{ConwayNortonConjecture}; Chapters~\ref{ch:17}--\ref{ch:19} built the Monster
Lie algebra $\chtwentym$, its Borcherds--Kac--Moody structure \emph{given} that $\chtwentym$ is generated by its
simple root spaces, the twisted denominator identities, complete replicability and the case-free genus-zero
argument.  At the end of Chapter~\ref{ch:19} the conjecture is reduced, in Lean, to exactly two inputs:
\begin{enumerate}[label=(\Roman*)]
\item \textbf{generation}: the Monster Lie algebra $\chtwentym$ is generated by its Cartan subalgebra and its
  simple root spaces (\chtwentylean{MonsterLieSimpleRootGeneration});
\item \textbf{finite order}: every automorphism of the moonshine vertex operator algebra $\Vnat$ has finite order
  (\lean{∀ g : AutVnat, IsOfFinOrder g}).
\end{enumerate}
This chapter proves both, with no hypotheses, and assembles the final theorem
\lean{M8.conwayNortonConjecture_holds : ConwayNortonConjecture}, which is re-exported under its official name
\lean{ConwayNorton.conwayNorton} (Chapter~\ref{ch:16}).  Along the way it documents the files on the
automorphism group $\Aut(\Vnat)$ itself: that every automorphism is determined by its action on the Griess
algebra $\Vnat_2$, the triality automorphism $\sigma$, and the theorem that $\Aut(\Vnat)$ is a \emph{finite} group.

\paragraph{The results in one paragraph.}  With no hypotheses and only Lean's three standard axioms, the library
proves: (a) restriction to the Griess algebra, $\Aut(\Vnat)\to\Aut(\Vnat_2,\times,\langle\cdot,\cdot\rangle)$, is
injective (\chtwentylean{vnatGriessRestrict_injective}); (b) the triality involution $\sigma^\natural$ of
Chapters~\ref{ch:12}--\ref{ch:14} is an automorphism of $\Vnat$ (\lean{sigmaAut}); (c) every automorphism of $\Vnat$
has finite order (\lean{autVnat_isOfFinOrder}); (d) $\Aut(\Vnat)$ is finite (\lean{finite_autVnat}); (e) the
Monster Lie algebra is generated by its simple root spaces (\chtwentylean{monsterLieSimpleRootGeneration}); and (f) the
genus-zero form of the Conway--Norton conjecture, as just described, holds for every automorphism of $\Vnat$
(\chtwentylean{conwayNortonConjecture_holds}).

\paragraph{Finite order and finiteness.}  (d) implies (c): in a finite group every element has finite order (Lagrange).
The library nevertheless proves (c) directly, and it is (c), not (d), that the final assembly uses
(\chtwentylean{conwayNortonConjecture_holds} $:=$ \chtwentylean{conwayNorton_of_isOfFinOrder} \lean{autVnat_isOfFinOrder};
\lean{finite_autVnat} appears in no proof of the Conway--Norton chain).  The reason is the order in which the results were
obtained: finite order was the last missing input of the conjecture and was proved first, by the periodicity lemma of
\S\ref{ch20:sec:periodic}; finiteness was proved afterwards, for the identification with the Monster, by a separate
algebraic-geometry lemma (\S\ref{ch20:sec:finite}) which does not use the periodicity lemma.  The two proofs share their
VOA ingredients (Ising vectors, rigidity, the faithfulness lemmas, $\sigma$).  A reorganised library could prove (d) first and
obtain (c) from it; the book documents the library as it is.
Statements (c) and (d) are proved \emph{without any group theory}: no fact about the Monster, its order, its
classes or its subgroups is used.  They rest on explicit computations with \emph{Ising vectors} in $\Vnat_2$ and on
two pieces of finite-dimensional algebra (a periodicity lemma and a finiteness lemma for ``rigid idempotents'' of a
commutative non-associative algebra).  Statement (e) is derived from the product formula for $J(p)-J(q)$ of
Chapter~\ref{ch:19} by a dimension count (``the sandwich''); there is no Lie-theoretic shortcut, for the reason given in
the warning at the end of \S\ref{ch20:sec:genmonster}.

\paragraph{What this chapter does.}  It
\begin{enumerate}[label=(\arabic*)]
\item recalls the automorphism group \lean{AutVnat} and its action on the Griess algebra, and proves faithfulness
  (\S\ref{ch20:sec:aut});
\item presents the triality automorphism $\sigma$ (\S\ref{ch20:sec:sigma});
\item proves that every automorphism has finite order: the abstract periodicity lemma
  (\S\ref{ch20:sec:periodic}), the explicit Ising vectors and their rigidity (\S\ref{ch20:sec:ising}), and the
  assembly (\S\ref{ch20:sec:finorder});
\item proves that $\Aut(\Vnat)$ is finite (\S\ref{ch20:sec:finite});
\item proves that $\chtwentym$ is generated by its simple root spaces, and explains how this feeds the twisted
  denominator identities (\S\ref{ch20:sec:generation});
\item assembles the Conway--Norton theorem and relates it to \lean{conwayNorton} (\S\ref{ch20:sec:assembly});
\item accounts for all $33$ files assigned to the chapter (\S\ref{ch20:sec:files}), collects the recorded traps
  (\S\ref{ch20:sec:traps}), and ends with the referee's checklist and the references.
\end{enumerate}
As everywhere in this book, our concern is the \emph{meaning} of the Lean statements; proofs are explained
mathematically, and we do not describe how the Lean proof scripts are organised.  The most important fidelity checks
for a referee are those of Theorem~\ref{ch20:thm:faithful} (what ``an automorphism fixing $\Vnat_2$ is the
identity'' says), Theorem~\ref{ch20:thm:finorder} (finite order), Theorem~\ref{ch20:thm:generation} (generation)
and Theorem~\ref{ch20:thm:final} (the final assembly).

\paragraph{Notation used throughout.}  We use the notation of Chapters~\ref{ch:02}--\ref{ch:14} and~\ref{ch:16}.
\begin{itemize}
\item $\Leech$ is the Leech lattice (\lean{leech}, Chapter~\ref{ch:02}), with form $\langle\cdot,\cdot\rangle$
  (\lean{leech.form}), even, positive definite, unimodular, with no vectors of norm $2$.  We write
  $\chtwentyLf=\{\alpha\in\Leech:\langle\alpha,\alpha\rangle=4\}$ for the \emph{minimal vectors}; there are $196560$
  of them, a finite set (\lean{normFour_finite}, Chapter~\ref{ch:04}).
\item $V_\Leech=M(1)\otimes\C[\Leech]$ is the lattice vertex operator algebra (Chapter~\ref{ch:03}); $h_a(-k-1)$,
  $a=1,\dots,24$, are the oscillators in an orthonormal frame; the coordinates of $\gamma\in\Leech$ in that frame are
  $\gamma_a$ (\lean{chargeCoord leech γ a}), so that $\sum_a\gamma_a\delta_a=\langle\gamma,\delta\rangle$.
\item $\varepsilon(\alpha,\beta)\in\{\pm1\}$ is the lattice cocycle (\lean{latticeCocycle leech α β},
  Chapter~\ref{ch:03}); it is bimultiplicative and $\varepsilon(\alpha,\beta)^2=1$.  Textbooks often normalise it
  so that $\varepsilon(\alpha,\alpha)=(-1)^{\langle\alpha,\alpha\rangle/2}$, which would make $\kappa_\alpha:=\varepsilon(\alpha,-\alpha)=+1$ for
  minimal $\alpha$.  The library's cocycle is only a bimultiplicative sign function built from a chosen basis, without that
  normalisation (any such cocycle gives an isomorphic lattice VOA, so nothing is lost); its signs are \emph{never evaluated},
  and in particular $\kappa_\alpha$ is not known to be $+1$.  The proofs below keep $\kappa_\alpha$ symbolic.
\item $\theta$ is the lift of $-1$; $\chtwentyVL$ (\lean{VPlus}, \lean{leechVOAPlus}) is the fixed-point
  subalgebra (Chapter~\ref{ch:04}); $x_\beta=e^\beta+e^{-\beta}\in\chtwentyVL$ (\lean{frameX0P β},
  Chapter~\ref{ch:11}); $x_0=2\cdot\chtwentyvac$; $x_{-\beta}=x_\beta$.
\item $q_{ab}=h_a(-1)h_b(-1)\chtwentyvac\in(\chtwentyVL)_2$ (\lean{qP a b}, Chapter~\ref{ch:13}); for
  $u,v\in\mathfrak h=\C\otimes\Leech$ we write $Q(u,v)=u(-1)v(-1)\chtwentyvac=\sum_{a,b}u_av_b\,q_{ab}$; on lattice
  vectors $Q(\alpha,\beta)$ is \lean{dV α β} (Chapter~\ref{ch:13}).
\item $\chtwentyTp$ (\lean{TP}) is the $\theta$-fixed part of the twisted module $V_\Leech^T=S[\theta]\otimes T$
  (Chapters~\ref{ch:06}--\ref{ch:08}), with its $\chtwentyVL$-module structure \lean{plusModule}.
\item $\Vnat=\chtwentyVL\oplus\chtwentyTp$ is the carrier \lean{Vnat} (pairs $(x,t)$), and \lean{VnatVOA} the
  vertex operator algebra on it (Chapter~\ref{ch:16}, Definition~\ref{ch16:def:vnatvoa}).  For $v\in\Vnat_2$ the
  \emph{zero mode} is $o(v)=v_{(1)}$; it preserves every weight space.
\item The \emph{Griess product} on $\Vnat_2$ is $u\times v=u_{(1)}v$ (\lean{griessMul}, Chapter~\ref{ch:L04}).  It is
  commutative because $\Vnat_1=0$.
\end{itemize}
\paragraph{A warning about letters.}  Some letters carry more than one meaning, always distinguished by decoration or
context: $e^\beta$ is a lattice exponential, $e_{\alpha,s}$ an Ising vector (\S\ref{ch20:sec:ising}), $e_\alpha$ an operator on
the module $T$, and a bare $e$ a generic idempotent (\S\ref{ch20:sec:periodic}) or a real root vector
(\S\ref{ch20:sec:generation}); $x_\beta$ is $e^\beta+e^{-\beta}$, while $x$ without subscript is the twisted generator of
$\chtwentyTp$; $\sigma$ is the triality automorphism except in \S\ref{ch20:sec:generation}, where it denotes an arbitrary
finite-order automorphism; $\Gamma$ is a degree in the PBW counting, $\Gamma_g$ a modular group; $\mathfrak m$ is the Monster
Lie algebra, $\mathfrak m_p$ a maximal ideal; $R(\Vnat)$ is the set of rigid idempotents, $R$ alone the radical of a form.

%======================================================================
\section{The automorphism group and the Griess algebra}\label{ch20:sec:aut}
%======================================================================

\subsection{The object}

The vertex operator algebra \lean{VnatVOA} and its automorphism group \lean{AutVnat} were defined in
Chapter~\ref{ch:16} (Definition~\ref{ch16:def:vnatvoa}); the file that defines them,
\leanpath{VOA/Moonshine/AutVnat.lean}, belongs to the present chapter.  We recall the definitions and add the facts
from that file that were not needed there.

\begin{definition}[The moonshine module as a named VOA; its automorphisms]\label{ch20:def:vnat}
$\Vnat$ is the vertex operator algebra on $\chtwentyVL\oplus\chtwentyTp$ with the block vertex operator, vacuum and
conformal vector of Chapters~\ref{ch:10}--\ref{ch:14}.  An automorphism of $\Vnat$ is a $\C$-linear bijection
$g:\Vnat\to\Vnat$ with $g(u_{(n)}v)=(gu)_{(n)}(gv)$ for all $u,v\in\Vnat$ and \emph{all} $n\in\Z$,
$g\chtwentyvac=\chtwentyvac$ and $g\omega=\omega$.  $\Aut(\Vnat)$ is the group of automorphisms.
\end{definition}
\begin{leancode}
theorem vnat_twistedLocality (t t' : TP) : Local (YnatT leechWbar t) (YnatT leechWbar t') := by

noncomputable def VnatVOA : VertexOperatorAlgebra Vnat := vnatVOA vnat_twistedLocality

theorem VnatVOA_isSimple : VnatVOA.IsSimple :=
  SimpleNat.vnatVOA_isSimple vnat_twistedLocality

abbrev AutVnat : Type := Aut VnatVOA

theorem VnatVOA_vacuum_ne_zero : VnatVOA.vacuum ≠ 0 :=
  SimpleNat.vnatVOA_vacuum_ne_zero vnat_twistedLocality
\end{leancode}
(\leanpath{VOA/Moonshine/AutVnat.lean}, namespace \lean{VOA.Moonshine.AutMonster}.)

\begin{dictionary}
\lean{vnat_twistedLocality} is the twisted$\times$twisted locality of the block vertex operator, extracted from the
existential theorem \lean{moonshine} of Chapter~\ref{ch:14} (its proof goes through the $\sigma$-key
\lean{sigmaKey_Dnat} of the standard triality data \lean{Dnat}; Chapter~\ref{ch:14}).  It is the only hypothesis of
the constructor \lean{vnatVOA} of Chapter~\ref{ch:10}, so \lean{VnatVOA} is a definition with \emph{no} hypotheses,
and it is the very object of \lean{moonshine} (the theorems \lean{VnatVOA_Y}, \lean{VnatVOA_vacuum},
\lean{VnatVOA_conformalVector}, \lean{VnatVOA_grading} of the same file identify its data by \lean{rfl}).
\lean{VnatVOA.IsSimple} is simplicity in the sense of Chapter~\ref{ch:10}; \lean{VnatVOA.vacuum ≠ 0} is needed to
define the Griess form below.  \lean{Aut VnatVOA} is the structure \lean{VOA.Aut} quoted in Chapter~\ref{ch:16}: a
linear equivalence with the three fields \lean{map_smul'} ($g(u_{(n)}v)=(gu)_{(n)}(gv)$ for every $n\in\Z$),
\lean{map_vacuum}, \lean{map_omega}.  So an automorphism in this library preserves the vacuum, the conformal vector and
\emph{all} modes, which is the textbook definition (FLM, Gannon \S5.2).  Because $g\omega=\omega$, $g$ commutes
with every $L_n=\omega_{(n+1)}$ and maps each $\Vnat_n$ onto itself (\lean{GriessAut.aut_map_grading}).
\end{dictionary}
\begin{physics}
$\Vnat$ is the chiral algebra of the holomorphic $c=24$ ``Monster'' conformal field theory; an automorphism is a
\emph{global internal symmetry}: an invertible operator on the space of states that fixes the vacuum, commutes with
the stress tensor (so it preserves energies) and intertwines every operator product expansion
$u(z)v(w)\sim\sum_n (u_{(n)}v)(w)/(z-w)^{n+1}$.  The requirement for \emph{all} $n$, singular and regular terms alike,
is what makes $g$ a symmetry of the full operator algebra rather than of its leading singularities.
\end{physics}
\begin{status}
Definitions and theorems with no hypotheses.  \lean{VnatVOA}, \lean{AutVnat} are used by every final theorem of the
part, whose \texttt{\#print axioms} is \lean{[propext, Classical.choice, Quot.sound]} (e.g.\ \lean{conwayNorton} in
\leanpath{VOA/Moonshine/MoonshineCNFinalAudit.lean}); they are covered by those audits.
\end{status}

\subsection{The Griess algebra and the restriction map}

\begin{definition}[Griess algebra, its automorphism group, restriction]\label{ch20:def:griessaut}
Let $W$ be a vertex operator algebra with $\chtwentyvac\neq0$.  On $W_2$ put the product $u\times v=u_{(1)}v$ and the
\emph{intrinsic form} $\langle u,v\rangle_3$ defined by $u_{(3)}v=\langle u,v\rangle_3\,\chtwentyvac$ (note
$u_{(3)}v\in W_0$).  $\mathrm{GriessAut}(W)$ is the group of linear automorphisms of $W_2$ that preserve $\times$ and
$\langle\cdot,\cdot\rangle_3$.  Restriction $g\mapsto g|_{W_2}$ is a group homomorphism
$\Aut(W)\to\mathrm{GriessAut}(W)$.  $W$ is \emph{generated by $W_2$} if every subspace containing $W_2$ and closed
under all $n$-th products is $W$.  For $W=\Vnat$ these are \lean{GriessAutVnat} and \lean{vnatGriessRestrict}.
\end{definition}
\begin{leancode}
noncomputable def griessAut (hvac : W.vacuum ≠ 0) : Subgroup (W.grading 2 ≃ₗ[ℂ] W.grading 2) where
  carrier := {g | (∀ u v, g (W.griessMul u v) = W.griessMul (g u) (g v)) ∧
    ∀ u v, griessForm hvac (g u) (g v) = griessForm hvac u v}

noncomputable def griessRestrict (hvac : W.vacuum ≠ 0) : Aut W →* griessAut hvac where
  toFun α := ⟨autRestrict α 2, fun u v => griessMul_aut α u v, fun u v => griessForm_aut hvac α u v⟩

def GeneratedByWeightTwo (W : VertexOperatorAlgebra 𝒱) : Prop :=
  ∀ S : Submodule ℂ 𝒱, W.grading 2 ≤ S → (∀ u ∈ S, ∀ v ∈ S, ∀ n : ℤ, W.smul u n v ∈ S) → S = ⊤

noncomputable def griessForm (hvac : W.vacuum ≠ 0) (u v : W.grading 2) : ℂ :=
  RigidBridge.vacCoeff W hvac (griessThree u v)

theorem griessForm_spec (hvac : W.vacuum ≠ 0) (u v : W.grading 2) :
    W.smul (u : 𝒱) 3 (v : 𝒱) = griessForm hvac u v • W.vacuum :=
\end{leancode}
(\leanpath{VOA/Core/GriessAutomorphism.lean}, namespace \lean{VOA.GriessAut}; the group axioms of \lean{griessAut}
and the monoid-homomorphism fields of \lean{griessRestrict} are omitted.)
\begin{leancode}
theorem vnat_griessMul_comm (u v : VnatVOA.grading 2) :
    VnatVOA.griessMul u v = VnatVOA.griessMul v u :=

noncomputable abbrev GriessAutVnat := GriessAut.griessAut VnatVOA_vacuum_ne_zero

noncomputable def vnatGriessRestrict : AutVnat →* GriessAutVnat :=
  GriessAut.griessRestrict VnatVOA_vacuum_ne_zero
\end{leancode}
(\leanpath{VOA/Moonshine/AutVnat.lean}.)
\begin{dictionary}
\lean{W.grading 2 ≃ₗ[ℂ] W.grading 2} is the type of $\C$-linear automorphisms of $W_2$; \lean{Subgroup} of it is a
subgroup of $\mathrm{GL}(W_2)$.  \lean{W.griessMul u v} is $u_{(1)}v$ (Chapter~\ref{ch:L04},
Definition~\ref{l04:def:griess}).  \lean{griessForm hvac u v} is the coefficient of $\chtwentyvac$ in $u_{(3)}v$
(\lean{griessForm_spec} states $u_{(3)}v=\langle u,v\rangle_3\chtwentyvac$; \lean{RigidBridge.vacCoeff} reads off the
coefficient, which needs $W_0=\C\chtwentyvac$ \emph{and} $\chtwentyvac\neq0$, whence \lean{hvac}).  Defining the form
intrinsically, through $u_{(3)}v$, means that its invariance under automorphisms is automatic
($g$ fixes $\chtwentyvac$ and intertwines the $(3)$-product); no uniqueness theorem for invariant forms is needed.
\lean{autRestrict α 2} is the restriction of $\alpha$ to $W_2$ (Chapter~\ref{ch:16}).  \lean{→*} is a monoid (hence
group) homomorphism.  \lean{GeneratedByWeightTwo W}: every subspace $S$ with $W_2\le S$ and $u_{(n)}v\in S$ for
$u,v\in S$, $n\in\Z$, is all of $W$ (\lean{⊤}); this is ``$W$ is generated by $W_2$ as a vertex algebra'' (note that
$\chtwentyvac\in S$ is not required: $\chtwentyvac=(2/c)\,\omega_{(3)}\omega$ when $c\ne0$).
\lean{vnat_griessMul_comm} is commutativity, from $\Vnat_1=0$ (\chtwentylean{griessMul_comm_of_grading1_eq_bot}: skew-symmetry
gives $u_{(1)}v=v_{(1)}u+\sum_{k\ge1}\pm\tfrac1{k!}L_{-1}^k\,v_{(1+k)}u$, and $v_{(2)}u\in\Vnat_1=0$, $v_{(3)}u\in\Vnat_0=\C\chtwentyvac$
is killed by $L_{-1}$, $v_{(m)}u=0$ for $m\ge4$ by weight).
\textbf{Could this be vacuous?}  If $\Vnat_0$ were not spanned by $\chtwentyvac$ the form would be ill-defined; it is
($\Vnat$ is of CFT type).  \lean{GriessAutVnat} is a subgroup of $\mathrm{GL}_{196884}(\C)$; nothing here asserts
what it is.
\end{dictionary}
\begin{physics}
$\Vnat_2$ is the space of weight-$2$ fields of the Monster CFT: the stress tensor $\omega$ and $196883$ Virasoro
primaries of dimension $2$.  The Griess product $u\times v$ is the coefficient of $(z-w)^{-2}$ in the OPE $u(z)v(w)$,
and $\langle u,v\rangle_3$ the coefficient of $(z-w)^{-4}$, i.e.\ the two-point function.  A global symmetry acts on
the weight-two fields preserving these OPE coefficients, hence by an automorphism of the Griess algebra; this is the
``Griess algebra'' on which Griess built the Monster.
\end{physics}
\begin{status}
Definitions; no hypotheses beyond \lean{hvac}, supplied for $\Vnat$ by \lean{VnatVOA_vacuum_ne_zero}.  The core file
\leanpath{VOA/Core/GriessAutomorphism.lean} belongs to Chapter~\ref{ch:L03} (assignment rule \texttt{VOA/Core/*.lean});
it is quoted here because the statements of this section use it.
\end{status}

\begin{theorem}[Faithfulness from generation]\label{ch20:thm:faithgen}
If $\Vnat$ is generated by $\Vnat_2$, then restriction $\Aut(\Vnat)\to\mathrm{GriessAut}(\Vnat)$ is injective.
\end{theorem}
\begin{leancode}
theorem vnatGriessRestrict_injective_of_gen (hgen : GriessAut.GeneratedByWeightTwo VnatVOA) :
    Function.Injective vnatGriessRestrict :=
  GriessAut.griessRestrict_injective _ hgen
\end{leancode}
(\leanpath{VOA/Moonshine/AutVnat.lean}.)
\begin{dictionary}
A conditional statement: the hypothesis \lean{hgen} is generation by the weight-two space, which the library does
\emph{not} prove for $\Vnat$.  The general theorem \lean{griessRestrict_injective} (same core file) is the
two-line argument below.  This declaration is superseded by the hypothesis-free
Theorem~\ref{ch20:thm:faithful}; it records the first route that was tried.
\end{dictionary}
\paragraph{Proof on paper.}  If $g|_{\Vnat_2}=1$, the fixed space $\{v:gv=v\}$ contains $\Vnat_2$ and is closed under
all products ($g(u_{(n)}v)=(gu)_{(n)}(gv)$), so by generation it is everything. \qed
\begin{physics}
``The OPE closure of the weight-two fields is the whole theory'': a symmetry that acts trivially on the generating
fields acts trivially on everything.
\end{physics}
\begin{status}
\textbf{Conditional} on \lean{hgen} (an unproved hypothesis); printed in \leanpath{VOA/Moonshine/MoonshineLRLiftAudit.lean} (all theorems of the \texttt{Aut*} files; green at the release, standard axioms, no \lean{sorryAx}); not used by any final theorem.
\end{status}

\subsection{Faithfulness without generation}

The weight-two space does generate $\Vnat$ (FLM), but that theorem is not in the library.  The library proves
faithfulness directly.  \emph{Reading guide:} the faithfulness theorem itself is \emph{not} used by the Conway--Norton theorem;
what is reused (in Lemma~\ref{ch20:lem:fixising}) are Steps 1--6 of its proof, with the hypothesis ``$g$ fixes $\Vnat_2$''
weakened to ``$g$ fixes the $(q_{ab},0)$ and the $(x_\alpha,0)$, $\alpha$ minimal'' (Steps 1--6 for the untwisted block) and to
``$g$ fixes $\chtwentyVL\times0$ and one twisted vector'' (Step 6, twisted part).

\begin{theorem}[$\Aut(\Vnat)$ acts faithfully on the Griess algebra]\label{ch20:thm:faithful}
An automorphism of $\Vnat$ that fixes $\Vnat_2$ pointwise is the identity.  Consequently restriction
$\Aut(\Vnat)\to\mathrm{GriessAut}(\Vnat)$, $g\mapsto g|_{\Vnat_2}$, is an injective group homomorphism.
\end{theorem}
\begin{leancode}
theorem aut_eq_one_of_fix_frame (g : AutVnat) (hg : ∀ v ∈ VnatVOA.grading 2, g v = v)
    (hx : ∀ β : leech.carrier, g ((frameX0P β, 0) : Vnat) = ((frameX0P β, 0) : Vnat)) : g = 1 := by

theorem aut_eq_one_of_fix_weightTwo (g : AutVnat) (hg : ∀ v ∈ VnatVOA.grading 2, g v = v) :
    g = 1 := by

theorem vnatGriessRestrict_injective : Function.Injective vnatGriessRestrict := by
\end{leancode}
(\leanpath{VOA/Moonshine/AutFaithfulAssembly.lean}, namespace \lean{VOA.Moonshine.AutMonster}.)
\begin{dictionary}
\lean{g v} applies the automorphism (coercion to a function); \lean{∀ v ∈ VnatVOA.grading 2, g v = v} is
$g|_{\Vnat_2}=\mathrm{id}$; \lean{g = 1} is equality with the identity of the group \lean{AutVnat}.
\lean{((frameX0P β, 0) : Vnat)} is the vector $(x_\beta,0)\in\chtwentyVL\oplus\chtwentyTp$, $x_\beta=e^\beta+e^{-\beta}$.
\lean{Function.Injective vnatGriessRestrict} is injectivity of the homomorphism of
Definition~\ref{ch20:def:griessaut}; by \lean{injective_iff_map_eq_one} it is equivalent to the second statement.
The first statement is an intermediate step with an extra hypothesis \lean{hx}, which the second discharges.
\textbf{Strength.}  The second and third statements have no hypotheses; they are exactly the textbook statements
``$\Aut(\Vnat)$ acts faithfully on $\Vnat_2$'' and ``$\Aut(\Vnat)\hookrightarrow\Aut(\Vnat_2,\times,\langle,\rangle)$''.
They say nothing about the image (that is Chapter~\ref{ch:22}).  \textbf{Could it be vacuous?}  The statement ``$g$ fixes
$\Vnat_2\Rightarrow g=1$'' is a genuine constraint on every automorphism, whatever the size of the group: it says that the
$196884$-dimensional Griess algebra already determines $g$ on the infinite-dimensional $\Vnat$.  (That $\Aut(\Vnat)$ is
non-trivial is not needed here and is not stated as a separate Lean theorem in this chapter; Chapter~\ref{ch:22} embeds
AxiomMath's group $\mathrm{Gmat}$ in it.)
\end{dictionary}

\paragraph{Proof on paper.}  Let $g$ fix $\Vnat_2$ pointwise.  The proof has six steps; the first four show that
$g$ multiplies each $(x_\beta,0)$ by a scalar $c_\beta$, the fifth that $c_\beta=1$, and the last propagates to all of
$\Vnat$.

\emph{Step 1: $g$ commutes with the zero modes of the quadratic states.}  $q_{ab}=h_a(-1)h_b(-1)\chtwentyvac$ has
weight two, so $(q_{ab},0)\in\Vnat_2$ is fixed by $g$, and
$g\bigl(o(q_{ab})v\bigr)=o(gq_{ab})(gv)=o(q_{ab})(gv)$ for every $v$.  Hence $g$ maps joint eigenvectors of the
operators $o(q_{ab})$ to joint eigenvectors with the same eigenvalues.

\emph{Step 2: $(x_\beta,0)$ is a joint eigenvector with eigenvalues $\beta_a\beta_b$.}  On $V_\Leech$ the vertex
operator of $q_{ab}$ is the normal-ordered product ${:}h_a(z)h_b(z){:}$, whose zero mode is
\begin{equation*}
  o(q_{ab})=h_a(0)h_b(0)+\sum_{n>0}\bigl(h_a(-n)h_b(n)+h_b(-n)h_a(n)\bigr).
\end{equation*}
On $p\otimes e^\gamma$ (polynomial $p$ in the oscillators, charge $\gamma$) the momentum operator $h_a(0)$ acts by the
scalar $\gamma_a=\langle h_a,\gamma\rangle$, so the first term gives $\gamma_a\gamma_b$; and $h_a(n)$, $n>0$, acts on the Fock
factor only, as $n\,\partial/\partial X_{n-1,a}$.  Hence $o(q_{ab})$ acts by
\begin{equation}\label{ch20:eq:zeromode}
  o(q_{ab})(p\otimes e^\gamma)=\gamma_a\gamma_b\,p\otimes e^\gamma+N_{ab}(p)\otimes e^\gamma,\qquad
  N_{ab}=\sum_{k\ge0}(k+1)\Bigl(X_{k,a}\frac{\partial}{\partial X_{k,b}}+X_{k,b}\frac{\partial}{\partial X_{k,a}}\Bigr),
\end{equation}
where $X_{k,a}$ is the polynomial variable for $h_a(-k-1)$.  (In the library's construction $V_\Leech=M(1)\otimes\C[\Leech]$
the Heisenberg field on the Fock factor $M(1)$ has \emph{zero} zero-mode, and the momentum is carried separately by the
charge factor; that is why the Lean formula, Lemma~\ref{ch20:lem:zeromode}, has the scalar $\gamma_a\gamma_b$ plus a purely
Fock operator and no cross terms.)  Since
$N_{ab}$ kills constants, $o(q_{ab})e^{\pm\beta}=\beta_a\beta_b e^{\pm\beta}$, hence $o(q_{ab})x_\beta=\beta_a\beta_bx_\beta$.

\emph{Step 3: in $\chtwentyVL$ the only such joint eigenvectors are the multiples of $x_\beta$.}  This rests on a
lemma about polynomials (Lemma~\ref{ch20:lem:jointeigen} below): \emph{if $p\neq0$ is a joint eigenvector of all the
operators $N^\mu_{ab}=\sum_k\mu_k(X_{k,a}\partial_{k,b}+X_{k,b}\partial_{k,a})$, $a,b$ in a set of at least two
directions, with weights $\mu_k$ that are injective in $k$ and satisfy $\mu_k+\mu_{k'}\neq0$, then $p$ is constant and
all eigenvalues are $0$.}  Granting it, decompose $v=\sum_\gamma p_\gamma\otimes e^\gamma$ into charge sectors, which
$o(q_{ab})$ preserves.  If $o(q_{ab})v=\beta_a\beta_bv$ for all $a,b$, every nonzero $p_\gamma$ is a joint eigenvector
of $N_{ab}$ ($\mu_k=k+1$) with eigenvalue $\beta_a\beta_b-\gamma_a\gamma_b$; so $p_\gamma$ is constant and
$\gamma_a\gamma_b=\beta_a\beta_b$ for all $a,b$, i.e.\ $\gamma\gamma^{\mathsf T}=\beta\beta^{\mathsf T}$ as rank-one
matrices, whence $\gamma=\pm\beta$.  So $v=c_1e^\beta+c_2e^{-\beta}$; if moreover $v\in\chtwentyVL$, invariance under
$\theta$ (which exchanges $e^\beta$ and $e^{-\beta}$) gives $c_1=c_2$ when $\beta\ne0$, i.e.\ $v\in\C x_\beta$.

\emph{Step 4: no such eigenvector in the twisted block.}  Recall (Chapters~\ref{ch:06}--\ref{ch:08}) that the twisted module is
$V_\Leech^T=S[\theta]\otimes T$, where $S[\theta]$ is the Fock space of the half-integrally moded oscillators $h_a(-k-\tfrac12)$
and $T$ is the $2^{12}$-dimensional irreducible module of the extraspecial group $2^{1+24}$ (no momentum); $\chtwentyTp$ is its
$\theta$-fixed part, and writing $t=\sum_ip_i\otimes\tau_i$ in a basis $(\tau_i)$ of $T$, the $p_i\in S[\theta]$ are the ``Fock
components'' of $t$.  On $\chtwentyTp$ the zero mode of
$q_{ab}$ is $\bigl(N^{\mathrm{tw}}_{ab}+\tfrac18\delta_{ab}\bigr)\otimes1$, where $N^{\mathrm{tw}}_{ab}$ is the operator
of the lemma with half-integral weights $\mu_k=k+\tfrac12$ (the variables now stand for $h_a(-k-\tfrac12)$).  The
constant $\tfrac18\delta_{ab}$ is a normal-ordering constant of the twisted vertex operator; it is forced by the known
twisted ground-state energy: since $\omega=\tfrac12\sum_aq_{aa}$, summing gives $L_0=\sum h(-n-\tfrac12)h(n+\tfrac12)+\tfrac12\cdot24\cdot\tfrac18
=\dots+\tfrac{24}{16}$, the vacuum energy of $24$ $\Z_2$-twisted bosons (\S\ref{ch20:sec:techfaith}).  If $t\in\chtwentyTp$ satisfied $o(q_{ab})t=\beta_a\beta_bt$ for all $a,b$ and $t\neq0$,
each Fock component of $t$ would be a joint eigenvector of $N^{\mathrm{tw}}_{ab}$ with eigenvalue
$\beta_a\beta_b-\tfrac18\delta_{ab}$, so by the lemma $\beta_a\beta_b=\tfrac18\delta_{ab}$ for all $a,b$.  In rank
$\ge2$ that is impossible: for $a\ne b$, $\beta_a\beta_b=0$ while $(\beta_a\beta_b)^2=\beta_a^2\beta_b^2=\tfrac1{64}$.
Hence $t=0$.  Combining steps 1--4: $g(x_\beta,0)$ is a joint eigenvector with eigenvalues $\beta_a\beta_b$, its
twisted component vanishes and its untwisted component is a multiple of $x_\beta$; so $g(x_\beta,0)=c_\beta(x_\beta,0)$
for a unique scalar $c_\beta$ ($x_\beta\neq0$).

\emph{Step 5: $c_\beta=1$.}  For $\langle\alpha,\gamma\rangle<0$ the \emph{leading product} is
\begin{equation}\label{ch20:eq:leading}
  (x_\alpha)_{(-\langle\alpha,\gamma\rangle-1)}\,x_\gamma=\varepsilon(\alpha,\gamma)\,x_{\alpha+\gamma}
\end{equation}
(Lemma~\ref{ch20:lem:leading}).  If $\alpha$ is minimal, $x_\alpha$ has weight $\langle\alpha,\alpha\rangle/2=2$, so
$(x_\alpha,0)\in\Vnat_2$ is fixed by $g$; applying $g$ to \eqref{ch20:eq:leading} gives
$c_\gamma\,\varepsilon\,x_{\alpha+\gamma}=c_{\alpha+\gamma}\,\varepsilon\,x_{\alpha+\gamma}$, so $c_{\alpha+\gamma}=c_\gamma$.  If
$\langle\alpha,\gamma\rangle\ge0$, apply the same to $-\alpha$ and $\gamma+\alpha$:
$\langle-\alpha,\gamma+\alpha\rangle=-4-\langle\alpha,\gamma\rangle<0$, and $(-\alpha)+(\gamma+\alpha)=\gamma$.  Hence
$c_{\gamma\pm\alpha}=c_\gamma$ for every $\gamma\in\Leech$ and every minimal $\alpha$.  Since the Leech lattice is
generated as an abelian group by its minimal vectors (Theorem~\ref{ch20:thm:leechspan}), $c_\beta=c_0$ for all $\beta$;
and $c_0=1$ because $x_0=2\chtwentyvac$ is fixed.  (In detail: the two cases give $c_{\gamma+\alpha}=c_\gamma$ for all
$\gamma$ and minimal $\alpha$; replacing $\gamma$ by $\gamma-\alpha$ gives $c_{\gamma-\alpha}=c_\gamma$, and the induction principle of
Theorem~\ref{ch20:thm:leechspan} applies.)

\emph{Step 6: propagation.}  The subspace $\{x\in\chtwentyVL: g(x,0)=(x,0)\}$ contains $\chtwentyvac$ and every
$x_\beta$, and is closed under products; by the generation theorem of Chapter~\ref{ch:04} (the $x_\beta$, $\beta\in\Leech$,
generate $\chtwentyVL$; \lean{leechPlus_generated}) it is all of $\chtwentyVL$.  On the twisted side, $\chtwentyTp=T^+_0\oplus T^+_1$ splits into the two
\emph{frame sectors} of Chapter~\ref{ch:11} (eigenspaces of a Golay-frame symmetry), and each $T^+_s$ is spanned by the vectors $u_{(n)}x_s$, $u\in\chtwentyVL$, $n\in\Z$, for a weight-two
vector $x_s\in T^+_s$ (Chapter~\ref{ch:11}, \lean{span_smul_frameX0T}); $g$ fixes $x_s\in\Vnat_2$ and fixes $\chtwentyVL$,
hence fixes every $u_{(n)}x_s$, hence $T^+_0\oplus T^+_1=\chtwentyTp$.  So $g=1$. \qed

\begin{physics}
The proof is a piece of honest CFT bookkeeping.  The $300$ states $q_{ab}$ ($a\le b$) are the bilinears $\partial X^a\partial X^b$
of the $24$ free bosons; their zero modes measure the ``momentum tensor'' $p_ap_b$ of a state plus an oscillator
contribution.  A symmetry fixing all weight-two fields commutes with these charges, so it maps the momentum-$\pm\beta$
cosine state $\cos(\beta\cdot X)$ to a multiple of itself, and it cannot move it into the twisted sector, where the
oscillators are half-integrally moded and the zero-point shift $\tfrac18\delta_{ab}$ (the twisted ground-state
energy $24/16$ split over the $24$ directions) makes the eigenvalue pattern of a momentum state impossible.  The
leading OPE coefficient $\cos(\alpha\cdot X)\times\cos(\gamma\cdot X)\to\cos((\alpha+\gamma)\cdot X)$ then transports the
phases along the lattice, and the vacuum fixes the normalisation.
\end{physics}
\begin{status}
\textbf{Unconditional}; no hypotheses.  All three statements are printed in \leanpath{VOA/Moonshine/MoonshineLRLiftAudit.lean} (Chapter~\ref{ch:22}), which
prints every theorem of the \texttt{Aut*} files; that audit built green with all of its $997$ reports within
\lean{[propext, Classical.choice, Quot.sound]} and no \lean{sorryAx}.  \chtwentylean{vnatGriessRestrict_injective}
is used by the lifting construction of \leanpath{VOA/Moonshine/MoonshineCNMonsterActionLift.lean} and
\leanpath{VOA/Moonshine/MoonshineCNMonsterActionConj.lean} (Chapter~\ref{ch:21}; used in Chapter~\ref{ch:22}).  The finite-order theorem does \emph{not} use these three statements;
it uses their ingredients (\S\ref{ch20:sec:techfaith}) with weakened hypotheses (Lemma~\ref{ch20:lem:fixising}).  The book's
authors did not re-run the build.
\end{status}

\subsection{The ingredients, stated}\label{ch20:sec:techfaith}

We state the ingredients of the proof as they are in Lean, because \S\ref{ch20:sec:finorder} reuses them.

\begin{lemma}[Zero mode of $q_{ab}$ on a lattice VOA]\label{ch20:lem:zeromode}
For every even positive-definite lattice $L$, frame directions $a,b$, polynomial $q$ in the oscillators and charge
$\gamma\in L$, formula \eqref{ch20:eq:zeromode} holds, with $N_{ab}$ given on polynomials by the finite sum over any
level bound $K$ of $q$.
\end{lemma}
\begin{leancode}
noncomputable abbrev quadState (a b : Fin (rank L)) : LatModuleGen L :=
  wordStateLatGen L [(0, a), (0, b)] 0

theorem smul_quadState_one (hEven : L.Even) (hPos : L.PosDef) (a b : Fin (rank L))
    (q : FockN (rank L)) (γ : L.carrier) :
    (latticeVOAGen L hEven hPos).smul (quadState L a b) 1 (latStateGen L q γ)
      = latStateGen L
          ((chargeCoord L γ a * chargeCoord L γ b) • q
            + nprod (currentN_isField (rank L) 1 a) (currentN_isField (rank L) 1 b) (-2) q) γ := by
\end{leancode}
(\leanpath{VOA/Moonshine/AutZeroModes.lean}, namespace \lean{VOA.LatGen}.)
\begin{leancode}
theorem nprod_currentN_neg_two (a b : Fin r) (q : FockN r) (K : ℕ)
    (hK : ∀ v ∈ q.vars, v.1 < K) :
    nprod (currentN_isField r 1 a) (currentN_isField r 1 b) (-2) q
      = ∑ k ∈ Finset.range K, ((k : ℂ) + 1) •
          (X (k, a) * pderiv (k, b) q + X (k, b) * pderiv (k, a) q) := by
\end{leancode}
(\leanpath{VOA/Moonshine/AutFockEigen.lean}, namespace \lean{VOA.HeisenbergN}.)
\begin{dictionary}
\lean{wordStateLatGen L [(0, a), (0, b)] 0} is the state $h_a(-1)h_b(-1)\chtwentyvac\otimes e^0$ (level index $0$
stands for mode $-1$).  \lean{latStateGen L q γ} is $q\otimes e^\gamma$, $q$ in the Fock space \lean{FockN (rank L)}
$=\C[X_{k,a}]$ (Chapter~\ref{ch:03}).  \lean{(latticeVOAGen L hEven hPos).smul u 1 v} is $u_{(1)}v$, the zero mode of
a weight-two $u$.  \lean{chargeCoord L γ a} is $\gamma_a$.  \lean{nprod (currentN_isField …) (…) (-2)} is the
$(-2)$-normal-ordered product of the two Heisenberg currents acting on the Fock space, i.e.\ the oscillator part
$N_{ab}$; \lean{nprod_currentN_neg_two} identifies it with
$\sum_{k<K}(k+1)(X_{k,a}\partial_{k,b}+X_{k,b}\partial_{k,a})$ whenever every variable of $q$ has level $<K$
(\lean{q.vars} is the set of variables occurring in $q$).  The two statements together are exactly
\eqref{ch20:eq:zeromode}.
\end{dictionary}
\begin{physics}
$o({:}\partial X^a\partial X^b{:})=p_ap_b+\sum_{k>0}(\alpha^a_{-k}\alpha^b_k+\alpha^b_{-k}\alpha^a_k)$: the familiar
mode expansion of the free-boson stress-tensor components; for $a=b$ summed over $a$ it is $2L_0$.
\end{physics}
\begin{status}
Unconditional.  Printed in \leanpath{VOA/Moonshine/MoonshineLRLiftAudit.lean} (all theorems of the \texttt{Aut*} files; green at the release, standard axioms, no \lean{sorryAx}).
\end{status}

\begin{lemma}[The joint-eigenvector lemma]\label{ch20:lem:jointeigen}
Let $r$ have two distinct directions, $\mu:\N\to\C$ injective with $\mu_k+\mu_{k'}\neq0$ for all $k,k'$, and
$N^\mu_{ab}=\sum_{k<K}\mu_k(X_{k,a}\partial_{k,b}+X_{k,b}\partial_{k,a})$ on $\C[X_{k,a}]$.  If $q\neq0$ has all
variables of level $<K$ and $N^\mu_{ab}q=c_{ab}q$ for all $a,b$, then $c_{ab}=0$ for all $a,b$ and $q$ is constant.
Consequently (Lean: \lean{eigen_quadState_span}), in a lattice VOA $V_L$ of rank $\ge2$, a joint eigenvector of all
$o(q_{ab})$ with eigenvalues $\beta_a\beta_b$ lies in $\C e^\beta+\C e^{-\beta}$; and (Lean: \lean{eigen_plus}) in
$\chtwentyVL$ it lies in $\C x_\beta$.
\end{lemma}
\begin{leancode}
noncomputable def quadOp (μ : ℕ → ℂ) (K : ℕ) (a b : Fin r) (q : FockN r) : FockN r :=
  ∑ k ∈ Finset.range K, μ k • (X (k, a) * pderiv (k, b) q + X (k, b) * pderiv (k, a) q)

theorem quadOp_joint_eigen (μ : ℕ → ℂ) (hsum : ∀ k k', μ k + μ k' ≠ 0)
    (hinj : ∀ k k', μ k = μ k' → k = k') (hr : ∀ b : Fin r, ∃ a, a ≠ b) (K : ℕ)
    (q : FockN r) (hK : ∀ v ∈ q.vars, v.1 < K) (c : Fin r → Fin r → ℂ)
    (h : ∀ a b, quadOp μ K a b q = c a b • q) (hq : q ≠ 0) :
    (∀ a b, c a b = 0) ∧ q = MvPolynomial.C (coeff 0 q) := by
\end{leancode}
(\leanpath{VOA/Moonshine/AutFockEigen.lean}.)
\begin{leancode}
theorem eigen_quadState_span (hEven : L.Even) (hPos : L.PosDef)
    (hr : ∀ b : Fin (rank L), ∃ a, a ≠ b) (β : L.carrier) (v : LatModuleGen L)
    (h : ∀ a b, (latticeVOAGen L hEven hPos).smul (quadState L a b) 1 v
      = (chargeCoord L β a * chargeCoord L β b) • v) :
    ∃ c₁ c₂ : ℂ, v = c₁ • latChargeGen L β + c₂ • latChargeGen L (-β) := by
\end{leancode}
(\leanpath{VOA/Moonshine/AutChargeEigen.lean}, namespace \lean{VOA.Moonshine.AutMonster}.)
\begin{leancode}
theorem eigen_plus (β : leech.carrier) (x : VPlus)
    (h : ∀ a b, leechVOAPlus.smul (qP a b) 1 x = (chargeCoord leech β a * chargeCoord leech β b) • x) :
    ∃ c : ℂ, x = c • frameX0P β := by
\end{leancode}
(\leanpath{VOA/Moonshine/AutFaithfulVnat.lean}.)
\begin{dictionary}
\lean{hr : ∀ b, ∃ a, a ≠ b} says the index set has at least two elements.  \lean{MvPolynomial.C (coeff 0 q)} is the
constant polynomial with value the constant term of $q$, so the second conjunct is ``$q$ is constant''.
\lean{latChargeGen L β} is $e^\beta=\chtwentyvac\otimes e^\beta$.  All hypotheses are the ones of the paper lemma.  The
weights used are $\mu_k=k+1$ (untwisted) and $\mu_k=k+\tfrac12$ (twisted; \lean{halfWeight_add_ne_zero},
\lean{halfWeight_injective} check the two conditions).
\end{dictionary}
\paragraph{Proof on paper.}  Write $\mathrm{Wt}_a(s)=\sum_k\mu_k s(k,a)$ for a monomial $X^s$.  (i) $N^\mu_{aa}$ is
diagonal on monomials: $N^\mu_{aa}X^s=2\mathrm{Wt}_a(s)X^s$; so every monomial in the support of $q$ has
$2\mathrm{Wt}_a(s)=c_{aa}$.  (ii) For $a\neq b$, $N^\mu_{ab}$ sends $X^s$ to monomials whose $\mathrm{Wt}_a$ differs from
$\mathrm{Wt}_a(s)$ by $\pm\mu_k\ne0$; since $N^\mu_{ab}q=c_{ab}q$ has support with constant $\mathrm{Wt}_a$, we get
$N^\mu_{ab}q=0$ and $c_{ab}=0$.  (iii) If some monomial $X^s$ of $q$ contains $X_{k,b}$, choose $a\neq b$ and look at
the coefficient of $X^t$, $t=s-e_{k,b}+e_{k,a}$, in $N^\mu_{ab}q$: the term $\mu_k s(k,b)\,\mathrm{coeff}_s(q)$ is
nonzero, and any other monomial $s'$ of $q$ contributing to $X^t$ would have
$\mathrm{Wt}_a(s')-\mathrm{Wt}_a(s)\in\{\mu_k-\mu_{k'}\ (k'\ne k),\ \mu_k+\mu_{k'}\}$, which is nonzero by the two hypotheses on
$\mu$, contradicting (i).  So $N^\mu_{ab}q\ne0$, contradicting (ii).  Hence $q$ is constant, and then
$c_{aa}=2\mathrm{Wt}_a(0)=0$.  The VOA consequences were explained in Steps 2--3 above. \qed
\begin{physics}
A state that has a definite value of every bilinear charge $p_ap_b+(\text{oscillators})$ equal to that of a pure momentum
state must be a pure momentum state: oscillator excitations always shift some ``level-weighted occupation number''
$\mathrm{Wt}_a$ and are therefore visible to the off-diagonal charges.
\end{physics}
\begin{status}
Unconditional.  Printed in \leanpath{VOA/Moonshine/MoonshineLRLiftAudit.lean} (all theorems of the \texttt{Aut*} files; green at the release, standard axioms, no \lean{sorryAx}).
\end{status}

\begin{lemma}[Twisted zero modes and twisted exclusion]\label{ch20:lem:twisted}
On the twisted block, $(q_{ab},0)_{(1)}(0,t)=\bigl(0,((N^{\mathrm{tw}}_{ab}+\tfrac18\delta_{ab})\otimes1)t\bigr)$.  If
$\beta\in\C^{24}$ and $(u,0)_{(1)}(0,t)=\beta_a\beta_b(0,t)$ for all $a,b$ and all $u\in\chtwentyVL$ whose underlying state is
$h_a(-1)h_b(-1)\chtwentyvac$, then $t=0$.
\end{lemma}
\begin{leancode}
theorem VnatVOA_smul_quad_one_snd (a b : Fin (rank leech)) (u : SkewVP)
    (hu : (u : LatModuleGen leech) = LatGen.quadState leech a b) (t : TP) :
    (VnatVOA.smul ((u, 0) : Vnat) 1 ((0, t) : Vnat)).1 = 0 ∧
      (((VnatVOA.smul ((u, 0) : Vnat) 1 ((0, t) : Vnat)).2 : TP) : Tw)
        = (quadNT 24 (cc a) (cc b)
            + (if a = b then (1 / 8 : ℂ) else 0) • (1 : End ℂ (FockT 24))).rTensor leechT (t : Tw) := by
\end{leancode}
(\leanpath{VOA/Moonshine/AutZeroModesTwisted.lean}, namespace \lean{VOA.Moonshine.AutZeroTw}.)
\begin{leancode}
theorem VnatVOA_twisted_outer_eigen_eq_zero (β' : Fin (rank leech) → ℂ) (t : TP)
    (h : ∀ a b : Fin (rank leech), ∀ u : SkewVP,
      (u : LatModuleGen leech) = LatGen.quadState leech a b →
        VnatVOA.smul ((u, 0) : Vnat) 1 ((0, t) : Vnat) = (β' a * β' b) • ((0, t) : Vnat)) :
    t = 0 := by
\end{leancode}
(\leanpath{VOA/Moonshine/AutTwistedExclusion.lean}.)
\begin{dictionary}
\lean{Tw} is $V_\Leech^T=S[\theta]\otimes T$ with $S[\theta]=$ \lean{FockT 24} $=\C[X_{k,c}]$, $X_{k,c}\leftrightarrow
h_c(-k-\tfrac12)$, and \lean{leechT} the $2^{12}$-dimensional module $T$ of the extraspecial group; \lean{TP} is the
subspace $\chtwentyTp$, coerced into \lean{Tw} by \lean{(t : Tw)}.  \lean{quadNT 24 c d} is $N^{\mathrm{tw}}_{cd}$, and
\lean{quadNT_apply} (same file) shows it equals \lean{quadOp (fun k => k + 1/2)} on polynomials of bounded level.
\lean{(…).rTensor leechT} is $A\otimes\mathrm{id}_T$.  \lean{cc} is the identification \lean{Fin (rank leech) ≃ Fin 24}
(a cast, \lean{rank leech = 24}).  \lean{.1}, \lean{.2} are the two components of a vector of
$\Vnat=\chtwentyVL\times\chtwentyTp$.  In the exclusion statement \lean{β'} is an arbitrary complex vector (not
necessarily a lattice vector); the conclusion is $t=0$, \emph{including} the case $\beta'=0$.
\end{dictionary}
\paragraph{Proof on paper.}  The twisted vertex operator of $q_{ab}$ is the normal-ordered product of two twisted
Heisenberg fields corrected by the twisted product formula; its mode $0$ picks up the contractions
$h_a(\tfrac12)\cdot h_b(-\tfrac12)$ with binomial coefficients, which sum to $(-\tfrac14+\tfrac38)\delta_{ab}=\tfrac18\delta_{ab}$
(file docstring of \leanpath{AutZeroModesTwisted.lean}).  Consistency check: with $\omega=\tfrac12\sum_aq_{aa}$ this gives
$L_0=\sum h(-n-\tfrac12)h(n+\tfrac12)+\tfrac{24}{16}$ on $V_\Leech^T$, the familiar twisted vacuum energy.  Charge-zero
states act on $S[\theta]\otimes T$ through $S[\theta]$ alone.  The exclusion: decompose $t=\sum_ip_i\otimes\tau_i$ along a
basis $(\tau_i)$ of $T$; the operators $A\otimes1$ act componentwise, so each nonzero $p_i$ is a joint eigenvector of
$N^{\mathrm{tw}}_{ab}$ with eigenvalue $\beta'_a\beta'_b-\tfrac18\delta_{ab}$, and Lemma~\ref{ch20:lem:jointeigen} with
$\mu_k=k+\tfrac12$ and the rank-one argument of Step 4 finish. \qed
\begin{physics}
In the $\Z_2$-twisted sector of $24$ free bosons the oscillators are half-integrally moded and the ground state has
energy $24/16$; the twisted sector carries no momentum at all.  So no twisted state can have the ``momentum tensor''
$\beta_a\beta_b$ of an untwisted momentum state: the symmetry cannot mix a cosine operator with a twist field.
\end{physics}
\begin{status}
Unconditional.  Printed in \leanpath{VOA/Moonshine/MoonshineLRLiftAudit.lean} (all theorems of the \texttt{Aut*} files; green at the release, standard axioms, no \lean{sorryAx}).
\end{status}

\begin{lemma}[The leading product]\label{ch20:lem:leading}
For $\alpha,\gamma\in\Leech$ with $\langle\alpha,\gamma\rangle<0$, \eqref{ch20:eq:leading} holds in $\chtwentyVL$, and
hence $(x_\alpha,0)_{(-\langle\alpha,\gamma\rangle-1)}(x_\gamma,0)=\varepsilon(\alpha,\gamma)(x_{\alpha+\gamma},0)$ in $\Vnat$.
\end{lemma}
\begin{leancode}
theorem frameX0P_smul_frameX0P_of_neg (α γ : leech.carrier) (h : leech.form α γ < 0) :
    leechVOAPlus.smul (frameX0P α) (-(leech.form α γ) - 1) (frameX0P γ)
      = latticeCocycle leech α γ • frameX0P (α + γ) := by
\end{leancode}
(\leanpath{VOA/Moonshine/AutFaithful.lean}.)
\begin{leancode}
theorem VnatVOA_smul_xN_of_neg (α γ : leech.carrier) (h : leech.form α γ < 0) :
    VnatVOA.smul ((frameX0P α, 0) : Vnat) (-(leech.form α γ) - 1) ((frameX0P γ, 0) : Vnat)
      = latticeCocycle leech α γ • ((frameX0P (α + γ), 0) : Vnat) := by
\end{leancode}
(\leanpath{VOA/Moonshine/AutFaithfulVnat.lean}.)
\begin{dictionary}
\lean{leech.form α γ} is $\langle\alpha,\gamma\rangle\in\Z$, and the mode index $-\langle\alpha,\gamma\rangle-1$ is an
integer; \lean{latticeCocycle leech α γ} is $\varepsilon(\alpha,\gamma)$.  The first statement lives in $\chtwentyVL$, the
second in $\Vnat$ (the untwisted block acts on itself by the $\chtwentyVL$ products, \lean{VnatVOA_smul_inl}).  Exactly
the paper statement.
\end{dictionary}
\paragraph{Proof on paper.}  $e^\alpha(z)e^\gamma=\varepsilon(\alpha,\gamma)z^{\langle\alpha,\gamma\rangle}\exp\bigl(\sum_{n>0}
\alpha(-n)z^n/n\bigr)e^{\alpha+\gamma}$, so $e^\alpha_{(m)}e^\gamma$ is $\varepsilon(\alpha,\gamma)$ times the Schur polynomial of
degree $-m-1-\langle\alpha,\gamma\rangle$; at $m=-\langle\alpha,\gamma\rangle-1$ that degree is $0$ and the result is
$\varepsilon(\alpha,\gamma)e^{\alpha+\gamma}$.  The cross terms $e^{\pm\alpha}_{(m)}e^{\mp\gamma}$ have degree $2\langle\alpha,\gamma\rangle<0$
and vanish (``no Fock dressing'').  The term $e^{-\alpha}_{(m)}e^{-\gamma}$ gives $\varepsilon(-\alpha,-\gamma)e^{-\alpha-\gamma}$ with
$\varepsilon(-\alpha,-\gamma)=\varepsilon(\alpha,\gamma)$ by bimultiplicativity.  Summing, and accounting for the projection
$\tfrac12(1+\theta)$ onto $\chtwentyVL$, gives \eqref{ch20:eq:leading}. \qed
\begin{physics}
The most singular term of the OPE $e^{i\alpha\cdot X}(z)e^{i\gamma\cdot X}(w)\sim(z-w)^{\langle\alpha,\gamma\rangle}e^{i(\alpha+\gamma)\cdot X}(w)$:
momentum is additive, and the leading coefficient is the cocycle phase.
\end{physics}
\begin{status}
Unconditional.  Both statements are printed in \leanpath{VOA/Moonshine/MoonshineLRLiftAudit.lean} (all theorems of the \texttt{Aut*} files; green at the release, standard axioms, no \lean{sorryAx}).
\end{status}

\begin{theorem}[The Leech lattice is generated by its minimal vectors]\label{ch20:thm:leechspan}
Every $\beta\in\Leech$ is an integral linear combination of vectors of norm $4$.  Equivalently: if a property $P$ of
lattice vectors holds at $0$ and is preserved by $\gamma\mapsto\gamma\pm\alpha$ for every minimal $\alpha$, it holds on
all of $\Leech$.
\end{theorem}
\begin{leancode}
theorem leech_mem_closure_minimal (β : leech.carrier) :
    β ∈ AddSubgroup.closure {α : leech.carrier | leech.form α α = 4} := by

theorem leech_closure_minimal_eq_top :
    AddSubgroup.closure {α : leech.carrier | leech.form α α = 4} = ⊤ :=

theorem leech_induction (P : leech.carrier → Prop) (h0 : P 0)
    (hadd : ∀ γ α : leech.carrier, leech.form α α = 4 → P γ → P (γ + α))
    (hsub : ∀ γ α : leech.carrier, leech.form α α = 4 → P γ → P (γ - α))
    (β : leech.carrier) : P β := by
\end{leancode}
(\leanpath{VOA/Moonshine/AutLeechSpan.lean}, namespace \lean{VOA.Moonshine.LeechSpan}.)
\begin{dictionary}
\lean{AddSubgroup.closure S} is the subgroup generated by $S$ (all finite integral combinations); \lean{= ⊤} says it is
all of $\Leech$.  The induction principle is the usual consequence.  Statement exactly as on paper.
\end{dictionary}
\paragraph{Proof on paper.}  We use the concrete model of Chapter~\ref{ch:02}: $\Leech$ is realised inside $\Z^{24}$, scaled
by $\sqrt8$ (form $x\cdot y/8$, so minimal vectors have $x\cdot x=32$), as $\Leech=2A\cup(u+2A)$, where
$A=\{a\in\Z^{24}: a\bmod2\text{ is a Golay codeword},\ 4\mid\sum_ia_i\}$ and $u=(5,1^{23})$ (the letter $A$ here has nothing to do
with Monster classes).  In this model one exhibits minimal generators and writes everything in terms of them:
the vectors $4e_i\pm4e_j$ generate $4D_{24}$; the octad vectors $2\cdot(\text{octad})$ are minimal, and the twelve
rows of the generator matrix of the Golay code (one of them a dodecad, handled by an explicit identity) give all of
$2A$ modulo $4D_{24}$; finally the odd coset is reached through the minimal vector $(-3,1^{23})$.  The finitely many
concrete facts about explicit integer vectors are checked by kernel computation (\lean{decide +kernel}; no
\lean{native_decide}). \qed
\begin{physics}
The charges of the cosine operators of weight two generate all charges: the minimal-momentum states already ``see''
the whole charge lattice.
\end{physics}
\begin{status}
Unconditional.  Printed in \leanpath{VOA/Moonshine/MoonshineLRLiftAudit.lean} (all theorems of the \texttt{Aut*} files; green at the release, standard axioms, no \lean{sorryAx}).
\end{status}

\begin{warning}
\textbf{A false induction criterion.}  A natural first idea for Step~5 is to show, by
induction on $|\beta|^2$, that the subspace of vectors fixed by $g$ contains every $x_\beta$, reaching $x_\beta$ from shorter
vectors through a product with some $x_\alpha$, $\alpha$ minimal.  This needs
the claim ``every $\beta$ with $|\beta|^2\ge6$ has a minimal $\alpha$ with $\langle\alpha,\beta\rangle\in\{3,4\}$''.  This is
\emph{false}: for $\beta=k\alpha_0$ with $k\ge5$ and $\alpha_0$ minimal, all inner products $\langle\alpha,\beta\rangle$ lie in
$k\Z$ (a numerical sample found $400$ of $4800$ failures, all of this shape).  The proof given above avoids the issue: it
propagates \emph{scalars} $c_\beta$ along $\gamma\mapsto\gamma\pm\alpha$ in both directions, which only needs that minimal
vectors \emph{generate} $\Leech$ as a group (Theorem~\ref{ch20:thm:leechspan}), not that norms decrease.
\end{warning}

\begin{lemma}[Scalars and their propagation]\label{ch20:lem:scalars}
If $g$ fixes $\Vnat_2$ pointwise then $g(x_\beta,0)=c(x_\beta,0)$ for some $c\in\C$, and the scalars at $\gamma$ and at
$\alpha+\gamma$ agree for every minimal $\alpha$.
\end{lemma}
\begin{leancode}
theorem aut_frameX0P_eigen (g : AutVnat) (hg : ∀ v ∈ VnatVOA.grading 2, g v = v)
    (β : leech.carrier) :
    ∃ c : ℂ, g ((frameX0P β, 0) : Vnat) = c • ((frameX0P β, 0) : Vnat) := by

theorem scalar_step (g : AutVnat) (hg : ∀ v ∈ VnatVOA.grading 2, g v = v)
    {α γ : leech.carrier} (hα : leech.form α α = 4) {c c' : ℂ}
    (hc : g ((frameX0P γ, 0) : Vnat) = c • ((frameX0P γ, 0) : Vnat))
    (hc' : g ((frameX0P (α + γ), 0) : Vnat) = c' • ((frameX0P (α + γ), 0) : Vnat)) : c = c' := by
\end{leancode}
(\leanpath{VOA/Moonshine/AutFaithfulVnat.lean}.)
\begin{dictionary}
Steps 1--4 and Step 5 of the proof of Theorem~\ref{ch20:thm:faithful}.  The hypothesis \lean{hg} is used only on the
vectors $(q_{ab},0)$ (first statement) and $(x_\alpha,0)$, $\alpha$ minimal (second statement); the finite-order proof
reuses private copies of these two lemmas with exactly these weakened hypotheses (\S\ref{ch20:sec:finorder}).
\end{dictionary}
\begin{physics}
A symmetry fixing the weight-two fields acts on each cosine operator by a phase, and the phases are transported by the
leading OPE coefficient.
\end{physics}
\begin{status}
Unconditional.  Printed in \leanpath{VOA/Moonshine/MoonshineLRLiftAudit.lean} (all theorems of the \texttt{Aut*} files; green at the release, standard axioms, no \lean{sorryAx}).
\end{status}

%======================================================================
\section{The triality automorphism}\label{ch20:sec:sigma}
%======================================================================

The extra symmetry of $\Vnat$ that is not visible in the lattice description is FLM's involution $\sigma$.

\paragraph{What the rest of the chapter uses.}  Only three facts about $\sigma$ are used below, and nothing else about it
matters for the finite-order and finiteness arguments:
\begin{enumerate}[label=($\sigma$\arabic*)]
\item $\sigma$ is an automorphism of $\Vnat$ (Theorem~\ref{ch20:thm:sigma});
\item $\sigma^2=1$ (\lean{sigNat_invol});
\item for the twisted generator $x\in\chtwentyTp$ of Lemma~\ref{ch20:lem:fixising}, $\sigma(0,x)=(\Psi x,0)$ lies in
  $\chtwentyVL\times0$ (\lean{sigNat_Ts}, Chapter~\ref{ch:14}).
\end{enumerate}
A reader willing to take these as given can skip to \S\ref{ch20:sec:periodic}.  The library never states $\sigma\neq1$ as a
theorem, and no argument of this chapter needs it.

\paragraph{Where $\sigma$ comes from.}  Chapters~\ref{ch:11}--\ref{ch:14} decompose $\Vnat$ along a Golay-code frame into
four blocks, $\Vnat=V_0^+\oplus V_3^+\oplus T_s^+\oplus T_{s+1}^+$: $\chtwentyVL=V_0^+\oplus V_3^+$ splits into the two
eigenspaces of a frame symmetry, and $\chtwentyTp=T_s^+\oplus T_{s+1}^+$ into the two frame sectors.  The library's
$\sigma$ is the operator \lean{sigNat D} of Chapter~\ref{ch:14}: on $V_0^+$ it is a ``code triality'' automorphism of $V_0^+$;
it exchanges $V_3^+$ and $T_s^+$ through mutually inverse linear isomorphisms $\Phi:V_3^+\to T_s^+$, $\Psi:T_s^+\to V_3^+$;
and on $T_{s+1}^+$ it is a linear involution $\Theta$.  These maps are part of a \emph{triality datum} \lean{SigmaData}; the
standard one, \lean{Dnat}, is built from a normalising scalar $c^\star\in\C$ (with $c^{\star2}\kappa_0=1$ for a fixed
cocycle sign $\kappa_0$) and the sign $\varepsilon=+1$ (Chapter~\ref{ch:14}), and its
\emph{$\sigma$-key} (\lean{SigmaKey.conj_eq}, the main identity proved in Chapters~\ref{ch:12}--\ref{ch:14}) says that
$\sigma$ conjugates the vertex operator of $x$ into that of $\sigma(0,x)$.  What was not stated there is that $\sigma$ is an
automorphism of the vertex operator algebra \lean{VnatVOA}.

\begin{theorem}[$\sigma^\natural\in\Aut(\Vnat)$]\label{ch20:thm:sigma}
The involution $\sigma^\natural=$ \lean{sigNat Dnat} is an automorphism of $\Vnat$: it is a linear involution, it
intertwines every product, and it fixes the vacuum and the conformal vector.
\end{theorem}
\begin{leancode}
noncomputable abbrev Dnat : SigmaData :=
  mkData cStar0 1 (c_ne_of_norm cStar0_norm) (by norm_num)

noncomputable def sigmaEquiv : Vnat ≃ₗ[ℂ] Vnat :=
  LinearEquiv.ofInvolutive (sigNat Dnat) sigNat_invol

noncomputable def sigmaAut : AutVnat where
  toLinearEquiv := sigmaEquiv

@[simp] theorem sigmaAut_apply (v : Vnat) : sigmaAut v = sigNat Dnat v := rfl
\end{leancode}
(\leanpath{VOA/Moonshine/AutSigma.lean} and, for \lean{Dnat}, \leanpath{VOA/Moonshine/AutVnat.lean}; namespace
\lean{VOA.Moonshine.AutMonster}; the three proof fields of \lean{sigmaAut} are omitted.)
\begin{dictionary}
\lean{Dnat} is the standard triality datum $D(c^\star,+1)$ (the sign $\varepsilon=+1$ of Chapter~\ref{ch:14}) used by the
unconditional theorem \lean{moonshine}.  \lean{LinearEquiv.ofInvolutive f hf} turns a linear map with $f\circ f=\mathrm{id}$
(\lean{sigNat_invol}, from \lean{sigNat_sq}) into a linear equivalence.  \lean{sigmaAut : AutVnat} is a term of the
structure \lean{Aut}; its existence \emph{is} the theorem: the three omitted fields are the proofs of
$\sigma(u_{(n)}v)=(\sigma u)_{(n)}(\sigma v)$ for all $u,v,n$, $\sigma\chtwentyvac=\chtwentyvac$, $\sigma\omega=\omega$.
\lean{sigmaAut_apply} says that, as a function, it is \lean{sigNat Dnat}.  So the statement is the textbook one.
It is used below only through $\sigma(0,x)=(\Psi x,0)$ for the twisted generator $x$ (\lean{sigNat_Ts}), and
$\sigma^2=1$.
\end{dictionary}
\paragraph{Proof on paper.}  Let $S$ be the space of $u\in\Vnat$ whose modes $\sigma$ intertwines
($\sigma u_{(n)}\sigma^{-1}=(\sigma u)_{(n)}$ for all $n$).  It is closed under products (Borcherds identity) and, $\sigma$
being an involution, $\sigma$-stable.  It contains $V_0^+\times0$ (Chapter~\ref{ch:14}: $\sigma$ intertwines the
$V_0^+$-modes) and the twisted generator $(0,x)$, $x=h_0(-\tfrac12)\chtwentyvac\otimes t_s$ (this is the $\sigma$-key
\lean{SigmaKey.conj_eq}, the content of Chapters~\ref{ch:12}--\ref{ch:14}).  Hence it contains $0\times T_s^+$, which is
spanned by $V_0^+$-modes applied to $x$; hence $V_3^+\times0=\sigma(0\times T_s^+)$; hence $\chtwentyVL\times0$; hence
$0\times\chtwentyTp$, which is spanned by $\chtwentyVL$-modes applied to $x$ (\lean{modeSpan_x0T}).  So $S=\Vnat$.  The
vacuum and $\omega$ lie in $V_0^+$, where $\sigma$ acts by an automorphism of $V_0^+$ fixing them. \qed
\begin{physics}
$\sigma$ is a symmetry of the orbifold theory that mixes the untwisted sector with the twisted sector: a ``quantum''
symmetry invisible in the free-boson description (it exchanges a cosine operator with a twist field).  Together with
the lattice symmetries it generates the Monster; in this library it is the one extra element needed, in
\S\ref{ch20:sec:finorder}, to reach the twisted sector from the untwisted one.
\end{physics}
\begin{status}
Unconditional definition (a term of \lean{AutVnat}); \lean{sigmaAut_apply} is printed in \leanpath{VOA/Moonshine/MoonshineLRLiftAudit.lean} (all theorems of the \texttt{Aut*} files; green at the release, standard axioms, no \lean{sorryAx}), and
\lean{sigmaAut} is used by \lean{autVnat_isOfFinOrder} and \lean{finite_autVnat}, whose audit lines in
\leanpath{VOA/Moonshine/MoonshineLRFinOrderAudit.lean} are standard.
\end{status}

%======================================================================
\section{Periodicity of rigid idempotents: an abstract lemma}\label{ch20:sec:periodic}
%======================================================================

\subsection{Why finite order, and the plan}

The twisted denominator identities of Chapter~\ref{ch:19} are proved for automorphisms of finite order, and the
analytic half of complete replicability needs the trace bound $|\tr(g|\Vnat_n)|\le\dim\Vnat_n$, which holds for
finite-order $g$ (Chapter~\ref{ch:16}, Theorem~\ref{ch16:thm:tracebounds}).  So the conjecture needs: \emph{every
automorphism of $\Vnat$ has finite order.}  Classically this is a corollary of $\Aut(\Vnat)=\mathbb M$; the library
does not have that identification unconditionally, and the proof follows a route that uses no group theory:
\begin{enumerate}[label=(\arabic*)]
\item an abstract lemma: in a finite-dimensional commutative algebra, every \emph{rigid idempotent} is periodic under
  every algebra automorphism (this section);
\item explicit \emph{Ising vectors} $e_{\alpha,\pm s}\in\Vnat_2$, one pair for each minimal $\alpha$, which are rigid
  idempotents of the Griess algebra (\S\ref{ch20:sec:ising});
\item faithfulness, with weakened hypotheses: an automorphism fixing all these Ising vectors fixes $\chtwentyVL$; conjugation
  by $\sigma$ reaches the twisted sector (\S\ref{ch20:sec:finorder}).
\end{enumerate}

\begin{definition}[Rigid idempotent]\label{ch20:def:rigid}
Let $B$ be a complex vector space with a bilinear product $\cdot$ (not necessarily associative or unital).  An element
$e\in B$ is a \emph{rigid idempotent} if $e\cdot e=2e$ and $e\cdot y=y$ implies $y=0$ (i.e.\ $1$ is not an eigenvalue of
left multiplication $L_e$).
\end{definition}
There is no separate Lean definition: the two conditions appear as hypotheses \lean{he : mul e e = (2 : ℂ) • e} and
\lean{hrig : ∀ y, mul e y = y → y = 0} of every statement below (and as the set \lean{vnatRigidSet} in
\S\ref{ch20:sec:finite}).  The normalisation $e\cdot e=2e$ (rather than $e\cdot e=e$) is the one of conformal vectors in
a Griess algebra: a Virasoro vector $\omega'$ with $\omega'_{(1)}\omega'=L'_0\omega'=2\omega'$.  Note that $e=0$ is a rigid
idempotent by this definition; this is harmless (the theorems below are about sets containing it, or about specific non-zero
vectors), but it shows that ``rigid idempotent'' is a purely algebraic notion, wider than ``stress tensor of a sub-theory''.

\begin{physics}
A rigid idempotent of the Griess algebra is a candidate stress tensor $T'(z)$ of a sub-theory whose embedding cannot be
deformed: an infinitesimal deformation $e+\epsilon y$ of the relation $e\cdot e=2e$ requires $2e\cdot y=2y$, i.e.\ a field
$y$ of conformal weight exactly $1$ with respect to $T'$ (a ``marginal'' direction for the embedding).  Rigidity says there
is none.  The stress tensor $\omega$ of $\Vnat$ is itself rigid ($L_0=2$ on $\Vnat_2$); the $c=1$ Virasoro vectors
$\tfrac12h(-1)^2\chtwentyvac$ are \emph{not} (they come in a continuous $O(24)$-family), which is why the argument must
use $c=\tfrac12$ Ising vectors.
\end{physics}

\subsection{The lemma and its proof}

\begin{theorem}[Periodicity of rigid idempotents]\label{ch20:thm:periodic}
Let $B$ be a finite-dimensional complex vector space with a commutative bilinear product, $A:B\to B$ a linear
automorphism with $A(x\cdot y)=Ax\cdot Ay$, and $e$ a rigid idempotent.  Then $A^Ke=e$ for some $K\ge1$.
\end{theorem}
\begin{leancode}
theorem pow_apply_eq_self_of_rigid [FiniteDimensional ℂ B]
    (mul : B →ₗ[ℂ] B →ₗ[ℂ] B) (hcomm : ∀ x y, mul x y = mul y x)
    (A : B ≃ₗ[ℂ] B) (hA : ∀ x y, A (mul x y) = mul (A x) (A y))
    (e : B) (he : mul e e = (2 : ℂ) • e)
    (hrig : ∀ y, mul e y = y → y = 0) :
    ∃ K : ℕ, 1 ≤ K ∧ (A ^ K) e = e := by
\end{leancode}
(\leanpath{VOA/Moonshine/MoonshineLRFinOrderPeriodic.lean}, namespace \chtwentylean{VOA.Moonshine.ConwayNorton.M8}; section
variables \lean{{B : Type*} [AddCommGroup B] [Module ℂ B]}.)
\begin{dictionary}
\lean{B →ₗ[ℂ] B →ₗ[ℂ] B} is a bilinear map, the product; \lean{hcomm} is commutativity.  \lean{A : B ≃ₗ[ℂ] B} is a linear
automorphism, \lean{hA} multiplicativity; \lean{A ^ K} is the $K$-th power in the group of linear automorphisms.
\lean{he}, \lean{hrig} are the two conditions of Definition~\ref{ch20:def:rigid}.  No associativity, no unit, no
form, no VOA: the lemma is pure finite-dimensional algebra (the file imports Mathlib only).  The conclusion
\lean{∃ K, 1 ≤ K ∧ (A ^ K) e = e} is periodicity of $e$ under $A$.  Exactly the paper statement.
\end{dictionary}

\paragraph{Proof on paper.}  Decompose $B=\bigoplus_\lambda W_\lambda$ into generalised eigenspaces of $A$ ($\lambda\neq0$
since $A$ is invertible), and write $x=\sum_\lambda x_\lambda$.
\begin{enumerate}[label=(\roman*)]
\item \emph{Multiplicativity:} $W_\lambda\cdot W_\mu\subseteq W_{\lambda\mu}$.  From $A(x\cdot y)=Ax\cdot Ay$ one gets the
  shifted product rule $(A-\lambda\mu)(x\cdot y)=((A-\lambda)x)\cdot Ay+\lambda\,x\cdot((A-\mu)y)$, and by induction
  $(A-\lambda\mu)^{n+m}(x\cdot y)=0$ when $(A-\lambda)^nx=0$, $(A-\mu)^my=0$.
\item \emph{Character derivations.}  Let $\varphi:\C^\times\to\C$ be a homomorphism from the multiplicative to the
  additive group, and $D_\varphi x=\sum_\lambda\varphi(\lambda)x_\lambda$.  By (i), $D_\varphi$ is a derivation:
  $D_\varphi(x\cdot y)=D_\varphi x\cdot y+x\cdot D_\varphi y$.
\item \emph{Derivations kill rigid idempotents.}  Applying a derivation $D$ to $e\cdot e=2e$ and using commutativity gives
  $2\,e\cdot De=2De$, so $De=0$ by rigidity.  Hence $\varphi(\lambda)e_\lambda=0$ for every $\lambda$ and every $\varphi$.
\item \emph{Separating characters.}  If $\lambda\in\C^\times$ has infinite order, the subgroup $\lambda^\Z\cong\Z$ carries the
  homomorphism $\lambda^n\mapsto n$; since $\C$ is a divisible, hence injective, abelian group (Baer's criterion), it
  extends to $\varphi:\C^\times\to\C$ with $\varphi(\lambda)=1$.  By (iii), $e_\lambda=0$.  So every $\lambda$ in the support
  of $e$ is a root of unity; let $K$ be the product of their orders.
\item \emph{The semisimple part.}  Let $S$ act by $\lambda$ on $W_\lambda$.  By (i) $S$ is multiplicative, it commutes with
  $A$, and $S^Ke=\sum\lambda^Ke_\lambda=e$.  $U=S^{-1}A$ is a multiplicative automorphism with $U-1$ nilpotent.
\item \emph{Unipotent automorphisms fix rigid idempotents.}  For a multiplicative $U$ with $(U-1)^de=0$, the orbit
  $k\mapsto U^ke=\sum_{j<d}\binom kj(U-1)^je$ is the restriction to $\N$ of a polynomial curve $P(t)$ (with coefficients
  $\binom tj=t(t-1)\cdots(t-j+1)/j!$).  Each $P(n)=U^ne$ is a rigid idempotent (multiplicative automorphisms preserve
  rigidity).  The identity $P(t)\cdot P(t)=2P(t)$ holds for $t\in\N$, hence for all $t$ (a polynomial with infinitely many
  roots is zero, applied coordinatewise through linear functionals).  Differentiating, $2P(t)\cdot P'(t)=2P'(t)$; at
  $t=n\in\N$ rigidity of $P(n)$ gives $P'(n)=0$; so $P'$ vanishes on $\N$, hence identically, and $P$ is constant:
  $Ue=P(1)=P(0)=e$.
\end{enumerate}
Applying (vi) to $U^K$ (still unipotent and multiplicative) and using that $S$ and $U$ commute,
$A^Ke=S^KU^Ke=S^Ke=e$. \qed

\begin{physics}
Heuristically: an eigenvalue $\lambda$ of a symmetry that is not a root of unity would generate, through
$\lambda^t=e^{t\log\lambda}$, a one-parameter family of symmetries, i.e.\ a conserved ``charge'' (the derivation $D_\varphi$
plays its role without logarithms).  A rigid idempotent cannot be moved by any such infinitesimal symmetry, so it can only
carry charges of finite order.  Step (vi) is the unipotent analogue: a symmetry of the form $\exp(N)$, $N$ nilpotent, is a
continuous deformation and cannot move a rigid object.
\end{physics}
\begin{status}
\textbf{Unconditional}, generic.  Printed in \leanpath{VOA/Moonshine/MoonshineLRFinOrderAudit.lean} with
\lean{[propext, Classical.choice, Quot.sound]}, 0 \lean{sorryAx} (build record of the release).
\end{status}

\subsection{The steps, in Lean}

The individual steps (i)--(vi) are separate declarations of the same file.  We quote the public ones that a referee
may want to check; their meaning is the corresponding step above.
\begin{leancode}
theorem maxGenEigenspace_mul_le (mul : B →ₗ[ℂ] B →ₗ[ℂ] B)
    (A : Module.End ℂ B) (hA : ∀ x y, A (mul x y) = mul (A x) (A y))
    (l μ : ℂ) {x y : B}
    (hx : x ∈ A.maxGenEigenspace l) (hy : y ∈ A.maxGenEigenspace μ) :
    mul x y ∈ A.maxGenEigenspace (l * μ) := by

theorem derivation_apply_eq_zero_of_rigid (mul : B →ₗ[ℂ] B →ₗ[ℂ] B)
    (hcomm : ∀ x y, mul x y = mul y x) (D : Module.End ℂ B)
    (hD : ∀ x y, D (mul x y) = mul (D x) y + mul x (D y))
    (e : B) (he : mul e e = (2 : ℂ) • e)
    (hrig : ∀ y, mul e y = y → y = 0) : D e = 0 := by

theorem exists_addChar_eq_one_of_not_isOfFinOrder {G : Type*} [CommGroup G]
    (g : G) (hg : ¬ IsOfFinOrder g) :
    ∃ φ : Additive G →+ ℂ, φ (Additive.ofMul g) = 1 := by

theorem isOfFinOrder_of_rigid_component (mul : B →ₗ[ℂ] B →ₗ[ℂ] B)
    (hcomm : ∀ x y, mul x y = mul y x)
    (A : B ≃ₗ[ℂ] B) (hA : ∀ x y, A (mul x y) = mul (A x) (A y))
    (e : B) (he : mul e e = (2 : ℂ) • e)
    (hrig : ∀ y, mul e y = y → y = 0)
    (l : ℂ) (hl : spectralDecomposition A e l ≠ 0) : IsOfFinOrder l := by

theorem polynomialCurve_eq_zero_value_of_rigid (mul : B →ₗ[ℂ] B →ₗ[ℂ] B)
    (hcomm : ∀ x y, mul x y = mul y x) (p : ι → Polynomial ℂ) (v : ι → B)
    (he : ∀ n : ℕ, mul (polynomialCurve p v n) (polynomialCurve p v n) =
      (2 : ℂ) • polynomialCurve p v n)
    (hrig : ∀ (n : ℕ) y, mul (polynomialCurve p v n) y = y → y = 0)
    (t : ℂ) : polynomialCurve p v t = polynomialCurve p v 0 := by

theorem unipotent_fix_of_rigid (mul : B →ₗ[ℂ] B →ₗ[ℂ] B)
    (hcomm : ∀ x y, mul x y = mul y x)
    (U : B ≃ₗ[ℂ] B) (hU : ∀ x y, U (mul x y) = mul (U x) (U y))
    (e : B) (he : mul e e = (2 : ℂ) • e)
    (hrig : ∀ y, mul e y = y → y = 0)
    (huni : ∃ d : ℕ, (((U : Module.End ℂ B) - 1) ^ d) e = 0) : U e = e := by
\end{leancode}
(\leanpath{VOA/Moonshine/MoonshineLRFinOrderPeriodic.lean}.)
\begin{dictionary}
\lean{A.maxGenEigenspace l} is the generalised eigenspace $W_\lambda=\bigcup_n\ker(A-\lambda)^n$.
\lean{Additive G →+ ℂ} is a homomorphism from the group $G$ written additively to $(\C,+)$, i.e.\ a ``logarithmic
character''; \lean{IsOfFinOrder g} is $\exists n\ge1,\ g^n=1$.  \lean{spectralDecomposition A e l} is the component
$e_\lambda\in W_\lambda$ (defined from Mathlib's \chtwentylean{iSup_maxGenEigenspace_eq_top} and
\chtwentylean{independent_maxGenEigenspace}).  \lean{polynomialCurve p v t} is $\sum_ip_i(t)v_i$, a polynomial curve in $B$
indexed by a finite type \lean{ι}; the theorem says that if its values at all natural numbers are rigid idempotents,
then it is constant.  \lean{unipotent_fix_of_rigid} has no finite-dimensionality hypothesis: local unipotence at $e$
(\lean{huni}) suffices.  All conditions match the steps above.
\end{dictionary}
\begin{physics}
These are the algebraic shadows of: ``charges are additive under the OPE'' (i), ``a $U(1)$ charge acts as a
derivation of the operator algebra'' (ii), and ``a rigid sub-theory carries zero continuous charge'' (iii).
\end{physics}
\begin{status}
All unconditional, generic, printed in \leanpath{VOA/Moonshine/MoonshineLRFinOrderAudit.lean}; standard axioms.
\end{status}

%======================================================================
\section{Ising vectors in the Griess algebra}\label{ch20:sec:ising}
%======================================================================

\subsection{The vectors}

\begin{definition}[Explicit Ising vectors]\label{ch20:def:ising}
For a minimal vector $\alpha\in\chtwentyLf$ and a scalar $s\in\C$ put
\begin{equation}\label{ch20:eq:ising}
  e_{\alpha,s}=\Bigl(\tfrac1{16}\,Q(\alpha,\alpha)+\tfrac s4\,x_\alpha,\;0\Bigr)\in\Vnat=\chtwentyVL\oplus\chtwentyTp,
  \qquad Q(\alpha,\alpha)=\alpha(-1)^2\chtwentyvac .
\end{equation}
The \emph{normalised} ones are those with $s^2\,\varepsilon(\alpha,-\alpha)=1$ (so $s\in\{\pm1\}$ if
$\varepsilon(\alpha,-\alpha)=1$, $s\in\{\pm i\}$ otherwise).
\end{definition}
\begin{leancode}
noncomputable def isingVec (α : leech.carrier) (s : ℂ) : Vnat :=
  ((1 / 16 : ℂ) • (dV α α : VPlus) + (s / 4) • frameX0P α, 0)

theorem isingVec_mem_grading_two {α : leech.carrier} (hα : leech.form α α = 4)
    (s : ℂ) : isingVec α s ∈ VnatVOA.grading 2 := by

theorem isingVec_one_self {α : leech.carrier} (hα : leech.form α α = 4) (s : ℂ)
    (hs : s ^ 2 * latticeCocycle leech α (-α) = 1) :
    VnatVOA.smul (isingVec α s) 1 (isingVec α s) = (2 : ℂ) • isingVec α s := by
\end{leancode}
(\leanpath{VOA/Moonshine/MoonshineLRFinOrderIsing.lean}, namespace \chtwentylean{VOA.Moonshine.ConwayNorton.M8}.)
\begin{dictionary}
\lean{dV α α} is the state $\alpha(-1)\alpha(-1)\chtwentyvac=Q(\alpha,\alpha)$ in $V_0^+\subseteq\chtwentyVL$ (it is
\lean{quadH}, defined as $h_\alpha{}_{(-1)}h_\alpha$; Chapter~\ref{ch:13}), coerced to \lean{VPlus}.  The pair \lean{(…, 0)}
places it in the untwisted block.  \lean{isingVec} is defined for every $s$; the normalisation appears only as the
hypothesis \lean{hs}, with $\kappa_\alpha=\varepsilon(\alpha,-\alpha)$ left symbolic.  \lean{isingVec_one_self} is $e\times e=2e$
for the Griess product ($e_{(1)}e$).  \textbf{What is not claimed.}  The Lean statements do not assert that $e$ generates a
Virasoro algebra of central charge $\tfrac12$ (that would need $e_{(3)}e=\tfrac14\chtwentyvac$ and the Virasoro relations), nor
that these are all the Ising vectors of $\Vnat$; the name ``Ising vector'' is motivation.  Only $e\times e=2e$ and the
rigidity of \S\ref{ch20:sec:rigid} are proved, and only they are used.
\end{dictionary}
\paragraph{Proof on paper.}  Three Griess products are needed (Lemma~\ref{ch20:lem:products} below):
$Q(\alpha,\alpha)\times Q(\alpha,\alpha)=16Q(\alpha,\alpha)$, $Q(\alpha,\alpha)\times x_\alpha=x_\alpha\times Q(\alpha,\alpha)=16x_\alpha$ and
$x_\alpha\times x_\alpha=\kappa_\alpha Q(\alpha,\alpha)$.  Hence
\begin{equation*}
  e\times e=\tfrac1{256}\cdot16Q+2\cdot\tfrac{s}{64}\cdot16x_\alpha+\tfrac{s^2}{16}\kappa_\alpha Q
  =\tfrac{1+s^2\kappa_\alpha}{16}Q+\tfrac s2x_\alpha=2e
\end{equation*}
exactly when $s^2\kappa_\alpha=1$. \qed
\begin{physics}
These are the classical Ising vectors $\omega^\pm(\alpha)=\tfrac1{16}\alpha(-1)^2\chtwentyvac\pm\tfrac14(e^\alpha+e^{-\alpha})$ of
Dong--Mason--Zhu and Miyamoto: the stress tensors of an Ising model (the $c=\tfrac12$ minimal model, one free Majorana
fermion) embedded in the $24$ free bosons.  The rank-one lattice theory of $\Z\alpha$ ($\langle\alpha,\alpha\rangle=4$, so that
$\cos(\alpha\cdot X)$ has dimension $2$) contains two commuting Ising models, $c=1=\tfrac12+\tfrac12$, and $\omega^\pm(\alpha)$ are
their two stress tensors.  Miyamoto showed that each Ising vector $e$ defines an involution $\tau_e$ (acting by $-1$ on the
$h=\tfrac1{16}$ sector of $\mathrm{Vir}_e$), and in $\Vnat$ these are the $2A$ involutions of the Monster.  None of this
is used by the library; it explains why these vectors are the natural rigid anchors.
\end{physics}
\begin{status}
\textbf{Unconditional} given \lean{hα}, \lean{hs} (which are the definition of the vectors, not open inputs).  Printed in
\leanpath{VOA/Moonshine/MoonshineLRFinOrderAudit.lean}; standard axioms.
\end{status}

\begin{warning}
\textbf{The naive Ising vector need not be idempotent.}  The literal vector
$\tfrac1{16}\alpha(-1)^2\chtwentyvac+\tfrac14x_\alpha$ satisfies $e\times e=2e$ only if $\varepsilon(\alpha,-\alpha)=1$.  The library's
cocycle comes from a basis chosen by \lean{Module.Free.chooseBasis} and its signs are never evaluated, so the scalar
$s$ is kept symbolic with the single hypothesis $s^2\varepsilon(\alpha,-\alpha)=1$.  The concrete choice used later is
$s=\pm$\lean{FrameT.sqrtEps leech α}, a chosen square root of $\varepsilon(\alpha,\alpha)$ (and $\varepsilon(\alpha,-\alpha)=\varepsilon(\alpha,\alpha)$
since $\varepsilon^2=1$).
\end{warning}

\subsection{Weight-two products and spanning sets}

\begin{lemma}[Products in $(\chtwentyVL)_2$]\label{ch20:lem:products}
Let $\alpha,\beta\in\chtwentyLf$.  (a) $\langle\alpha,\beta\rangle\in\{-4,-2,-1,0,1,2,4\}$.  (b) $x_\alpha\times x_\alpha=\varepsilon(\alpha,-\alpha)Q(\alpha,\alpha)$;
$x_\alpha\times x_\beta=\varepsilon(\alpha,\beta)x_{\alpha+\beta}$ if $\langle\alpha,\beta\rangle=-2$; $x_\alpha\times x_\beta=\varepsilon(\alpha,-\beta)x_{\beta-\alpha}$ if
$\langle\alpha,\beta\rangle=2$; $x_\alpha\times x_\beta=0$ if $|\langle\alpha,\beta\rangle|\le1$.  (c) $Q(u,v)\times x_\beta=\langle u,\beta\rangle\langle v,\beta\rangle x_\beta$
for all $u,v,\beta\in\Leech$, and $x_\alpha\times Q(u,v)=\langle u,\alpha\rangle\langle v,\alpha\rangle x_\alpha$.  (d)
$q_{ab}\times q_{cd}=\delta_{bc}q_{ad}+\delta_{bd}q_{ac}+\delta_{ac}q_{bd}+\delta_{ad}q_{bc}$; equivalently, by bilinearity, for
$u,v,w,z\in\mathfrak h$,
\begin{equation*}
  Q(u,v)\times Q(w,z)=\langle v,w\rangle Q(u,z)+\langle v,z\rangle Q(u,w)+\langle u,w\rangle Q(v,z)+\langle u,z\rangle Q(v,w),
\end{equation*}
so in particular $Q(\alpha,\alpha)\times Q(\alpha,\alpha)=16\,Q(\alpha,\alpha)$ and $Q(\alpha,\alpha)\times Q(\alpha,v')=4\,Q(\alpha,v')$ for
$v'\perp\alpha$, $\langle\alpha,\alpha\rangle=4$ (Lean: \lean{isingQuad_one_isingQuad}, \lean{quadH_one_quadH_self}).
\end{lemma}
\begin{leancode}
theorem leech_form_cases_of_normFour {α β : leech.carrier}
    (hα : leech.form α α = 4) (hβ : leech.form β β = 4) :
    leech.form α β = -4 ∨ leech.form α β = -2 ∨ leech.form α β = -1 ∨
      leech.form α β = 0 ∨ leech.form α β = 1 ∨ leech.form α β = 2 ∨
      leech.form α β = 4 := by

theorem frameX0P_one_frameX0P_self {α : leech.carrier} (hα : leech.form α α = 4) :
    leechVOAPlus.smul (frameX0P α) 1 (frameX0P α) =
      latticeCocycle leech α (-α) • (dV α α : VPlus) := by

theorem frameX0P_one_frameX0P_of_form_small {α β : leech.carrier}
    (hlo : -1 ≤ leech.form α β) (hhi : leech.form α β ≤ 1) :
    leechVOAPlus.smul (frameX0P α) 1 (frameX0P β) = 0 := by

theorem quadH_one_frameX0P (u v β : leech.carrier) :
    leechVOAPlus.smul (dV u v : VPlus) 1 (frameX0P β) =
      ((leech.form u β : ℂ) * (leech.form v β : ℂ)) • frameX0P β :=

theorem quadState_one_quadState (a b c d : Fin (rank leech)) :
    leechVOAPlus.smul (qP a b) 1 (qP c d) =
      (if b = c then qP a d else 0) + (if b = d then qP a c else 0) +
      (if a = c then qP b d else 0) + (if a = d then qP b c else 0) := by
\end{leancode}
(\leanpath{VOA/Moonshine/MoonshineLRFinOrderIsing.lean}; the cases $\pm2$ are
\chtwentylean{frameX0P_one_frameX0P_of_form_neg_two}, \chtwentylean{frameX0P_one_frameX0P_of_form_two}, and the reversed products
\lean{frameX0P_one_quadH}, \lean{quadH_one_quadH_self} in the same file.)
\begin{dictionary}
\lean{leechVOAPlus.smul x 1 y} is $x_{(1)}y$ in $\chtwentyVL$, the Griess product on weight two.  \lean{(leech.form u β : ℂ)}
casts the integer $\langle u,\beta\rangle$ to $\C$.  In (a) the Lean disjunction lists the seven values.  The statements
are the paper ones; note that (b) keeps the cocycle sign $\varepsilon(\alpha,-\alpha)$ unevaluated.
\end{dictionary}
\paragraph{Proof on paper.}  (a) $|\alpha\mp\beta|^2=8\mp2\langle\alpha,\beta\rangle$ is even, non-negative, never $2$ (no
roots) and $0$ only if $\beta=\pm\alpha$; so $8\mp2\langle\alpha,\beta\rangle\in\{0\}\cup[4,\infty)$ for both signs.  (b) The
$(1)$-product $e^{\alpha}_{(1)}e^{\beta}$ is $\varepsilon(\alpha,\beta)$ times the Schur polynomial of degree
$-2-\langle\alpha,\beta\rangle$ applied to $e^{\alpha+\beta}$: zero if $\langle\alpha,\beta\rangle\ge-1$, the vacuum term if
$\langle\alpha,\beta\rangle=-2$, and $\tfrac12(\alpha(-1)^2+\alpha(-2))\chtwentyvac$ if $\beta=-\alpha$; the two terms with $\beta=-\alpha$ add
to $\varepsilon(\alpha,-\alpha)Q(\alpha,\alpha)$ (the $\alpha(-2)$ parts cancel).  (c) is \eqref{ch20:eq:zeromode} at charge $\pm\beta$ and
commutativity.  (d) is $N_{ab}$ of \eqref{ch20:eq:zeromode} applied to $X_{0,c}X_{0,d}$. \qed
\begin{physics}
OPE coefficients of the weight-two fields $\cos(\beta\cdot X)$ and $\partial X^a\partial X^b$: momenta add, and the
coefficient of $(z-w)^{-2}$ in $\cos\times\cos$ is non-zero only for momenta at angle $\langle\alpha,\beta\rangle=\mp2$ (or
$\beta=-\alpha$, where the identity channel produces $\partial X\partial X$).
\end{physics}
\begin{status}
Unconditional; printed in \leanpath{VOA/Moonshine/MoonshineLRFinOrderAudit.lean}; standard axioms.
\end{status}

\begin{lemma}[Spanning sets of the weight-two blocks]\label{ch20:lem:span}
$(\chtwentyVL)_2$ is spanned by the $q_{ab}$ and the $x_\alpha$, $\alpha\in\chtwentyLf$; $(\chtwentyTp)_2$ is spanned by the vectors
$h_v(-\tfrac12)\otimes t$, $v\in\C^{24}$, $t\in T$.
\end{lemma}
\begin{leancode}
theorem leechVOAPlus_grading_two_eq_span :
    leechVOAPlus.grading 2 = Submodule.span ℂ
      (Set.range (fun p : Fin (rank leech) × Fin (rank leech) => qP p.1 p.2) ∪
        frameX0P '' {α : leech.carrier | leech.form α α = 4}) := by

theorem twistedPlus_grading_two_eq_span :
    plusModule.weight (2 : ℂ) = Submodule.span ℂ
      {y | ∃ (v : Fin (rank leech) → ℂ) (t : TLeechT), y = hP v t} := by
\end{leancode}
(\leanpath{VOA/Moonshine/MoonshineLRFinOrderIsing.lean}.)
\begin{dictionary}
\lean{Submodule.span ℂ S} is the linear span; \lean{Set.range f} the image of $f$; \lean{frameX0P '' S} the image of
the set of minimal vectors.  \lean{plusModule.weight (2 : ℂ)} is the weight-two space of $\chtwentyTp$ as a
$\chtwentyVL$-module, and \lean{hP v t} is $h_v(-\tfrac12)\chtwentyvac\otimes t$ with $h_v=\sum_cv_ch_c$ and $t\in T$
(\lean{TLeechT}, Chapter~\ref{ch:13}).  The dimensions are $300+98280=98580$ and $24\cdot2^{12}=98304$, summing to $196884$
(Chapter~\ref{ch:04}); the lemma asserts only spanning.
\end{dictionary}
\begin{physics}
The $196884$ weight-two fields of the Monster CFT are: $300$ bilinears $\partial X^a\partial X^b$, $98280$ cosines
$\cos(\alpha\cdot X)$ with $\alpha$ minimal (up to sign), and $98304=24\cdot4096$ twisted states (one half-integral oscillator on
the $4096$ twisted ground states).
\end{physics}
\begin{status}
Unconditional; printed in the FinOrder audit; standard axioms.
\end{status}

\subsection{Rigidity}\label{ch20:sec:rigid}

\begin{theorem}[The spectral lemma and rigidity]\label{ch20:thm:rigid}
Let $\alpha\in\chtwentyLf$, $s^2\varepsilon(\alpha,-\alpha)=1$, $e=e_{\alpha,s}$ and $X=o(e)$.  Then
\begin{enumerate}[label=(\alph*)]
\item on $(\chtwentyVL)_2$: $X(X-2)(X-\tfrac12)(X-\tfrac1{16})=0$;
\item on $(\chtwentyTp)_2$: $X(X-\tfrac1{16})(X-\tfrac12)=0$; explicitly
  $X\bigl(h_v(-\tfrac12)\otimes t\bigr)=\tfrac1{32}\bigl[h_{(1+8P_\alpha)v}(-\tfrac12)\otimes t+h_{(1-8P_\alpha)v}(-\tfrac12)\otimes
  Et\bigr]$, where $P_\alpha v=\tfrac14\langle v,\alpha\rangle\alpha$ and $E=s\,e_\alpha$ on $T$ satisfies $E^2=1$;
\item $e$ is rigid: if $y\in\Vnat_2$ and $e\times y=y$ then $y=0$.
\end{enumerate}
\end{theorem}
\begin{leancode}
noncomputable def isingPlusMode (α : leech.carrier) (s : ℂ) : Module.End ℂ VPlus :=
  mode (leechVOAPlus.Y (isingVec α s).1) 1

noncomputable def isingTwistedMode (α : leech.carrier) (s : ℂ) : Module.End ℂ TP :=
  mode (plusModule.YM (isingVec α s).1) 1

noncomputable def isingSpectralPolynomial {B : Type*} [AddCommGroup B] [Module ℂ B]
    (X : Module.End ℂ B) : Module.End ℂ B :=
  X * (X - (2 : ℂ) • 1) * (X - (1 / 2 : ℂ) • 1) * (X - (1 / 16 : ℂ) • 1)

theorem isingVec_zeroMode_plus_annihilated {α : leech.carrier}
    (hα : leech.form α α = 4) (s : ℂ)
    (hs : s ^ 2 * latticeCocycle leech α (-α) = 1)
    (y : VPlus) (hy : y ∈ leechVOAPlus.grading 2) :
    isingSpectralPolynomial (isingPlusMode α s) y = 0 := by

noncomputable def isingChargeInvolution (α : leech.carrier) (s : ℂ) :
    Module.End ℂ TLeechT := s • eT leech leech_even leech_isUnit_det_gramMat α

theorem isingTwistedMode_hP {α : leech.carrier} (hα : leech.form α α = 4)
    (s : ℂ) (v : Fin (rank leech) → ℂ) (t : TLeechT) :
    isingTwistedMode α s (hP v t) = (1 / 32 : ℂ) •
      (hP (v + (2 * ∑ a, chargeCoord leech α a * v a) • chargeCoord leech α) t +
        hP (v - (2 * ∑ a, chargeCoord leech α a * v a) • chargeCoord leech α)
          (isingChargeInvolution α s t)) := by

theorem isingVec_zeroMode_twisted_annihilated {α : leech.carrier}
    (hα : leech.form α α = 4) (s : ℂ)
    (hs : s ^ 2 * latticeCocycle leech α (-α) = 1)
    (y : TP) (hy : y ∈ plusModule.weight (2 : ℂ)) :
    (isingTwistedMode α s * (isingTwistedMode α s - (1 / 16 : ℂ) • 1) *
      (isingTwistedMode α s - (1 / 2 : ℂ) • 1)) y = 0 := by

theorem isingVec_rigid {α : leech.carrier} (hα : leech.form α α = 4) (s : ℂ)
    (hs : s ^ 2 * latticeCocycle leech α (-α) = 1) :
    ∀ y : VnatVOA.grading 2,
      VnatVOA.griessMul ⟨isingVec α s, isingVec_mem_grading_two hα s⟩ y = y → y = 0 := by
\end{leancode}
(\leanpath{VOA/Moonshine/MoonshineLRFinOrderRigid.lean}, namespace \chtwentylean{VOA.Moonshine.ConwayNorton.M8}.)
\begin{dictionary}
\lean{mode f 1} applied to a vertex operator is the mode $u_{(1)}$ (the library's convention \lean{mode f n = f (-n-1)}
makes \lean{mode (Y u) n} equal to $u_{(n)}$); so \lean{isingPlusMode α s} is $o(e)$ restricted to the untwisted block
($e$ acts there through $\chtwentyVL$) and \lean{isingTwistedMode α s} is $o(e)$ on $\chtwentyTp$ (through the module map
\lean{plusModule.YM}).  \lean{(isingVec α s).1} is the $\chtwentyVL$-component of $e$.  In
\lean{isingSpectralPolynomial}, \lean{*} is composition of endomorphisms and \lean{(2 : ℂ) • 1} is $2\cdot\mathrm{id}$.  The
annihilation statements hold on the weight-two subspaces (\lean{hy}); the operators themselves act on all of
$\chtwentyVL$, $\chtwentyTp$.  \lean{eT … α} is the operator $e_\alpha$ on the $2^{12}$-dimensional space $T$
(Chapter~\ref{ch:06}), and \lean{isingChargeInvolution α s} $=E=s\,e_\alpha$.  In \lean{isingTwistedMode_hP} the vector
$v\pm(2\sum_a\alpha_av_a)\alpha$ is $(1\pm8P_\alpha)v$, since $\sum_a\alpha_av_a=\langle\alpha,v\rangle$ and $\langle\alpha,\alpha\rangle=4$.
\lean{isingVec_rigid} is rigidity in the Griess algebra $\Vnat_2$ (the subtype \lean{⟨isingVec α s, …⟩}), exactly the
hypothesis \lean{hrig} of Theorem~\ref{ch20:thm:periodic}.  \textbf{Strength.}  The spectral statements are slightly
stronger than needed (they locate all eigenvalues of $o(e)$ on $\Vnat_2$ in $\{0,2,\tfrac12,\tfrac1{16}\}$, Miyamoto's
list); only (c) is used.  \textbf{Could (c) be vacuous?}  No: it is a statement about the explicit operator $o(e)$ on the
$196884$-dimensional space $\Vnat_2$, and $e\ne0$.
\end{dictionary}

\paragraph{Proof on paper.}  (a) By Lemma~\ref{ch20:lem:span} it suffices to check the spanning vectors; write
$Q=Q(\alpha,\alpha)$, $X=\tfrac1{16}o(Q)+\tfrac s4o(x_\alpha)$.
\begin{itemize}
\item \emph{Quadratic block.}  For $u,v\in\mathfrak h$ write $u=\tfrac{a}{4}\alpha+u'$, $v=\tfrac b4\alpha+v'$ with $u',v'\perp\alpha$,
  $a=\langle\alpha,u\rangle$, $b=\langle\alpha,v\rangle$.  By Lemma~\ref{ch20:lem:products}: $XQ(u',v')=0$;
  $XQ(\alpha,v')=\tfrac12Q(\alpha,v')$; on $S_\alpha=\mathrm{span}\{Q,x_\alpha\}$, $XQ=Q+4sx_\alpha$ and
  $Xx_\alpha=x_\alpha+\tfrac{s\kappa_\alpha}4Q$, so $(X-1)^2=s^2\kappa_\alpha=1$ there, i.e.\ $X(X-2)=0$ on $S_\alpha$.
\item \emph{Charged block}, $\beta\in\chtwentyLf$, $\beta\ne\pm\alpha$, by the cases of Lemma~\ref{ch20:lem:products}(a):
  $\langle\alpha,\beta\rangle=0$: $Xx_\beta=0$; $\langle\alpha,\beta\rangle=\pm1$: $Xx_\beta=\tfrac1{16}x_\beta$ (from $\tfrac1{16}\langle\alpha,\beta\rangle^2$);
  $\langle\alpha,\beta\rangle=\pm2$: with $\beta'=\beta\mp\alpha\in\chtwentyLf$, $Xx_\beta=\tfrac14x_\beta+\tfrac{sc}4x_{\beta'}$ and
  $Xx_{\beta'}=\tfrac14x_{\beta'}+\tfrac{sc'}4x_\beta$ with $cc'=\kappa_\alpha$ (bimultiplicativity), so $X^2=\tfrac12X$ on
  $\mathrm{span}\{x_\beta,x_{\beta'}\}$.
\end{itemize}
Every case is a factor of $t(t-2)(t-\tfrac12)(t-\tfrac1{16})$.  (b) $o(Q)$ sends $h_v(-\tfrac12)\otimes t$ to
$h_{\frac12(1+8P_\alpha)v}(-\tfrac12)\otimes t$ (the zero mode of Lemma~\ref{ch20:lem:twisted} summed against
$\alpha_a\alpha_b$), and each of $o(e^{\pm\alpha})$ acts on level $\tfrac12$ by $2^{-4}h_{(1-8P_\alpha)v}(-\tfrac12)\otimes e_\alpha t$; with
the factors $\tfrac1{16}$ and $\tfrac s4$ this gives the displayed formula ($\tfrac s4\cdot2\cdot2^{-4}=\tfrac s{32}$).  $E^2=s^2e_\alpha^2=s^2\varepsilon(\alpha,\alpha)=1$.  On
$(1-P_\alpha)\mathfrak h\otimes T$, $32X=1+E\in\{0,2\}$; on $P_\alpha\mathfrak h\otimes T$, $32X=9-7E\in\{2,16\}$.  So
$X\in\{0,\tfrac1{16},\tfrac12\}$ on commuting eigenspaces of $E$.  (c) $o(e)$ preserves both blocks of $\Vnat_2$.  If
$Xy=y$, then $p(X)y=p(1)y$ for any polynomial $p$; with the polynomials of (a), (b), $p(1)=1\cdot(-1)\cdot\tfrac12\cdot\tfrac{15}{16}\ne0$
and $1\cdot\tfrac{15}{16}\cdot\tfrac12\ne0$, so $y=0$. \qed
\begin{physics}
The eigenvalues of $L'_0=o(e)$ on weight-two fields are the conformal weights $h\in\{0,\tfrac12,\tfrac1{16}\}$ of the three
irreducible representations of the Ising model, plus $h=2$ for $e$ itself: exactly what the representation theory of
the $c=\tfrac12$ Virasoro algebra predicts (Miyamoto 1996).  Here they are obtained by direct computation, without that
representation theory.  Since $1$ is not among them, the Ising sub-theory has no marginal deformations inside $\Vnat$.
\end{physics}
\begin{status}
\textbf{Unconditional} given the normalisation \lean{hs}.  Printed in \leanpath{VOA/Moonshine/MoonshineLRFinOrderAudit.lean};
standard axioms.
\end{status}

%======================================================================
\section{Every automorphism of the moonshine module has finite order}\sectionmark{Finite order}\label{ch20:sec:finorder}
%======================================================================

\subsection{From Ising vectors to the untwisted block}

\begin{lemma}[Ising vectors determine the untwisted block]\label{ch20:lem:fixising}
(a) If $g\in\Aut(\Vnat)$ fixes, for every minimal $\alpha$, both $e_{\alpha,s}$ and $e_{\alpha,-s}$ for some normalised $s$, then
$g$ fixes every $(q_{ab},0)$ and every $(x_\alpha,0)$, $\alpha\in\chtwentyLf$.  (b) If $g$ fixes every $(q_{ab},0)$ and every
$(x_\alpha,0)$, $\alpha\in\chtwentyLf$, then $g$ fixes $\chtwentyVL\times0$.  (c) If $g$ fixes $\chtwentyVL\times0$ and the
twisted generator $(0,x)$, then $g=1$.
\end{lemma}
\begin{leancode}
theorem aut_fix_quad_minimal_of_fix_ising (g : AutVnat)
    (hg : ∀ α : leech.carrier, leech.form α α = 4 →
      ∃ s : ℂ, s ^ 2 * latticeCocycle leech α (-α) = 1 ∧
        g (isingVec α s) = isingVec α s ∧ g (isingVec α (-s)) = isingVec α (-s)) :
    (∀ a b : Fin (rank leech), g ((qP a b, 0) : Vnat) = ((qP a b, 0) : Vnat)) ∧
      (∀ α : leech.carrier, leech.form α α = 4 →
        g ((frameX0P α, 0) : Vnat) = ((frameX0P α, 0) : Vnat)) := by

theorem aut_fix_untwisted_of_fix_quad_minimal (g : AutVnat)
    (hq : ∀ a b : Fin (rank leech), g ((qP a b, 0) : Vnat) = ((qP a b, 0) : Vnat))
    (hx : ∀ α : leech.carrier, leech.form α α = 4 →
      g ((frameX0P α, 0) : Vnat) = ((frameX0P α, 0) : Vnat)) :
    ∀ x : VPlus, g ((x, 0) : Vnat) = ((x, 0) : Vnat) := by

theorem aut_eq_one_of_fix_untwisted_twistedGen (g : AutVnat)
    (hU : ∀ x : VPlus, g ((x, 0) : Vnat) = ((x, 0) : Vnat))
    (hx : g ((0, xGen) : Vnat) = ((0, xGen) : Vnat)) : g = 1 := by
\end{leancode}
(\leanpath{VOA/Moonshine/MoonshineLRFinOrderRigid.lean}.)
\begin{leancode}
noncomputable abbrev xGen : TP := x0T (tF Dnat.s)
\end{leancode}
(\leanpath{VOA/Moonshine/AutSigma.lean}.)
\begin{dictionary}
\lean{xGen} is the twisted generator $x=h_0(-\tfrac12)\chtwentyvac\otimes t_s$ of the frame sector $T^+_s$ for the standard
data (\lean{x0T}, \lean{tF}; Chapters~\ref{ch:12}--\ref{ch:14}).  All three statements are as on paper.  (b) and (c) are
the two halves of Theorem~\ref{ch20:thm:faithful} with the hypothesis ``$g$ fixes $\Vnat_2$'' replaced by the weaker
hypotheses actually used.
\end{dictionary}
\paragraph{Proof on paper.}  (a) From the definition,
$x_\alpha=\tfrac2s(e_{\alpha,s}-e_{\alpha,-s})$ and $Q(\alpha,\alpha)=8(e_{\alpha,s}+e_{\alpha,-s})$ (in the untwisted block).  The
$q_{ab}$ are recovered from the $Q(\gamma,\gamma)$, $\gamma$ minimal, by polarisation inside the lattice: with $w_j$ the vector
$8e_j$ of the $\sqrt8$-model and $\gamma_\pm=4e_j\pm4e_k\in\chtwentyLf$ ($j\ne k$), $Q(w_j,w_k)=Q(\gamma_+,\gamma_+)-Q(\gamma_-,\gamma_-)$ and
$Q(w_j,w_j)+Q(w_k,w_k)=2(Q(\gamma_+,\gamma_+)+Q(\gamma_-,\gamma_-))$, so three indices isolate $Q(w_j,w_j)$, and the $q_{ab}$ lie in the
span of the $Q(w_j,w_k)$ (Chapter~\ref{ch:13}, \lean{qV_eq_sum_dV}).  (b) Steps 2--6 of the proof of
Theorem~\ref{ch20:thm:faithful}: they use the hypothesis on $\Vnat_2$ only at the $q_{ab}$ (to commute with $o(q_{ab})$)
and at the $x_\alpha$ with $\alpha$ minimal (to propagate scalars).  (c) $\chtwentyTp$ is spanned by the vectors
$u_{(n)}x$, $u\in\chtwentyVL$ (\lean{modeSpan_x0T}), and $g(u_{(n)}x)=(gu)_{(n)}(gx)=u_{(n)}x$. \qed
\begin{physics}
Two Ising models (the two stress tensors $\omega^\pm(\alpha)$) per minimal momentum pin down every cosine operator and every
bilinear; through the OPE they pin down the whole untwisted orbifold.
\end{physics}
\begin{status}
Unconditional; printed in the FinOrder audit; standard axioms.
\end{status}

\subsection{The theorem}

\begin{theorem}[Every automorphism of $\Vnat$ has finite order]\label{ch20:thm:finorder}
For every $g\in\Aut(\Vnat)$ there is $n\ge1$ with $g^n=1$.
\end{theorem}
\begin{leancode}
theorem aut_isingVec_periodic (g : AutVnat) {α : leech.carrier}
    (hα : leech.form α α = 4) (s : ℂ)
    (hs : s ^ 2 * latticeCocycle leech α (-α) = 1) :
    ∃ K : ℕ, 1 ≤ K ∧ (g ^ K) (isingVec α s) = isingVec α s := by

theorem aut_pow_fix_untwisted (g : AutVnat) :
    ∃ K : ℕ, 1 ≤ K ∧ ∀ x : VPlus, (g ^ K) ((x, 0) : Vnat) = ((x, 0) : Vnat) := by

theorem autVnat_isOfFinOrder : ∀ g : AutVnat, IsOfFinOrder g := by
\end{leancode}
(\leanpath{VOA/Moonshine/MoonshineLRFinOrderRigid.lean}, namespace \chtwentylean{VOA.Moonshine.ConwayNorton.M8}.)
\begin{dictionary}
\lean{IsOfFinOrder g} is Mathlib's ``$g$ has finite order'': $\exists n\ge1,\ g^n=1$ in the group \lean{AutVnat}
(\chtwentylean{isOfFinOrder_iff_pow_eq_one}).  \lean{g ^ K} is the $K$-th power in \lean{AutVnat} (composition).  The headline has
\textbf{no hypotheses}: the quantifier \lean{∀ g : AutVnat} is in the statement.  It is exactly the input
\lean{hfin} of the reductions of Chapter~\ref{ch:19}.  It is \emph{weaker} than finiteness of $\Aut(\Vnat)$ (an infinite
torsion group is conceivable a priori); finiteness is Theorem~\ref{ch20:thm:finite}.  \textbf{Could it be vacuous?}  Only if
\lean{AutVnat} were trivial, which it is not (it contains $\sigma$); and its strength does not depend on the size of the group.
\end{dictionary}
\paragraph{Proof on paper.}  Fix $g$.  (1) Let $A=g|_{\Vnat_2}$; it is a linear automorphism of the commutative Griess
algebra preserving the product (\lean{griessMul_aut}).  For each minimal $\alpha$ and $s=\pm c_\alpha$, where $c_\alpha$
is a fixed square root of $\varepsilon(\alpha,\alpha)$, $e_{\alpha,s}$ is a rigid idempotent (Theorems~\ref{ch20:thm:rigid} and
\lean{isingVec_one_self}), so by Theorem~\ref{ch20:thm:periodic} some $g^{k}$, $k\ge1$, fixes it
(\lean{aut_isingVec_periodic}).  There are finitely many minimal vectors (\lean{normFour_finite}), so the product
$K_1$ of all these periods works for all of them at once; by Lemma~\ref{ch20:lem:fixising}(a),(b), $g^{K_1}$ fixes
$\chtwentyVL\times0$ (\lean{aut_pow_fix_untwisted}).  (2) Apply (1) to $h=\sigma g\sigma^{-1}$: some $h^{K_2}=\sigma g^{K_2}\sigma^{-1}$
fixes $\chtwentyVL\times0$.  Since $\sigma(0,x)=(\Psi x,0)\in\chtwentyVL\times0$ (\lean{sigNat_Ts}), $g^{K_2}$ fixes
$\sigma^{-1}(\Psi x,0)=(0,x)$.  (3) $g^{K_1K_2}$ fixes $\chtwentyVL\times0$ and $(0,x)$, so it is $1$ by
Lemma~\ref{ch20:lem:fixising}(c). \qed

\begin{physics}
Every global symmetry of the Monster CFT is of finite order.  For a unitary CFT with finite-dimensional weight spaces one
expects the symmetry group to be compact (it preserves the positive inner product on each $V_n$, so it sits inside
$\prod_nU(V_n)$), with Lie algebra the conserved currents, i.e.\ $V_1$; since $\Vnat_1=0$ there are no continuous symmetries,
so the group is discrete and compact, hence finite.  The library's proof uses neither unitarity nor compactness: rigid
Ising sub-theories play the role of the missing currents, and the triality $\sigma$ carries the argument into the twisted
sector.
\end{physics}
\begin{status}
\textbf{Unconditional.}  Printed in \leanpath{VOA/Moonshine/MoonshineLRFinOrderAudit.lean}; \texttt{\#print axioms} $=$
\lean{[propext, Classical.choice, Quot.sound]}, also printed in \leanpath{VOA/Moonshine/MoonshineCNFinalAudit.lean}
(targeted build green, 0 \lean{sorryAx}).  Uses no group theory of the Monster.
\end{status}

\begin{warning}
\textbf{Two routes that do not work.}  The finite-order argument needs (a) some
vectors of $\Vnat_2$ that every automorphism permutes with finite orbits, and (b) a way to see that an automorphism fixing the
untwisted block fixes the twisted one.  Two tempting choices fail.  (a) Starting from $c=1$ Virasoro vectors
$\tfrac12h(-1)^2\chtwentyvac$ fails: they form a continuous $O(24)$-family, are not rigid, and need not be periodic under
an automorphism.  Only $c=\tfrac12$ rigidity starts the argument.  (b) Diagonal zero modes do not separate the twisted sector:
the eigenvalues of $o(q_{aa})$ on $h_{e_c}(-\tfrac12)\otimes t$ are $(\tfrac98$ at $a=c$, $\tfrac18$ elsewhere$)$, and on $x_n$
with $n$ odd of norm $4$ in a $V_3^+$ block they are $n_a^2/8$, the same pattern; any argument of the form ``a map commuting
with all $o(q_{aa})$ has no twisted-to-untwisted component'' is false.  The library uses the off-diagonal $o(q_{ab})$
(Lemma~\ref{ch20:lem:twisted}) and the $\sigma$-trick.  Also not used: a route through derivations of the
infinite-dimensional $\Vnat$ and a Jordan decomposition of $g$ on $\Vnat$; the finite-dimensional lemma of
\S\ref{ch20:sec:periodic} replaces it.
\end{warning}

%======================================================================
\section{The automorphism group of the moonshine module is finite}\sectionmark{Aut of the moonshine module is finite}\label{ch20:sec:finite}
%======================================================================

The Conway--Norton theorem needs only finite order.  Finiteness of the group was proved afterwards, as a
route-independent piece of the identification of $\Aut(\Vnat)$ with the Monster.  The proof again uses no group theory.

\subsection{Finitely many rigid idempotents}

\begin{theorem}[Rigid idempotents are finite in number; derivations vanish on their span]\label{ch20:thm:finrigid}
Let $B$ be a finite-dimensional complex vector space with a commutative bilinear product.  (a) The set of rigid
idempotents $\{e: e\cdot e=2e,\ (e\cdot y=y\Rightarrow y=0)\}$ is finite.  (b) If the rigid idempotents span $B$, every
derivation of $(B,\cdot)$ is zero.
\end{theorem}
\begin{leancode}
theorem finite_setOf_rigid_idempotent [FiniteDimensional ℂ B] :
    {e : B | mul e e = (2 : ℂ) • e ∧ ∀ y, mul e y = y → y = 0}.Finite := by

theorem derivation_eq_zero_of_span_rigid [FiniteDimensional ℂ B]
    (hspan : Submodule.span ℂ
      {e : B | mul e e = (2 : ℂ) • e ∧ ∀ y, mul e y = y → y = 0} = ⊤)
    (D : Module.End ℂ B)
    (hD : ∀ x y, D (mul x y) = mul (D x) y + mul x (D y)) : D = 0 := by
\end{leancode}
(\leanpath{VOA/Moonshine/MoonshineLRFiniteRigid.lean}, namespace \chtwentylean{VOA.Moonshine.ConwayNorton.M8}; section variables
\lean{(mul : B →ₗ[ℂ] B →ₗ[ℂ] B) (hcomm : ∀ x y, mul x y = mul y x)}, both included in each statement by
\lean{include hcomm in}.)
\begin{dictionary}
\lean{S.Finite} is Mathlib's \lean{Set.Finite}: $S$ is a finite set.  The set in (a) is exactly the set of rigid
idempotents of Definition~\ref{ch20:def:rigid}; it is \emph{not} restricted to explicit Ising vectors, nor to vectors of any
special shape.  In (b), \lean{hspan} says the rigid idempotents span $B$, \lean{hD} that $D$ is a derivation.  Statements
exactly as on paper; the file docstring stresses that the product need be neither associative nor unital.
\end{dictionary}
\paragraph{Proof on paper.}  (a) Choose a basis and write the product with structure constants $c_{ijk}$.  The
idempotent equations $f_i(X)=\sum_{j,k}c_{ijk}X_jX_k-2X_i$ generate an ideal $J\subseteq R=\C[X_1,\dots,X_n]$; idempotents
are the points of the affine variety $V(J)$.  Let $p$ be a rigid idempotent and $\mathfrak m_p$ its maximal ideal.  Taylor
expansion at $p$ gives $f(p+y)=2(L_p-1)y+(\text{quadratic in }y)$ (using commutativity, the linear part of
$(p+y)\cdot(p+y)$ is $2p\cdot y$); the Jacobian $2(L_p-1)$ is invertible by rigidity, so each coordinate $X_i-p_i$ is a linear
combination of the $f_j$ modulo $\mathfrak m_p^2$: $\mathfrak m_p\subseteq J+\mathfrak m_p^2$.  Hence $\mathfrak m_p$ is a
\emph{minimal} prime over $J$: if $J\subseteq Q\subseteq\mathfrak m_p$ with $Q$ prime, the image $N$ of $\mathfrak m_p$ in the
Noetherian domain $R/Q$ satisfies $N=N^2$, so $N=0$ by Nakayama (Krull), i.e.\ $Q=\mathfrak m_p$.  A Noetherian ring has
finitely many minimal primes over an ideal (\chtwentylean{Ideal.finite_minimalPrimes_of_isNoetherianRing}), and $p\mapsto\mathfrak m_p$
is injective.  (b) Each rigid idempotent is killed by every derivation (step (iii) of Theorem~\ref{ch20:thm:periodic}), and
they span. \qed
\begin{physics}
Rigid sub-theories are isolated points of the ``moduli space of embedded stress tensors'': they cannot be deformed
(\S\ref{ch20:sec:periodic}), so they come in a discrete, and by algebraic geometry finite, set.  A symmetry can only permute
them.
\end{physics}
\begin{status}
\textbf{Unconditional}, generic (Mathlib only).  Printed in \leanpath{VOA/Moonshine/MoonshineLRFinOrderAudit.lean};
standard axioms.
\end{status}

\begin{warning}
\textbf{A version of this argument that does not work.}  One could try to
count
\emph{all Ising vectors} of $\Vnat$, which needs rigidity at every Ising vector, not only at the explicit ones; for an
arbitrary Ising vector the only known proof goes through $c=\tfrac12$ representation theory (vanishing of a singular vector,
needing unitarity or the classification of Ising vectors).  The version in the library sidesteps this by counting
\emph{rigid} idempotents, rigidity being part of the definition of the set, and by using only the explicit Ising vectors
for faithfulness (Theorem~\ref{ch20:thm:finite}).  \textbf{Also not done:} part (b) is \emph{not} applied to $\Vnat_2$; that
would need ``the rigid idempotents span $\Vnat_2$'', which the library does not prove.  So $\mathrm{Der}(\Vnat_2)=0$ and
$\mathrm{Der}(\Vnat)=0$ are not theorems of the library.
\end{warning}

\subsection{\texorpdfstring{$\Aut(\Vnat)$}{Aut(V)} acts faithfully on a finite set}

\begin{definition}[The rigid idempotents of the Griess algebra]\label{ch20:def:rigidset}
$R(\Vnat)\subseteq\Vnat_2$ is the set of rigid idempotents of the Griess algebra $(\Vnat_2,\times)$.
\end{definition}
\begin{leancode}
def vnatGriessBilin :
    VnatVOA.grading 2 →ₗ[ℂ] VnatVOA.grading 2 →ₗ[ℂ] VnatVOA.grading 2 :=
  LinearMap.mk₂ ℂ VnatVOA.griessMul VnatVOA.griessMul_add_left
    (fun c x y => Subtype.ext (VnatVOA.smul_smul_left c x 1 y))
    VnatVOA.griessMul_add_right
    (fun c x y => Subtype.ext (VnatVOA.smul_smul_right x c 1 y))

def vnatRigidSet : Set (VnatVOA.grading 2) :=
  {e | vnatGriessBilin e e = (2 : ℂ) • e ∧ ∀ y, vnatGriessBilin e y = y → y = 0}

theorem autRestrict_mem_vnatRigidSet (g : AutVnat) {e : VnatVOA.grading 2}
    (he : e ∈ vnatRigidSet) : GriessAut.autRestrict g 2 e ∈ vnatRigidSet := by

theorem isingVec_mem_vnatRigidSet {α : leech.carrier} (hα : leech.form α α = 4) (s : ℂ)
    (hs : s ^ 2 * latticeCocycle leech α (-α) = 1) :
    (⟨isingVec α s, isingVec_mem_grading_two hα s⟩ : VnatVOA.grading 2) ∈ vnatRigidSet :=
\end{leancode}
(\leanpath{VOA/Moonshine/MoonshineLRFiniteAut.lean}, namespace \chtwentylean{VOA.Moonshine.ConwayNorton.M8}.)
\begin{dictionary}
\lean{vnatGriessBilin} packages \lean{griessMul} as a bilinear map (\lean{LinearMap.mk₂} with the four linearity proofs),
so that the generic theorems apply; \lean{vnatGriessBilin_apply} states it is \lean{griessMul} by \lean{rfl}.
\lean{vnatRigidSet} is $R(\Vnat)$ verbatim.  \chtwentylean{autRestrict_mem_vnatRigidSet}: $\Aut(\Vnat)$ permutes $R(\Vnat)$.
\lean{isingVec_mem_vnatRigidSet}: the normalised explicit Ising vectors lie in it (this is
\lean{isingVec_one_self} and \lean{isingVec_rigid} together).  The conformal vector $\omega$ also lies in $R(\Vnat)$
($o(\omega)=L_0=2$ on $\Vnat_2$), although the library does not state this.
\end{dictionary}
\begin{physics}
The finite set of rigid embedded stress tensors of weight two, permuted by every symmetry of the theory.
\end{physics}
\begin{status}
Definitions and unconditional lemmas; printed in the FinOrder audit; standard axioms.
\end{status}

\begin{theorem}[$\Aut(\Vnat)$ is finite]\label{ch20:thm:finite}
(a) An automorphism of $\Vnat$ that fixes every rigid idempotent of $\Vnat_2$ is the identity.  (b) If $R(\Vnat)$ is finite,
$\Aut(\Vnat)$ is finite.  (c) $\Aut(\Vnat)$ is finite.
\end{theorem}
\begin{leancode}
theorem eq_one_of_fix_rigid (k : AutVnat)
    (hk : ∀ e ∈ vnatRigidSet, GriessAut.autRestrict k 2 e = e) : k = 1 := by

theorem finite_autVnat_of_finite_rigid (hR : vnatRigidSet.Finite) : Finite AutVnat := by

theorem finite_autVnat : Finite AutVnat := by
\end{leancode}
(\leanpath{VOA/Moonshine/MoonshineLRFiniteAut.lean}.)
\begin{dictionary}
\lean{Finite AutVnat} is Mathlib's type class ``the type \lean{AutVnat} has finitely many elements'' (there is a bijection
with some \lean{Fin n}).  \lean{GriessAut.autRestrict k 2 e} is $k(e)$ for $e\in\Vnat_2$.  (b) is conditional on \lean{hR};
(c) has \textbf{no hypotheses}: it applies (b) with \lean{hR} supplied by Theorem~\ref{ch20:thm:finrigid}(a) for the Griess
algebra (finite-dimensional because $\Vnat_2$ is, \lean{gradingNat 2}).  This is stronger than
Theorem~\ref{ch20:thm:finorder} and gives no bound on the order: the library does not compute $|\Aut(\Vnat)|$.
\end{dictionary}
\paragraph{Proof on paper.}  (a) If $k$ fixes $R(\Vnat)$ pointwise, it fixes every normalised $e_{\alpha,\pm s}$, so by
Lemma~\ref{ch20:lem:fixising}(a),(b) it fixes $\chtwentyVL\times0$.  Since $\sigma^{-1}$ permutes $R(\Vnat)$, the conjugate
$\sigma k\sigma^{-1}$ also fixes $R(\Vnat)$ pointwise, hence fixes $\chtwentyVL\times0$, hence $k$ fixes $\sigma^{-1}(\Psi x,0)=(0,x)$.
By Lemma~\ref{ch20:lem:fixising}(c), $k=1$.  (b) The map $\Aut(\Vnat)\to\{\text{maps }R(\Vnat)\to R(\Vnat)\}$, $g\mapsto g|_{R}$, is
injective by (a) (if $g|_R=h|_R$ then $h^{-1}g$ fixes $R$); a finite set has finitely many self-maps.  (c) (b) and
Theorem~\ref{ch20:thm:finrigid}(a). \qed
\begin{physics}
The global symmetry group of the Monster CFT is finite: it acts faithfully by permutations on the finite set of rigid
weight-two stress tensors.  (Classically one deduces this from $\Aut(\Vnat)=\mathbb M$; here it comes first, and the
identification with AxiomMath's Monster is the subject of Chapters~\ref{ch:21}--\ref{ch:22}.)
\end{physics}
\begin{status}
(c) \textbf{unconditional}; (b) conditional on \lean{hR}, discharged in (c).  Printed in
\leanpath{VOA/Moonshine/MoonshineLRFinOrderAudit.lean}; \texttt{\#print axioms} $=$ \lean{[propext, Classical.choice, Quot.sound]}
(build record of the release).  Not used by \lean{conwayNorton}.  No group
theory is used.
\end{status}

%======================================================================
\section{The Monster Lie algebra is generated by its simple root spaces}\sectionmark{Generation by simple root spaces}\label{ch20:sec:generation}
%======================================================================

\subsection{What has to be proved, and why it is not a Lie-theoretic statement}

Chapter~\ref{ch:18} constructs the Monster Lie algebra $\chtwentym$ (the physical states of $\Vnat\otimes V_{I\!I_{1,1}}$
modulo null states), graded by $I\!I_{1,1}$, with root spaces $\chtwentym_{(m,n)}\cong\Vnat_{1+mn}$ in Borcherds' coordinates,
a Cartan subalgebra $\chtwentym_{(0,0)}\cong\C^2$, and an invariant bilinear form; and it proves that these data form
\emph{Borcherds data} (a Borcherds--Kac--Moody algebra of rank two) \emph{provided} $\chtwentym$ is generated by its Cartan
subalgebra and its simple root spaces.  Every consequence that Chapter~\ref{ch:19} draws (the free positive part, the
denominator identity, the twisted denominator identities) passes through that structure.  This section discharges the
proviso.

The simple roots, in Borcherds' coordinates $(m,n)\in\Z^2$ with norm $-2mn$, are the real root $(1,-1)$ (root space
$\Vnat_0=\C\chtwentyvac$) and the imaginary roots $(1,j)$, $j\ge1$ (root space $\Vnat_{j+1}$, of dimension $c(j)$, the
coefficient of $q^j$ in $J$).  An important fact: \emph{given
the other Borcherds data, generation of $\chtwentym$ by these simple root spaces is equivalent to the product formula}
\begin{equation}\label{ch20:eq:knz}
  \prod_{m,n\ge1}(1-p^mq^n)^{c(mn)}=1-\sum_{a,b\ge1}c(a+b-1)\,p^aq^b
\end{equation}
(the Koike--Norton--Zagier identity, ``KNZ'', which is $p^{-1}(1-p/q)\prod_{m,n\ge1}(1-p^mq^n)^{c(mn)}=J(p)-J(q)$
rewritten).  So a proof of generation \emph{is} a proof of a statement about the modular function $J$, and it cannot come
from Lie theory or vertex algebra alone.  The library proves \eqref{ch20:eq:knz} analytically (Chapter~\ref{ch:19},
\lean{knz_product}) and derives generation from it by a dimension count.

\subsection{The objects from Chapters~\ref{ch:18} and~\ref{ch:19}}

We quote the declarations that the statements of this section use; their full presentation is in the chapters named.

\begin{leancode}
def borcherdsRealRoot : ℤ × ℤ := (1, -1)

def borcherdsImaginaryRoot (j : ℤ) : ℤ × ℤ := (1, j)

def borcherdsPositiveCone (γ : ℤ × ℤ) : Prop := 1 ≤ γ.1 ∧ 1 ≤ γ.2

def borcherdsPositiveRoot (γ : ℤ × ℤ) : Prop :=
  γ = borcherdsRealRoot ∨ borcherdsPositiveCone γ

def borcherdsRootSupport (γ : ℤ × ℤ) : Prop :=
  γ = 0 ∨ borcherdsPositiveRoot γ ∨ borcherdsPositiveRoot (-γ)

def borcherdsSimpleRoot (γ : ℤ × ℤ) : Prop :=
  γ = borcherdsRealRoot ∨ ∃ j : ℤ, 1 ≤ j ∧ γ = borcherdsImaginaryRoot j

def borcherdsWeightVector (γ : ℤ × ℤ) : ℂ × ℂ := (γ.1, γ.2)

def borcherdsCartanPairing (ξ μ : ℂ × ℂ) : ℂ := -(ξ.1 * μ.2 + μ.1 * ξ.2)
\end{leancode}
(\leanpath{VOA/Moonshine/MoonshineCNBorcherdsRoots.lean}, namespace \chtwentylean{VOA.Moonshine.ConwayNorton}.)
\begin{leancode}
def borcherdsSimpleGenerators (wt : ℤ × ℤ → Submodule ℂ g) : Set g :=
  {x | x ∈ wt 0 ∨ ∃ α, borcherdsSimpleRoot α ∧ (x ∈ wt α ∨ x ∈ wt (-α))}

structure BorcherdsData (g : Type*) [LieRing g] [LieAlgebra ℂ g] where
  ...
  generated : LieSubalgebra.lieSpan ℂ g (borcherdsSimpleGenerators weightSpace) = ⊤
\end{leancode}
(\leanpath{VOA/Moonshine/MoonshineCNBorcherdsData.lean}.)
\begin{leancode}
def monsterLieBorcherdsWeight (γ : ℤ × ℤ) : Submodule ℂ monsterLieAlgebra :=
  monsterLieMomentum (borcherdsMomentum γ)

def MonsterLieSimpleRootGeneration : Prop :=
  LieSubalgebra.lieSpan ℂ monsterLieAlgebra (borcherdsSimpleGenerators monsterLieBorcherdsWeight) =
    ⊤

def monsterLieBorcherdsData (hgen : MonsterLieSimpleRootGeneration) :
    BorcherdsData monsterLieAlgebra where

def borcherdsMomentum (γ : ℤ × ℤ) : lorentzianLattice.carrier := ![γ.1, -γ.2]

def monsterLieBorcherdsRootEquiv (γ : ℤ × ℤ) (hγ : γ ≠ 0) :
    monsterLieBorcherdsWeight γ ≃ₗ[ℂ] VnatVOA.grading (1 + γ.1 * γ.2) :=
\end{leancode}
(\leanpath{VOA/Moonshine/MoonshineCNMonsterBorcherdsData.lean}; \lean{monsterLieAlgebra} is
\chtwentylean{moonshineMatterNearVOA.primaryLieQuotient}, \leanpath{VOA/Moonshine/MoonshineCNMonsterLie.lean}.)
\begin{dictionary}
\textbf{Roots.}  A weight is a pair $\gamma=(m,n)\in\Z^2$ (\lean{γ.1}, \lean{γ.2}); \lean{borcherdsCartanPairing} is
the $\C$-bilinear extension of $\langle(m,n),(m',n')\rangle=-(mn'+m'n)$, so $\langle\gamma,\gamma\rangle=-2mn$.  The root
support is $0$, $\pm(1,-1)$ and $\pm$(open positive cone $m,n\ge1$); simple roots are $(1,-1)$ and $(1,j)$, $j\ge1$.
\textbf{Generators.}  \lean{borcherdsSimpleGenerators wt} is the \emph{set} $\mathfrak g_0\cup\bigcup_{\alpha\text{
simple}}(\mathfrak g_\alpha\cup\mathfrak g_{-\alpha})$; \lean{LieSubalgebra.lieSpan ℂ g S} is the Lie subalgebra generated by $S$;
\lean{= ⊤} says it is all of $\mathfrak g$.  \textbf{Borcherds data.}  The fields are: a grading by $\Z^2$
(\lean{DirectSum.IsInternal}: $\mathfrak g=\bigoplus_\gamma\mathfrak g_\gamma$), finite-dimensional pieces, brackets add
degrees, a two-dimensional Cartan $\mathfrak g_0$ acting on $\mathfrak g_\gamma$ by the pairing with $\gamma$, a symmetric
invariant form pairing $\mathfrak g_\gamma$ only with $\mathfrak g_{-\gamma}$, nondegenerately, and equal to the pairing on
the Cartan, roots in the support, one-dimensional real root spaces, and \lean{generated}.  \textbf{The Monster.}
\lean{monsterLieBorcherdsWeight γ} is the root space of $\chtwentym$ at the library momentum \lean{borcherdsMomentum γ}
$=(m,-n)$ (the library's lattice $I\!I_{1,1}$ has norm $+2xy$; Borcherds' $(m,n)$ is the library's $(m,-n)$; trap recorded in
Chapter~\ref{ch:16}, \S\ref{ch16:sec:traps}).  \chtwentylean{monsterLieBorcherdsRootEquiv} is the no-ghost isomorphism
$\chtwentym_{(m,n)}\cong\Vnat_{1+mn}$ for $(m,n)\ne0$ (Chapter~\ref{ch:18}).  \chtwentylean{MonsterLieSimpleRootGeneration} is
precisely ``$\chtwentym$ is generated by $\chtwentym_0$ and the $\chtwentym_{\pm\alpha}$, $\alpha$ simple''.
\end{dictionary}
\begin{physics}
$\chtwentym$ is the Lie algebra of physical states (BRST cohomology at ghost number one) of the bosonic string in $26$
dimensions with $24$ dimensions compactified on the moonshine orbifold and two on a Lorentzian torus; the momentum
$(m,n)\in I\!I_{1,1}$ is the light-cone momentum/winding, and the physical states at momentum $(m,n)$ are counted by the
transverse states $\Vnat_{1+mn}$ (no-ghost).  The real simple root $(1,-1)$ is the tachyon-like state $\Vnat_0$; the
imaginary simple roots $(1,j)$ are the states of one unit of $m$.
\end{physics}
\begin{status}
Definitions (Chapter~\ref{ch:18}); the root definitions and \lean{BorcherdsData} contain no VOA input.
\end{status}

\begin{leancode}
noncomputable def coneLog (e : ℕ → ℕ → ℂ) : PowerSeries (LaurentSeries ℂ) :=
  PowerSeries.mk (coneLogCoeff e)

noncomputable def coneProd (e : ℕ → ℕ → ℂ) : PowerSeries (LaurentSeries ℂ) :=
  odeSol (coneLog e)

theorem knz_product : coneProd (fun m n => cJ (m * n)) = knzSeries cJ := by

theorem coneProd_exponent_unique {e₁ e₂ : ℕ → ℕ → ℂ} (h : coneProd e₁ = coneProd e₂) :
    ∀ m n : ℕ, 0 < m → 0 < n → e₁ m n = e₂ m n := by
\end{leancode}
(\leanpath{VOA/Moonshine/MoonshineLRM5Product.lean}.)
\begin{leancode}
noncomputable def knzCoeff (c : ℕ → ℂ) (a : ℕ) : LaurentSeries ℂ :=
  if a = 0 then 1 else
    -HahnSeries.ofPowerSeries ℤ ℂ (PowerSeries.mk fun b => if b = 0 then 0 else c (a + b - 1))

noncomputable def knzSeries (c : ℕ → ℂ) : PowerSeries (LaurentSeries ℂ) :=
  PowerSeries.mk (knzCoeff c)
\end{leancode}
(\leanpath{VOA/Moonshine/MoonshineLRM5KNZ.lean}.)
\begin{leancode}
noncomputable def cJ (n : ℕ) : ℂ := PowerSeries.coeff (n + 1) JShift
\end{leancode}
(\leanpath{VOA/Moonshine/MoonshineLRM5Replicable.lean}.)
\begin{dictionary}
The ring is \lean{PowerSeries (LaurentSeries ℂ)} $=\C((q))[[p]]$.  \lean{coneLog e} is
$-\sum_{k,m,n\ge1}m\,e(m,n)\,p^{mk}q^{nk}$ (the $p\partial_p$-logarithmic derivative of the product), and \lean{coneProd e}
is the unique power series with constant term $1$ solving $p\,\partial_pF=F\cdot\mathrm{coneLog}(e)$ (\lean{odeSol},
\lean{coneProd_spec}); for natural exponents it is the formal product $\prod_{m,n\ge1}(1-p^mq^n)^{e(m,n)}$, and this
definition makes sense for complex (in particular negative) exponents without logarithms.  \lean{knzSeries c} is
$1-\sum_{a,b\ge1}c(a+b-1)p^aq^b$.  \lean{cJ n} is the coefficient of $q^{n+1}$ in \lean{JShift} $=qJ$, i.e.\ the coefficient
$c(n)$ of $q^n$ in $J$ ($c(0)=0$, $c(1)=196884$); equivalently $c(n)=\dim\Vnat_{n+1}$.  \lean{knz_product} is
\eqref{ch20:eq:knz}.  \lean{coneProd_exponent_unique}: the exponents of a cone product are determined by the product (on the
open cone).
\end{dictionary}
\begin{physics}
The product $\prod(1-p^mq^n)^{c(mn)}$ is the ``second-quantised'' partition function of the string states counted by
$\Vnat$; that it equals $p(J(p)-J(q))/(1-p/q)$ is the denominator identity of $\chtwentym$, a modular identity
(heuristically related to Hecke operators acting on $J$; Chapter~\ref{ch:19}).
\end{physics}
\begin{status}
Proved in Chapter~\ref{ch:19} (\lean{knz_product}: unconditional, standard axioms, in the M5 audit).
\end{status}

\subsection{Pre-data and the generic generation theorem}

\begin{definition}[Borcherds pre-data; the subalgebra generated by simple roots; the KNZ hypothesis]\label{ch20:def:predata}
\emph{Borcherds pre-data} on a complex Lie algebra $\mathfrak g$ are all the fields of Borcherds data except
\lean{generated}.  For pre-data $P$, $G=P.\mathrm{simpleSpan}$ is the Lie subalgebra generated by the Cartan and the simple
root spaces.  $P$ satisfies \emph{KNZ} if
\begin{equation*}
  \prod_{m,n\ge1}(1-p^mq^n)^{\dim\mathfrak g_{(m,n)}}=1-\sum_{a,b\ge1}\dim\mathfrak g_{(1,a+b-1)}\,p^aq^b .
\end{equation*}
\end{definition}
\begin{leancode}
structure BorcherdsPreData (g : Type*) [LieRing g] [LieAlgebra ℂ g] where
  weightSpace : ℤ × ℤ → Submodule ℂ g
  weight_internal : DirectSum.IsInternal weightSpace
  weight_finite : ∀ γ, FiniteDimensional ℂ (weightSpace γ)
  bracket_weight : ∀ {γ δ} {x y : g},
    x ∈ weightSpace γ → y ∈ weightSpace δ → ⁅x, y⁆ ∈ weightSpace (γ + δ)
  cartanMap : (ℂ × ℂ) →ₗ[ℂ] g
  cartan_injective : Function.Injective cartanMap
  cartan_range : LinearMap.range cartanMap = weightSpace 0
  cartan_action : ∀ (ξ : ℂ × ℂ) (γ : ℤ × ℤ) (x : g), x ∈ weightSpace γ →
    ⁅cartanMap ξ, x⁆ = borcherdsCartanPairing ξ (borcherdsWeightVector γ) • x
  form : LinearMap.BilinForm ℂ g
  form_symmetric : ∀ x y, form x y = form y x
  form_invariant : ∀ x y z, form ⁅x, y⁆ z = form x ⁅y, z⁆
  form_orthogonal : ∀ γ δ, γ + δ ≠ 0 → ∀ x ∈ weightSpace γ,
    ∀ y ∈ weightSpace δ, form x y = 0
  form_cartan : ∀ ξ μ, form (cartanMap ξ) (cartanMap μ) = borcherdsCartanPairing ξ μ
  form_nondegenerate : ∀ γ, ∀ x ∈ weightSpace γ,
    (∀ y ∈ weightSpace (-γ), form x y = 0) → x = 0
  root_support : ∀ γ, weightSpace γ ≠ ⊥ → borcherdsRootSupport γ
  real_dimension : Module.finrank ℂ (weightSpace borcherdsRealRoot) = 1 ∧
    Module.finrank ℂ (weightSpace (-borcherdsRealRoot)) = 1

def simpleSpan : LieSubalgebra ℂ g :=
  LieSubalgebra.lieSpan ℂ g (borcherdsSimpleGenerators P.weightSpace)

def KNZ : Prop :=
  coneProd (fun m n => (finrank ℂ (P.weightSpace ((m : ℤ), (n : ℤ))) : ℂ)) =
    knzSeries fun j => (finrank ℂ (P.weightSpace (1, (j : ℤ))) : ℂ)
\end{leancode}
(\leanpath{VOA/Moonshine/MoonshineLRGenerationPreData.lean};\newline \lean{KNZ} is in
\leanpath{VOA/Moonshine/MoonshineLRGenerationSandwich.lean}; namespace \chtwentylean{VOA.Moonshine.ConwayNorton.M8.BorcherdsPreData}.)
\begin{dictionary}
The fields of \lean{BorcherdsPreData} are, word for word, those of \lean{BorcherdsData} minus \lean{generated}; the
definition \lean{toData} (same file) adds a generation proof back.  \lean{simpleSpan} is $G$.  In \lean{KNZ},
\lean{finrank ℂ (P.weightSpace ((m : ℤ), (n : ℤ)))} is $\dim\mathfrak g_{(m,n)}$ cast to $\C$, and the right side is
\lean{knzSeries} of $j\mapsto\dim\mathfrak g_{(1,j)}$ (the value at $j=0$ is irrelevant: \lean{knzCoeff} only uses
$c(a+b-1)$ with $a,b\ge1$).  So \lean{KNZ} is the displayed identity.
\end{dictionary}
\begin{physics}
Pre-data are the string-theory data (graded physical states, momentum zero-modes, the Zamolodchikov-type invariant pairing)
before one knows which states are ``elementary''.  KNZ is the statement that the free second-quantised spectrum generated by
the simple-root states reproduces the full spectrum.
\end{physics}
\begin{status}
Definitions; printed in \leanpath{VOA/Moonshine/MoonshineLRGenerationAudit.lean}; standard axioms.
\end{status}

\begin{theorem}[Generation from KNZ]\label{ch20:thm:genKNZ}
Borcherds pre-data satisfying KNZ are generated by their simple root spaces: $G=\mathfrak g$.
\end{theorem}
\begin{leancode}
theorem simpleSpan_eq_top_of_KNZ (hK : P.KNZ) : P.simpleSpan = ⊤ := by

theorem lieSpan_eq_top_of_KNZ (hK : P.KNZ) :
    LieSubalgebra.lieSpan ℂ g (borcherdsSimpleGenerators P.weightSpace) = ⊤ :=
  P.simpleSpan_eq_top_of_KNZ hK
\end{leancode}
(\leanpath{VOA/Moonshine/MoonshineLRGenerationSandwich.lean}; section variable \lean{(P : BorcherdsPreData g)} for an
arbitrary complex Lie algebra \lean{g}.)
\begin{dictionary}
The single hypothesis \lean{hK} is KNZ for $P$.  The conclusion is the \lean{generated} field, so \lean{P.toData} with this
proof is Borcherds data.  Generic: no VOA, no Monster.
\end{dictionary}

\paragraph{Proof on paper (the sandwich).}  Let $G$ be the subalgebra generated by the Cartan and the simple root
spaces, $R=\{x\in G:(x,G)=0\}$ the radical of the restricted form, and $\mathfrak g'=G/R$.
\begin{enumerate}[label=(\arabic*)]
\item \emph{$G$ and $R$ are graded.}  The generators are homogeneous, so $G=\bigoplus_\gamma(G\cap\mathfrak g_\gamma)$.  $R$ is
  an ideal by invariance, and homogeneous because the form pairs degree $\gamma$ only with $-\gamma$
  (\lean{form_decompose_eq}).  Hence $\mathfrak g'$ is graded by the images $\mathfrak g'_\gamma$, and the form descends to a
  form nondegenerate on $\mathfrak g'_\gamma\times\mathfrak g'_{-\gamma}$.
\item \emph{$\mathfrak g'$ is Borcherds data, generated by construction} (\lean{quotientData}).  If $\mathfrak g_\gamma$ and
  $\mathfrak g_{-\gamma}$ both lie in $G$ (true for the Cartan, the real roots and the simple imaginary roots), then
  $R\cap\mathfrak g_\gamma=0$ by nondegeneracy, so $\mathfrak g'_\gamma\cong\mathfrak g_\gamma$ (\lean{finrank_quotientWeight}).
\item \emph{The character identity for any Borcherds data} (\chtwentylean{coneProd_neg_mul_knzSeries}):
  \begin{equation*}
    \prod_{m,n\ge1}(1-p^mq^n)^{-\dim\mathfrak g'_{(m,n)}}\cdot\Bigl(1-\sum_{a,b\ge1}\dim\mathfrak g'_{(1,a+b-1)}p^aq^b\Bigr)=1 .
  \end{equation*}
  Reason: by Chapter~\ref{ch:18}, the positive part $\mathfrak u^+=\bigoplus_{m,n\ge1}\mathfrak g'_{(m,n)}$ is free on the
  generating space $U=\bigoplus_{j\ge1}\bigoplus_{0\le k<j}(\mathrm{ad}\,e)^k\mathfrak g'_{(1,j)}$ (the real-root strings through
  the imaginary simple roots), so $U(\mathfrak u^+)\cong T(U)$ (\chtwentylean{positiveTensorGeneratorToEnveloping_bijective}).  The
  graded dimension of $T(U)$ is $1/(1-\mathrm{ch}\,U)$, that of $U(\mathfrak u^+)$ is $\prod(1-p^mq^n)^{-\dim\mathfrak u^+_{(m,n)}}$
  (graded PBW), and $\dim U_{(a,b)}=\dim\mathfrak g'_{(1,a+b-1)}$.  For the last point: $\mathrm{ad}\,e$ moves the weight by the
real root $(1,-1)$, so the string through $(1,j)$ has weights $(1+k,j-k)$; it is cut off at $k<j$ because $(1+k,j-k)$ must stay
in the open cone (for $k=j$ the weight $(1+j,0)$ is not a root), and $\mathrm{ad}\,e$ is injective on these pieces by
$\mathfrak{sl}_2$-theory (the imaginary simple root spaces are lowest-weight spaces for the real $\mathfrak{sl}_2$, of weight
$j-1$, so the string has exactly $j$ terms).  A weight $(a,b)$ with $a,b\ge1$ occurs exactly once, for $j=a+b-1$, $k=a-1$.
\item \emph{Comparison.}  By (2), the simple imaginary root spaces of $\mathfrak g'$ and $\mathfrak g$ have the same
  dimensions, so the second factor in (3) is \lean{knzSeries} of $j\mapsto\dim\mathfrak g_{(1,j)}$, which by KNZ for $P$ equals
  $\prod(1-p^mq^n)^{\dim\mathfrak g_{(m,n)}}$.  Hence $\prod(1-p^mq^n)^{\dim\mathfrak g_{(m,n)}-\dim\mathfrak g'_{(m,n)}}=1$, and by
  uniqueness of exponents (\lean{coneProd_exponent_unique}, \chtwentylean{eq_of_coneProd_neg_mul_eq_one})
  $\dim\mathfrak g'_\gamma=\dim\mathfrak g_\gamma$ on the open positive cone (\chtwentylean{finrank_quotientWeight_of_KNZ}).
\item \emph{The sandwich.}  For $\gamma$ in the cone, $\dim\mathfrak g_\gamma=\dim\mathfrak g'_\gamma\le\dim(G\cap\mathfrak g_\gamma)\le
  \dim\mathfrak g_\gamma$, so $G\supseteq\mathfrak g_\gamma$ (\chtwentylean{weightSpace_le_simpleSpan_of_KNZ}).
\item \emph{The negative cone.}  The opposite pre-data (grading $\gamma\mapsto-\gamma$, Cartan composed with $-1$;
  \lean{opposite}) have the same $G$ and satisfy KNZ, because $\dim\mathfrak g_{-\gamma}=\dim\mathfrak g_\gamma$ (the form pairs
  them nondegenerately); applying (5) to them gives $G\supseteq\mathfrak g_{-\gamma}$.
\item The remaining root spaces ($0$, $\pm(1,-1)$) are generators.  By the root support, $G$ contains every
  $\mathfrak g_\gamma$, hence $G=\mathfrak g$. \qed
\end{enumerate}
\begin{physics}
The sandwich is a counting argument familiar from string theory: the states built freely from the ``elementary'' string
states (the simple roots) can only be fewer than all physical states; if their generating function already equals the
full partition function (KNZ), there is nothing else.
\end{physics}
\begin{status}
\textbf{Unconditional}, given its displayed hypothesis \lean{hK}.  Printed in
\leanpath{VOA/Moonshine/MoonshineLRGenerationAudit.lean} (\chtwentylean{BorcherdsPreData.lieSpan_eq_top_of_KNZ}); standard axioms.
\end{status}

\subsection{The supporting statements, in Lean}

\begin{leancode}
theorem coneProd_neg_mul_knzSeries (e f : g) (he : e ∈ D.weightSpace borcherdsRealRoot)
    (hf : f ∈ D.weightSpace (-borcherdsRealRoot)) (hef : D.form e f = 1) :
    coneProd (fun m n => -(finrank ℂ (D.weightSpace ((m : ℤ), (n : ℤ))) : ℂ)) *
      knzSeries (fun j => (finrank ℂ (D.weightSpace (1, (j : ℤ))) : ℂ)) = 1 := by

theorem finrank_quotientWeight_of_KNZ (hK : P.KNZ) (m n : ℕ) (hm : 0 < m) (hn : 0 < n) :
    finrank ℂ (P.quotientWeight ((m : ℤ), (n : ℤ))) =
      finrank ℂ (P.weightSpace ((m : ℤ), (n : ℤ))) := by

theorem weightSpace_le_simpleSpan_of_KNZ (hK : P.KNZ) (m n : ℕ) (hm : 0 < m) (hn : 0 < n) :
    P.weightSpace ((m : ℤ), (n : ℤ)) ≤ P.simpleSpan.toSubmodule := by

def opposite : BorcherdsPreData g where
  weightSpace γ := P.weightSpace (-γ)
\end{leancode}
(\leanpath{VOA/Moonshine/MoonshineLRGenerationSandwich.lean}; \lean{D : BorcherdsData g}, \lean{P : BorcherdsPreData g}.)
\begin{leancode}
def simpleRadical : LieIdeal ℂ P.simpleSpan where
  carrier := {x | ∀ y, P.simpleForm x y = 0}

abbrev simpleQuotient : Type _ := P.simpleSpan ⧸ P.simpleRadical

def quotientWeight (γ : ℤ × ℤ) : Submodule ℂ P.simpleQuotient :=
  (P.simpleSpanWeight γ).map (P.quotientMap : P.simpleSpan →ₗ[ℂ] P.simpleQuotient)
\end{leancode}
(\leanpath{VOA/Moonshine/MoonshineLRGenerationRadical.lean}.)
\begin{leancode}
theorem finrank_quotientWeight {γ : ℤ × ℤ} (hγ : P.weightSpace γ ≤ P.simpleSpan.toSubmodule)
    (hγ' : P.weightSpace (-γ) ≤ P.simpleSpan.toSubmodule) :
    finrank ℂ (P.quotientWeight γ) = finrank ℂ (P.weightSpace γ) := by

theorem quotient_generated :
    LieSubalgebra.lieSpan ℂ P.simpleQuotient (borcherdsSimpleGenerators P.quotientWeight) = ⊤ := by

def quotientData : BorcherdsData P.simpleQuotient where
  weightSpace := P.quotientWeight
\end{leancode}
(\leanpath{VOA/Moonshine/MoonshineLRGenerationQuotientData.lean}.)
\begin{leancode}
theorem positivePart_coneProd_mul_one_sub (e f : g) (he : e ∈ D.weightSpace borcherdsRealRoot)
    (hf : f ∈ D.weightSpace (-borcherdsRealRoot)) (hef : D.form e f = 1) :
    coneProd (fun m n => -(finrank ℂ (positiveWeight D (m, n)) : ℂ)) *
      (1 - finrankSeries (positiveGenWeight D e)) = 1 := by
\end{leancode}
(\leanpath{VOA/Moonshine/MoonshineLRGenerationPositive.lean}.)
\begin{leancode}
theorem eq_of_coneProd_neg_mul_eq_one {c d : ℕ → ℕ → ℂ} {X : PowerSeries (LaurentSeries ℂ)}
    (h₁ : coneProd (fun m n => -d m n) * (1 - X) = 1) (h₂ : coneProd c = 1 - X) :
    ∀ m n : ℕ, 0 < m → 0 < n → c m n = d m n := by
\end{leancode}
(\leanpath{VOA/Moonshine/MoonshineLRGenerationCone.lean}.)
\begin{dictionary}
\lean{e}, \lean{f} with \lean{hef} are a normalised pair of real root vectors ($e\in\mathfrak g_{(1,-1)}$,
$f\in\mathfrak g_{(-1,1)}$, $(e,f)=1$); such a pair exists for any Borcherds data (\chtwentylean{exists_normalized_real_pair},
Chapter~\ref{ch:18}), and it enters only to define the generating space $U$ (\lean{positiveGeneratorSpace}).
\lean{positiveWeight D (m, n)} is $\mathfrak u^+_{(m,n)}$ (the root space viewed inside the positive part), and
\lean{finrankSeries (positiveGenWeight D e)} is $\mathrm{ch}\,U=\sum\dim U_{(a,b)}p^aq^b$.  \lean{P.quotientWeight γ} is
$\mathfrak g'_\gamma$; \lean{P.simpleSpanWeight γ} is $G\cap\mathfrak g_\gamma$ as a subspace of $G$; \lean{simpleRadical} is $R$
(an ideal: the omitted fields prove closure and \lean{lie_mem} from invariance); \lean{⧸} is the quotient Lie algebra.  In
\chtwentylean{eq_of_coneProd_neg_mul_eq_one}, \lean{h₁} and \lean{h₂} are the two identities of step (4) with
$X=\mathrm{ch}\,U$; the conclusion is equality of exponents on the cone.  \lean{opposite} replaces $\mathfrak g_\gamma$ by
$\mathfrak g_{-\gamma}$ (the other fields are omitted).  Each statement is the corresponding step of the proof above.
\end{dictionary}
\begin{status}
All unconditional (given their displayed hypotheses), printed in \leanpath{VOA/Moonshine/MoonshineLRGenerationAudit.lean}
(185 reports, standard axioms, 0 \lean{sorryAx}; build record of the release).
\end{status}

\paragraph{The generic layers underneath.}  Steps (3) and (4) rest on five generic files, which are pure algebra and
combinatorics (no Borcherds data):
\begin{itemize}
\item \leanpath{MoonshineLRGenerationCount.lean}: for a family of basis vectors with degrees in the open cone $\N_{\ge1}^2$,
  finitely many in each degree, the number $N(\Gamma)$ of exponent vectors of total degree $\Gamma$ satisfies the
  logarithmic-derivative recursion $\Gamma_1N(\Gamma)=\sum_i\sum_{k\ge1}(\deg i)_1N(\Gamma-k\deg i)$ (\lean{count_recursion}; a
  double count);
\item \leanpath{MoonshineLRGenerationSeries.lean}: hence $\sum_\Gamma N(\Gamma)p^aq^b=\prod(1-p^mq^n)^{-d(m,n)}$, $d$ the number of
  basis vectors per degree (\lean{expSeries_eq_coneProd}; both sides solve the same first-order equation in $p$);
\item \leanpath{MoonshineLRGenerationGraded.lean}: graded PBW: for a Lie algebra graded by a monoid with a homogeneous basis,
  ordered monomials form a basis of $U(L)$, $U(L)$ is graded (via the coaction $U(L)\to U(L)[G]$), and in the cone
  $\dim U(L)_\Gamma=N(\Gamma)$ (\chtwentylean{finrank_envGraded_eq_expCount});
\item \leanpath{MoonshineLRGenerationTensor.lean}: if $T(S)\to U(L)$ is bijective for a graded subspace $S$, the words in a
  homogeneous basis of $S$ are a basis of $U(L)$, and
  $\prod(1-p^mq^n)^{-\dim L_{(m,n)}}\cdot(1-\mathrm{ch}\,S)=1$ (\chtwentylean{coneProd_mul_one_sub_basisDegSeries});
\item \leanpath{MoonshineLRGenerationCone.lean}: the same in dimension form for cone-graded spaces with finite-dimensional
  pieces (\chtwentylean{coneProd_finrank_mul_one_sub_finrankSeries}), homogeneous bases of such spaces (\lean{coneBasis}), and
  the comparison of exponents (\chtwentylean{eq_of_coneProd_neg_mul_eq_one});
\end{itemize}
and on three files about quotients: \leanpath{MoonshineLRGenerationQuotient.lean} (\lean{form_decompose_eq}: a form pairing
$\gamma$ only with $-\gamma$ satisfies $B(x_\gamma,y)=B(x,y_{-\gamma})$; \lean{isInternal_map_mkQ}: a quotient by a homogeneous submodule is
graded), \leanpath{MoonshineLRGenerationRadical.lean} (the radical, the quotient, the descended form) and
\leanpath{MoonshineLRGenerationQuotientData.lean} (the fields of \lean{quotientData}: Cartan injective with range
$\mathfrak g'_0$, action, symmetric invariant nondegenerate form, support, one-dimensional real root spaces, generation).  We
quote the three generic headlines:
\begin{leancode}
theorem count_recursion (Γ : ℕ × ℕ) :
    Γ.1 * expCount deg low Γ =
      ∑ i ∈ low Γ, ∑ k ∈ Finset.Icc 1 Γ.1,
        if k • deg i ≤ Γ then (deg i).1 * expCount deg low (Γ - k • deg i) else 0 := by
\end{leancode}
(\leanpath{VOA/Moonshine/MoonshineLRGenerationCount.lean}.)
\begin{leancode}
theorem expSeries_eq_coneProd :
    expSeries deg low = coneProd (fun m n => -(basisDegCount deg low (m, n) : ℂ)) := by
\end{leancode}
(\leanpath{VOA/Moonshine/MoonshineLRGenerationSeries.lean}.)
\begin{leancode}
theorem finrank_envGraded_eq_expCount (Γ : ℕ × ℕ) :
    finrank ℂ (envGraded wt hbr Γ) = expCount deg low Γ := by
\end{leancode}
(\leanpath{VOA/Moonshine/MoonshineLRGenerationGraded.lean}.)
\begin{dictionary}
In all three, \lean{deg : ι → ℕ × ℕ} (or \lean{A → ℕ × ℕ}) are the degrees of the basis vectors and \lean{low Γ} is a
finite set containing exactly the indices of degree $\le\Gamma$ (hypothesis \lean{hlow}); the section hypotheses
\lean{hdeg : ∀ i, 1 ≤ (deg i).1 ∧ 1 ≤ (deg i).2} put all degrees in the open cone.  \lean{expCount deg low Γ} is $N(\Gamma)$,
\lean{basisDegCount deg low γ} is $d(\gamma)$, \lean{expSeries} the generating series of $N$, \lean{envGraded wt hbr Γ} the
degree-$\Gamma$ part of $U(L)$ for the grading \lean{wt} of $L$ (with \lean{hbr}: brackets add degrees, and \lean{hb}: the
basis is homogeneous).  These are the classical PBW character formulas, stated carefully enough to be used with complex
exponents.
\end{dictionary}
\begin{physics}
The Fock-space counting of a free (bosonic) second-quantised theory: the partition function of free particles with $d(\gamma)$
species at each momentum $\gamma$ is $\prod(1-e^\gamma)^{-d(\gamma)}$.
\end{physics}
\begin{status}
Generic, unconditional, printed in the Generation audit; standard axioms.
\end{status}

\subsection{The Monster instance}\label{ch20:sec:genmonster}

\begin{theorem}[The Monster Lie algebra is generated by its simple root spaces]\label{ch20:thm:generation}
$\chtwentym$ is generated by $\chtwentym_0$ and the root spaces $\chtwentym_{\pm(1,-1)}$, $\chtwentym_{\pm(1,j)}$, $j\ge1$.
\end{theorem}
\begin{leancode}
def monsterLiePreData : BorcherdsPreData monsterLieAlgebra where
  weightSpace := monsterLieBorcherdsWeight
  weight_internal := monsterLieBorcherdsWeight_isInternal
  weight_finite := monsterLieBorcherdsWeight_finite

theorem finrank_monsterLieBorcherdsWeight_cone (m n : ℕ) (hm : 0 < m) (_hn : 0 < n) :
    (finrank ℂ (monsterLieBorcherdsWeight ((m : ℤ), (n : ℤ))) : ℂ) = cJ (m * n) := by

theorem monsterLiePreData_KNZ : monsterLiePreData.KNZ := by

theorem monsterLieSimpleRootGeneration : MonsterLieSimpleRootGeneration :=
  monsterLiePreData.lieSpan_eq_top_of_KNZ monsterLiePreData_KNZ
\end{leancode}
(\leanpath{VOA/Moonshine/MoonshineLRGenerationMonster.lean}, namespace \chtwentylean{VOA.Moonshine.ConwayNorton.M8}; the remaining
fields of \lean{monsterLiePreData} are the corresponding fields of \lean{monsterLieBorcherdsData}, copied one by one.)
\begin{dictionary}
\lean{monsterLiePreData} uses exactly the grading, Cartan and form of Chapter~\ref{ch:18}'s \lean{monsterLieBorcherdsData}, so
\chtwentylean{monsterLiePreData.weightSpace} is \lean{monsterLieBorcherdsWeight} \emph{by definition} and the conclusion of
\lean{lieSpan_eq_top_of_KNZ} is literally \chtwentylean{MonsterLieSimpleRootGeneration} (Lean accepts the term).  The dimension
statement says $\dim\chtwentym_{(m,n)}=c(mn)$ for $m,n\ge1$: by the no-ghost isomorphism
$\chtwentym_{(m,n)}\cong\Vnat_{1+mn}$ and $\dim\Vnat_{1+mn}=[q^{1+mn}]JShift=c(mn)$ (\chtwentylean{finrank_gradingNat_eq_coeff_JShift}).
(The underscore in \lean{_hn} marks a hypothesis not used in the proof; it is harmless.)  KNZ for $\chtwentym$ is then
\lean{knz_product} after rewriting both sides with these dimensions.  \textbf{The headline has no hypotheses.}
\textbf{Could it be vacuous?}  $\chtwentym$ is the specific Lie algebra of Chapter~\ref{ch:18}, with root spaces of the
dimensions just stated, so ``generated by'' is a substantial statement; and the sign convention (Borcherds' $(m,n)$ is the
library momentum $(m,-n)$) is built into \lean{monsterLieBorcherdsWeight}, so the cone $m,n\ge1$ is the right one (the root
spaces at library momenta $x,y\ge1$ are $\Vnat_{1-xy}$, which is $0$ except at $(x,y)=(1,1)$, the real root $(1,-1)$ in
Borcherds' coordinates).
\end{dictionary}
\paragraph{Proof on paper.}  Theorem~\ref{ch20:thm:genKNZ} for the pre-data of $\chtwentym$; KNZ holds because
$\dim\chtwentym_{(m,n)}=c(mn)$ on the cone and $\dim\chtwentym_{(1,j)}=c(j)$, so both sides of KNZ are those of
\eqref{ch20:eq:knz}, which is \lean{knz_product}. \qed
\begin{physics}
The physical states of the moonshine string are generated, as a Lie algebra under the string vertex (the bracket of
BRST classes), by the tachyon state and the states of unit $m$-momentum.  This is the input that turns $\chtwentym$ into a
generalised Kac--Moody algebra with a Weyl--Kac--Borcherds denominator formula, i.e.\ into a tool for computing twisted
partition functions.
\end{physics}
\begin{status}
\textbf{Unconditional}.  Printed in \leanpath{VOA/Moonshine/MoonshineLRGenerationAudit.lean} and in
\leanpath{VOA/Moonshine/MoonshineCNFinalAudit.lean}: \texttt{\#print axioms} $=$ \lean{[propext, Classical.choice, Quot.sound]}
(build record of the release).
\end{status}

\begin{warning}
\textbf{Do not look for a Lie-theoretic proof of generation.}  Given the other Borcherds
data, generation of $\chtwentym$ by the known simple roots is \emph{equivalent} to KNZ: the sandwich gives
``KNZ $\Rightarrow$ generation'' (this section), and conversely generation makes the character identity (3) a statement about
$\chtwentym$ itself, which is KNZ.  So any proof that $\chtwentym$ has no further simple roots off the line $m=1$ is a proof of
the product formula for $J(p)-J(q)$, a statement about the modularity of $J$.  The Casimir argument of Chapter~\ref{ch:18}
cannot exclude extra simple roots $(m,n)$, $m\ge2$; it only excludes relations among the known ones.  The equivalence is
a paper statement; the library proves only the direction it needs.
\end{warning}
\begin{warning}
\textbf{Stale docstring.}  The docstring of \chtwentylean{MonsterLieSimpleRootGeneration} in
\leanpath{VOA/Moonshine/MoonshineCNMonsterBorcherdsData.lean} still says that generation ``is the only Borcherds-data field not
yet proved''; it is now proved (\chtwentylean{monsterLieSimpleRootGeneration}).  Its second clause (``it is
equivalent, given the other fields, to the product formula'') is the paper equivalence above; only one direction is in Lean.
The Lean statements are unaffected.
\end{warning}

\subsection{How generation feeds the twisted denominator identities}

With generation, Chapter~\ref{ch:18}'s \lean{monsterLieBorcherdsData monsterLieSimpleRootGeneration} is Borcherds data, and
Chapter~\ref{ch:19} applies to it.  For every automorphism $\sigma$ of $\Vnat$ of finite order, $\sigma$ acts on $\chtwentym$
by graded Lie algebra automorphisms preserving the form and fixing the Cartan and the real root spaces; taking
$\sigma$-equivariant graded traces in the homology of the free Lie algebra $\mathfrak u^+$ gives the twisted denominator
identity, whose coefficients are the traces $\tr(\sigma^k|\Vnat_n)$.  In Lean:
\begin{leancode}
theorem twistedDenominatorIdentity_of_simpleRootGeneration (hgen : MonsterLieSimpleRootGeneration)
    (σ : AutVnat) (hσ : IsOfFinOrder σ) : TwistedDenominatorIdentity (mckayThompsonFamily σ) :=

theorem completelyReplicable_of_simpleRootGeneration (hgen : MonsterLieSimpleRootGeneration)
    (hfin : ∀ g : AutVnat, IsOfFinOrder g) : CompletelyReplicable :=

theorem conwayNorton_of_simpleRootGeneration (hgen : MonsterLieSimpleRootGeneration)
    (hfin : ∀ g : AutVnat, IsOfFinOrder g) : ConwayNortonConjecture :=
  M8.conwayNorton_of_finOrder_of_completelyReplicable hfin
    (completelyReplicable_of_simpleRootGeneration hgen hfin)
\end{leancode}
(\leanpath{VOA/Moonshine/MoonshineCNTwistedDenominator.lean}, Chapter~\ref{ch:19}.)
\begin{dictionary}
\lean{TwistedDenominatorIdentity c} (Chapter~\ref{ch:19}) is a formal power-series identity in the coefficient array
\lean{mckayThompsonFamily σ} $=(k,n)\mapsto\tr(\sigma^k|\Vnat_{n+1})$; \lean{CompletelyReplicable} is the proposition of
Chapter~\ref{ch:16} (Definition~\ref{ch16:def:replicable}).  The hypotheses \lean{hgen} and \lean{hfin} are \emph{exactly} the two
inputs (I), (II) of the introduction; this chapter supplies both.  Generation is used once, to build the Borcherds data;
finite order is used for the twisted identities (\lean{hσ}) and for the analytic replication bounds.
\end{dictionary}
\begin{physics}
The twisted denominator identity is the statement that the $\sigma$-twisted second-quantised string partition function
equals its free-field (simple-root) expression; generation is what makes the simple-root expression complete.
\end{physics}
\begin{status}
Conditional on \lean{hgen}, \lean{hfin} as displayed; both discharged in this chapter.  Printed in
\leanpath{VOA/Moonshine/MoonshineCNAudit18.lean}; standard axioms (Chapter~\ref{ch:16}).
\end{status}

%======================================================================
\section[The Conway--Norton theorem for every automorphism]{The Conway--Norton theorem for every automorphism of the moonshine module}\sectionmark{The Conway--Norton theorem}\label{ch20:sec:assembly}
%======================================================================

\subsection{From complete replicability and finite order}

Chapter~\ref{ch:19} proves the genus-zero theorem in the form ``finite order, complete replicability, the prime Hecke
identity for the McKay--Thompson families and their non-injectivity modulo translations imply the conjecture'', and
separately proves the last two inputs for \emph{every} replicable family.  The first assembly step, in a file of this
chapter, combines them.

\begin{theorem}[Conway--Norton from finite order and complete replicability]\label{ch20:thm:cnrep}
If every automorphism of $\Vnat$ has finite order and the McKay--Thompson series are completely replicable, then the
Conway--Norton conjecture holds for every automorphism of $\Vnat$.
\end{theorem}
\begin{leancode}
theorem conwayNorton_of_finOrder_of_completelyReplicable
    (hfin : ∀ g : AutVnat, IsOfFinOrder g) (hrep : CompletelyReplicable) :
    ConwayNortonConjecture :=
  conwayNorton_of_completelyReplicable_of_families hfin hrep
    (fun _ F k => ReplicableFamily_primeHeckeIdentity F k)
    (fun _ hN F _ => ReplicableFamily_notInjectiveModTranslations hN F)
\end{leancode}
(\leanpath{VOA/Moonshine/MoonshineLRConwayNorton.lean}, namespace \chtwentylean{VOA.Moonshine.ConwayNorton.M8}.)  The three
ingredients from Chapter~\ref{ch:19}:
\begin{leancode}
theorem conwayNorton_of_completelyReplicable_of_families (hfin : ∀ g : AutVnat, IsOfFinOrder g)
    (hrep : CompletelyReplicable)
    (hB1 : ∀ (N : ℕ) (F : ReplicableFamily N) (k : ℕ), PrimeHeckeIdentity F k)
    (hD7 : ∀ (N : ℕ), 0 < N → ∀ F : ReplicableFamily N, (∀ k, PrimeHeckeIdentity F k) →
      F.NotInjectiveModTranslations) :
    ConwayNortonConjecture :=
\end{leancode}
(\leanpath{VOA/Moonshine/MoonshineCNGenusZeroMain.lean}.)
\begin{leancode}
theorem ReplicableFamily_primeHeckeIdentity {N : ℕ} (F : ReplicableFamily N) (k : ℕ) :
    PrimeHeckeIdentity F k := by
\end{leancode}
(\leanpath{VOA/Moonshine/MoonshineLRHecke.lean}.)
\begin{leancode}
theorem ReplicableFamily_notInjectiveModTranslations {N : ℕ} (hN : 0 < N)
    (F : ReplicableFamily N) : F.NotInjectiveModTranslations := by
\end{leancode}
(\leanpath{VOA/Moonshine/MoonshineLRNontrivial.lean}.)
\begin{dictionary}
\lean{CompletelyReplicable} (Chapter~\ref{ch:16}, Definition~\ref{ch16:def:replicable}): for every $g$ and $m\ge1$ there is a
monic polynomial $P$ of degree $m$ with $P(T_g(\tau))=\sum_{ad=m}\sum_{0\le b<d}T_{g^a}\bigl(\tfrac{a\tau+b}d\bigr)$.
\lean{ReplicableFamily N} is the abstract family of $q$-series of period $N$ of Chapter~\ref{ch:16}
(Definition~\ref{ch16:def:family}).  \lean{PrimeHeckeIdentity F k} is a Hecke-orbit identity at primes $p\equiv1\pmod N$, and
\chtwentylean{F.NotInjectiveModTranslations} says that no member of the family is injective on $\mathbb H/\Z$ (Chapter~\ref{ch:19}).
The two lambda terms show that the hypotheses \lean{hB1}, \lean{hD7} of the genus-zero theorem are discharged by
hypothesis-free theorems for all families; so the only hypotheses left are \lean{hfin} and \lean{hrep}, as stated.
\end{dictionary}
\paragraph{Proof on paper.}  The proof term \emph{is} the proof: Chapter~\ref{ch:19}'s genus-zero theorem with its two
family-level hypotheses filled in. \qed
\begin{physics}
``Complete replicability'' is the statement that the twisted partition functions of $g$ and all its powers are related by
Hecke operators, as expected from the second-quantised (symmetric-product) string theory; the genus-zero argument
converts this into ``each twisted partition function is a Hauptmodul''.
\end{physics}
\begin{status}
Conditional on \lean{hfin}, \lean{hrep} (both later discharged).  Printed in
\leanpath{VOA/Moonshine/MoonshineLRConwayNortonAudit.lean} (its single report) and in
\leanpath{VOA/Moonshine/MoonshineCNAudit18.lean}; standard axioms.
\end{status}

\subsection{The final theorem}

\begin{theorem}[Conway--Norton for $\Aut(\Vnat)$; Borcherds]\label{ch20:thm:final}
For every automorphism $g$ of $\Vnat$ there are a subgroup $\Gamma_g\le\mathrm{SL}_2(\mathbb R)$ and $N\ge1$ such that
$\Gamma_0(N)\le\Gamma_g$, $\Gamma_g$ is commensurable with $\mathrm{SL}_2(\Z)$, and $T_g=q^{-1}+\sum_{n\ge1}\tr(g|\Vnat_{n+1})q^n$ is
the normalised Hauptmodul of $\Gamma_g$.  (This is the qualitative genus-zero form: $\Gamma_g$ is not identified with
Conway and Norton's group and $N$ is not related to the order of $g$; it is weaker than the table-level conjecture.)
\end{theorem}
\begin{leancode}
theorem conwayNorton_of_isOfFinOrder (hfin : ∀ g : AutVnat, IsOfFinOrder g) :
    ConwayNortonConjecture :=
  conwayNorton_of_simpleRootGeneration monsterLieSimpleRootGeneration hfin
\end{leancode}
(\leanpath{VOA/Moonshine/MoonshineLRGenerationMonster.lean}.)
\begin{leancode}
theorem conwayNortonConjecture_holds : ConwayNortonConjecture :=
  conwayNorton_of_isOfFinOrder autVnat_isOfFinOrder
\end{leancode}
(\leanpath{VOA/Moonshine/MoonshineLRConwayNortonFinal.lean}, namespace \chtwentylean{VOA.Moonshine.ConwayNorton.M8}.)
\begin{leancode}
theorem conwayNorton : ConwayNortonConjecture :=
  M8.conwayNortonConjecture_holds
\end{leancode}
(\leanpath{VOA/Moonshine/MoonshineCNFinalAudit.lean}, namespace \chtwentylean{VOA.Moonshine.ConwayNorton}; Chapter~\ref{ch:16}.)
\begin{dictionary}
\lean{ConwayNortonConjecture} is the proposition of Chapter~\ref{ch:16} (Definition~\ref{ch16:def:cn}); its clause-by-clause
comparison with Conway and Norton's conjecture is Chapter~\ref{ch:16}, \S\ref{ch16:sec:fidelity}: it is the genus-zero
(qualitative) form, for the full automorphism group of the library's $\Vnat$, weaker than the table-level conjecture in that
$\Gamma_g$ is not identified with Conway--Norton's group, not required to lie in the normaliser of $\Gamma_0(N)$, and $N$ is not
related to the order of $g$.  The three declarations above are one-line proof terms:
\chtwentylean{conwayNorton_of_isOfFinOrder} plugs generation (Theorem~\ref{ch20:thm:generation}) into
\chtwentylean{conwayNorton_of_simpleRootGeneration}; \chtwentylean{conwayNortonConjecture_holds} plugs in finite order
(Theorem~\ref{ch20:thm:finorder}); and \lean{conwayNorton} is \emph{definitionally} the same proof under the
name fixed in the statement file.  So \lean{conwayNorton} and \chtwentylean{M8.conwayNortonConjecture_holds} are the same theorem;
the former is the official name.  \textbf{No hypotheses}, no section variables, no type-class assumptions.
\end{dictionary}

\paragraph{The whole chain, unwound.}  Reading the proof terms from the top:
\begin{enumerate}[label=(\arabic*)]
\item \lean{conwayNorton} $:=$ \chtwentylean{conwayNortonConjecture_holds} $:=$ \chtwentylean{conwayNorton_of_isOfFinOrder} applied to
  \lean{autVnat_isOfFinOrder} (\S\ref{ch20:sec:finorder}; Ising rigidity, periodicity, the ingredients of the faithfulness
  proof with weakened hypotheses, and $\sigma$; not the faithfulness theorem itself);
\item \chtwentylean{conwayNorton_of_isOfFinOrder} $:=$ \chtwentylean{conwayNorton_of_simpleRootGeneration} applied to
  \chtwentylean{monsterLieSimpleRootGeneration} (\S\ref{ch20:sec:generation}; the sandwich and \lean{knz_product});
\item \chtwentylean{conwayNorton_of_simpleRootGeneration} (Chapter~\ref{ch:19}) $:=$
  \chtwentylean{conwayNorton_of_finOrder_of_completelyReplicable} applied to \chtwentylean{completelyReplicable_of_simpleRootGeneration},
  i.e.\ to the twisted denominator identities of the Borcherds data \lean{monsterLieBorcherdsData} (Chapters~\ref{ch:18},
  \ref{ch:19});
\item \chtwentylean{conwayNorton_of_finOrder_of_completelyReplicable} (this chapter) $:=$ the genus-zero theorem
  \chtwentylean{conwayNorton_of_completelyReplicable_of_families} with the family-level Hecke and non-injectivity theorems
  (Chapter~\ref{ch:19});
\item which in turn is the $\Gamma_1$ genus-zero argument followed by the bridge $\Gamma_1\Rightarrow\Gamma_0$
  \lean{conwayNorton_of_gammaOne} (Chapter~\ref{ch:16}, \S\ref{ch16:sec:gammazero}).
\end{enumerate}
The inputs from Chapters~\ref{ch:18}--\ref{ch:19} that this chapter uses are exactly the declarations quoted with their
paths in \S\ref{ch20:sec:generation} and in this section; everything else in the chain is this chapter's.

\begin{physics}
Every global symmetry $g$ of the Monster CFT has a twisted torus partition function $Z\bigl[{}^{1}_{g}\bigr](\tau)=\tr(g\,q^{L_0-1})$
that is a global coordinate on a genus-zero modular curve.  In Borcherds' string-theory proof: compactify the bosonic
string on the moonshine orbifold times a Lorentzian torus; its physical states form $\chtwentym$; the denominator formula of
$\chtwentym$ is the product formula for $J(p)-J(q)$; generation (which is that formula) and finite order of $g$ (which follows
here from the rigidity of the Ising sub-theories) give the twisted denominator formulas, i.e.\ complete replicability; and
complete replicability forces genus zero.
\end{physics}
\begin{status}
\textbf{Unconditional.}  \texttt{\#print axioms} for \lean{conwayNorton}, \chtwentylean{M8.conwayNortonConjecture_holds},
\chtwentylean{M8.monsterLieSimpleRootGeneration}, \lean{M8.autVnat_isOfFinOrder} (the four commands of
\leanpath{VOA/Moonshine/MoonshineCNFinalAudit.lean}): \lean{[propext, Classical.choice, Quot.sound]}, no \lean{sorryAx},
targeted build and the umbrella \lean{VOA.Moonshine} green (build record of the release).
\chtwentylean{conwayNortonConjecture_holds} is also printed in \leanpath{VOA/Moonshine/MoonshineLRFinOrderAudit.lean}, and
\chtwentylean{conwayNorton_of_isOfFinOrder} in \leanpath{VOA/Moonshine/MoonshineLRGenerationAudit.lean}.  The book's authors did not
re-run the builds.  The repository-wide aggregate \texttt{lake build} had not been run for this part at the time of the
build record; the targeted and umbrella builds are the machine evidence.  The theorem
does \emph{not} depend on \lean{GmatEqAutO}: it is about $\Aut(\Vnat)$, whatever that group is.
\end{status}

\subsection{What the Conway--Norton theorem does and does not use from this chapter}

\begin{itemize}
\item \textbf{Used:} \chtwentylean{monsterLieSimpleRootGeneration}, \lean{autVnat_isOfFinOrder}, and through the latter the
  periodicity lemma, the Ising vectors and their rigidity, the faithfulness ingredients (zero modes, joint eigenvectors,
  twisted exclusion, leading products, the Leech span), and \lean{sigmaAut}.
\item \textbf{Not used:} \chtwentylean{vnatGriessRestrict_injective} and \chtwentylean{aut_eq_one_of_fix_weightTwo} themselves (their
  hypotheses are too strong for the finite-order argument; weakened copies are used), \chtwentylean{vnatGriessRestrict_injective_of_gen},
  and the finiteness theorem \lean{finite_autVnat} with its generic lemmas.  These serve the identification of $\Aut(\Vnat)$
  with the Monster (Chapters~\ref{ch:21}--\ref{ch:22}).
\end{itemize}

%======================================================================
\section{The files of this chapter}\label{ch20:sec:files}
%======================================================================

All $33$ files of the chapter (the families \texttt{MoonshineLR\{FinOrder*,} \texttt{Generation*,}
\texttt{ConwayNorton*,} \texttt{FiniteRigid,} \texttt{FiniteAut\}} and \texttt{Aut*}) are under \texttt{VOA/Moonshine/}.  ``Individually'' means every
headline declaration is presented under the Fidelity Protocol; ``collectively'' means the file's content is described with
its role and its main statements quoted.
\begin{center}\small
\begin{longtable}{p{0.3\textwidth}p{0.45\textwidth}p{0.15\textwidth}}
\toprule
file & content & coverage\\
\midrule
\endhead
\chtwentyfile{AutVnat.lean} & \lean{VnatVOA}, \lean{AutVnat}, Griess restriction, faithfulness from generation & individually
  (\S\ref{ch20:sec:aut}; also Ch.~\ref{ch:16})\\
\chtwentyfile{AutSigma.lean} & $\sigma^\natural\in\Aut(\Vnat)$ & individually (\S\ref{ch20:sec:sigma})\\
\chtwentyfile{AutFaithful.lean} & the leading product \eqref{ch20:eq:leading} in $\chtwentyVL$ & individually
  (Lemma~\ref{ch20:lem:leading})\\
\chtwentyfile{AutFaithfulVnat.lean} & $o(q_{ab})x_\beta$, \lean{eigen_plus}, block form, \lean{aut_frameX0P_eigen}, scalar steps &
  individually (Lemmas~\ref{ch20:lem:jointeigen}, \ref{ch20:lem:scalars})\\
\chtwentyfile{AutFaithfulAssembly.lean} & faithfulness, unconditionally & individually (Thm.~\ref{ch20:thm:faithful})\\
\chtwentyfile{AutChargeEigen.lean} & joint eigenvectors in $V_L$, charge sector by charge sector & headline individually
  (Lemma~\ref{ch20:lem:jointeigen}), rest collectively\\
\chtwentyfile{AutFockEigen.lean} & $N_{ab}$ in closed form; \lean{quadOp_joint_eigen} & headlines individually
  (Lemmas~\ref{ch20:lem:zeromode}, \ref{ch20:lem:jointeigen}); coefficient lemmas collectively\\
\chtwentyfile{AutLeechSpan.lean} & $\Leech$ is generated by its minimal vectors & individually (Thm.~\ref{ch20:thm:leechspan})\\
\chtwentyfile{AutZeroModes.lean} & $o(q_{ab})$ on $V_L$ & headline individually (Lemma~\ref{ch20:lem:zeromode}); two-direction
  normal-ordering lemmas collectively\\
\chtwentyfile{AutZeroModesTwisted.lean} & $o(q_{ab})=N^{\mathrm{tw}}_{ab}+\tfrac18\delta_{ab}$ on the twisted sector & headline
  individually (Lemma~\ref{ch20:lem:twisted}); the twisted product-formula computation collectively\\
\chtwentyfile{AutTwistedExclusion.lean} & no twisted joint eigenvector of eigenvalue $\beta\beta^{\mathsf T}$ & headline
  individually (Lemma~\ref{ch20:lem:twisted}); tensor-component lemmas collectively\\
\chtwentyfile{MoonshineLRFinOrderPeriodic.lean} & periodicity of rigid idempotents & individually (\S\ref{ch20:sec:periodic})\\
\chtwentyfile{MoonshineLRFinOrderIsing.lean} & Ising vectors, Griess products, spanning sets & individually
  (\S\ref{ch20:sec:ising}); coordinate lemmas (\lean{isingQuad}, \ldots) collectively\\
\chtwentyfile{MoonshineLRFinOrderRigid.lean} & spectral lemma, rigidity, finite order & individually
  (\S\S\ref{ch20:sec:rigid}, \ref{ch20:sec:finorder}); private lemmas collectively\\
\chtwentyfile{MoonshineLRFiniteRigid.lean} & finitely many rigid idempotents; derivations & individually
  (Thm.~\ref{ch20:thm:finrigid}); private commutative-algebra lemmas collectively\\
\chtwentyfile{MoonshineLRFiniteAut.lean} & $\Aut(\Vnat)$ finite & individually (\S\ref{ch20:sec:finite})\\
\chtwentyfile{MoonshineLRGenerationCount.lean} & counting exponent vectors & collectively, headline quoted\\
\chtwentyfile{MoonshineLRGenerationSeries.lean} & count series $=$ cone product & collectively, headline quoted\\
\chtwentyfile{MoonshineLRGenerationGraded.lean} & graded PBW for $U(L)$ & collectively, headline quoted\\
\chtwentyfile{MoonshineLRGenerationTensor.lean} & character identity from $T(S)\cong U(L)$ & collectively\\
\chtwentyfile{MoonshineLRGenerationCone.lean} & cone-graded bases; exponent comparison & collectively, one statement quoted\\
\chtwentyfile{MoonshineLRGenerationPositive.lean} & character identity for $\mathfrak u^+$ of Borcherds data & collectively, headline
  quoted\\
\chtwentyfile{MoonshineLRGenerationPreData.lean} & \lean{BorcherdsPreData}, \lean{simpleSpan} & individually
  (Def.~\ref{ch20:def:predata})\\
\chtwentyfile{MoonshineLRGenerationQuotient.lean} & graded quotients; opposite-degree forms & collectively\\
\chtwentyfile{MoonshineLRGenerationRadical.lean} & radical $R$, $G/R$, descended form & collectively, definitions quoted\\
\chtwentyfile{MoonshineLRGenerationQuotientData.lean} & $G/R$ is Borcherds data & collectively, headlines quoted\\
\chtwentyfile{MoonshineLRGenerationSandwich.lean} & KNZ, the sandwich, opposite data, generation from KNZ & individually
  (Thm.~\ref{ch20:thm:genKNZ})\\
\chtwentyfile{MoonshineLRGenerationMonster.lean} & the Monster instance; \chtwentylean{conwayNorton_of_isOfFinOrder} & individually
  (Thms.~\ref{ch20:thm:generation}, \ref{ch20:thm:final})\\
\chtwentyfile{MoonshineLRConwayNorton.lean} & Conway--Norton from finite order and replicability & individually
  (Thm.~\ref{ch20:thm:cnrep})\\
\chtwentyfile{MoonshineLRConwayNortonFinal.lean} & \chtwentylean{conwayNortonConjecture_holds} & individually (Thm.~\ref{ch20:thm:final})\\
\chtwentyfile{MoonshineLRFinOrderAudit.lean} & audit & list below\\
\chtwentyfile{MoonshineLRGenerationAudit.lean} & audit & list below\\
\chtwentyfile{MoonshineLRConwayNortonAudit.lean} & audit & list below\\
\bottomrule
\end{longtable}
\end{center}

\paragraph{The technical lemmas, collectively.}  Besides the statements presented above, the \texttt{Aut*} files contain the
computational lemmas behind them: the two-direction normal-ordered products of lattice currents on charged Fock states
(\leanpath{AutZeroModes.lean}: \lean{nprod_lift_ZfGen₂}, \lean{nprod_alphaLatGen_tmul₂}, \lean{YfullGen_quadState}); coefficient
extraction for $X\partial$ operators on polynomials (\leanpath{AutFockEigen.lean}: \lean{coeff_pderiv}, \lean{lvlWt},
\lean{coeff_quadOp_self}, \lean{coeff_quadOp_ne}); charge coordinates and the rank-one argument $\gamma\gamma^{\mathsf T}=\beta\beta^{\mathsf T}
\Rightarrow\gamma=\pm\beta$ (\leanpath{AutChargeEigen.lean}: \lean{eq_of_chargeCoord_eq}, \lean{chargeCoord_rankOne},
\lean{coordEα_of_eigen}); the twisted closed form $N^{\mathrm{tw}}_{ab}=\sum_k(k+\tfrac12)(X_{k,a}\partial_{k,b}+X_{k,b}\partial_{k,a})$ and its
transport to the Leech twisted module, $\chtwentyTp$ and the twisted block of $\Vnat$ (\leanpath{AutZeroModesTwisted.lean}:
\lean{quadNT_apply}, \lean{mode_Ywt_quad_zero}, \lean{YML_charge_zero}, \lean{YML_quadSt_two}, \lean{plusModule_qP_one},
\lean{VnatVOA_smul_quad_one}); componentwise action of $A\otimes1$ and the rank-two impossibility
$\beta\beta^{\mathsf T}\ne\tfrac18\mathbf 1$ (\leanpath{AutTwistedExclusion.lean}: \lean{component_rTensor},
\lean{rTensor_joint_eigen}, \lean{not_outer_eq_delta}, \lean{quadNT_joint_eigen}); and the explicit $\Z^{24}$ computations of
\leanpath{AutLeechSpan.lean} (\lean{pair_mem_minT}, \lean{four_smul_mem_minC}, \lean{two_smul_R_mem}, \lean{mem_minC_of_mem_carrier}).
Each is part of one of the paper steps of Theorem~\ref{ch20:thm:faithful}; their physical content is the free-boson mode
algebra (untwisted and $\Z_2$-twisted).  The \texttt{MoonshineLRFinOrder*} files contain, besides the quoted statements, the
coordinate form $Q(u,v)=\sum u_av_bq_{ab}$ of the quadratic letters (\lean{isingQuadRow}, \lean{isingQuad},
\lean{isingQuad_one_isingQuad}), the spectral-decomposition operators of the periodicity lemma (\lean{spectralDecomposition},
\lean{spectralWeight}, \lean{semisimplePart}, \lean{characterDerivation}), and private lemmas carrying out the case analysis
of Theorem~\ref{ch20:thm:rigid} and the scalar propagation with weakened hypotheses.  \leanpath{MoonshineLRFiniteRigid.lean}
contains private lemmas on point ideals ($\ker(\mathrm{eval}_p)=(X_i-p_i)$), the Nakayama step for isolated points, and the
coordinate form of a bilinear product.

\paragraph{The audit modules.}
\begin{itemize}
\item \leanpath{VOA/Moonshine/MoonshineLRFinOrderAudit.lean}: imports \chtwentyfile{MoonshineLRFinOrderIsing}, \chtwentyfile{MoonshineLRFinOrderPeriodic},
  \chtwentyfile{MoonshineLRFinOrderRigid}, \chtwentyfile{MoonshineLRConwayNortonFinal}, \chtwentyfile{MoonshineLRFiniteAut}; $86$ \texttt{\#print axioms} reports at the release,
  including \chtwentylean{pow_apply_eq_self_of_rigid}, \lean{isingVec_rigid}, \lean{autVnat_isOfFinOrder},
  \chtwentylean{conwayNortonConjecture_holds}, \chtwentylean{finite_setOf_rigid_idempotent}, \chtwentylean{derivation_eq_zero_of_span_rigid},
  \lean{finite_autVnat}.  Per the build record all reports are \lean{[propext, Classical.choice, Quot.sound]}, 0
  \lean{sorryAx}.  (An earlier build message reported $118$ reports for this audit, another $86$; the file at the
  release contains $86$ commands.  The record does not explain the difference.)
\item \leanpath{VOA/Moonshine/MoonshineLRGenerationAudit.lean}: imports all \texttt{MoonshineLRGeneration*} files; $185$
  reports, including \chtwentylean{BorcherdsPreData.lieSpan_eq_top_of_KNZ}, \chtwentylean{monsterLieSimpleRootGeneration},
  \chtwentylean{conwayNorton_of_isOfFinOrder}; standard axioms, 0 \lean{sorryAx}.
\item \leanpath{VOA/Moonshine/MoonshineLRConwayNortonAudit.lean}: one report,
  \chtwentylean{conwayNorton_of_finOrder_of_completelyReplicable}; standard axioms.
\end{itemize}
The \texttt{Aut*} files have no audit module of their own: every theorem in them is printed in
\leanpath{VOA/Moonshine/MoonshineLRLiftAudit.lean} (Chapter~\ref{ch:22}; green, $997$
reports in that audit, all within \lean{[propext, Classical.choice, Quot.sound]}, no \lean{sorryAx}).  As in Chapter~\ref{ch:16}: a clean audit line certifies the absence of extra axioms and of \lean{sorry}; it does
not certify that a declaration is unconditional, nor that it means what it should.

%======================================================================
\section{Recorded traps and dead ends}\label{ch20:sec:traps}
%======================================================================

The traps specific to a section are recorded where they arise: the false induction criterion
(\S\ref{ch20:sec:techfaith}), the cocycle sign of the naive Ising vector (\S\ref{ch20:sec:ising}), $c=1$ vectors and diagonal
zero modes (\S\ref{ch20:sec:finorder}), the rejected count of all Ising vectors (\S\ref{ch20:sec:finite}), and the equivalence of
generation with KNZ (\S\ref{ch20:sec:generation}).  We add the cross-cutting ones.

\begin{warning}
\textbf{Sign of the momenta} (Chapter~\ref{ch:16}, \S\ref{ch16:sec:traps}).  Borcherds' root $(m,n)$ is the library momentum
$(m,-n)$, and the root space at library momentum $(x,y)$ is $\Vnat_{1-xy}$.  A product over library momenta $x,y\ge1$ would be
wrong (those root spaces vanish, except the real root space at library $(1,1)$).  The generation files work exclusively in Borcherds coordinates through
\lean{monsterLieBorcherdsWeight}, which applies the sign once.
\end{warning}

\begin{warning}
\textbf{Finite order is not finiteness.}  An infinite group can have every element of finite
order; Theorem~\ref{ch20:thm:finorder} does not imply Theorem~\ref{ch20:thm:finite}, which needed the separate algebraic-geometry
argument.  The periodicity lemma alone gives periodicity of each element, not finiteness of the group.  The converse
implication does hold (a finite group is torsion), as explained in the introduction.  Neither theorem is used to identify
$\Aut(\Vnat)$ with the Monster.
\end{warning}

\begin{warning}
\textbf{Authors' working notes that are out of date.}  (i) The dependency notes,
node M1, still lists ``FINITE ORDER \lean{autVnat_isOfFinOrder}'' with the open marker and the planned route ``their span +
$\sigma$-images $=\Vnat_2$; \chtwentylean{vnatGriessRestrict_injective}'', although the theorem is proved (the same file says so at its
top) and the proof uses neither a spanning statement nor \chtwentylean{vnatGriessRestrict_injective}.  (ii)
the notes on the Monster, \S0 and \S5, say \lean{Finite AutVnat} ``is not delivered''; it is now
\lean{finite_autVnat}.  (iii) The docstring of \chtwentylean{MonsterLieSimpleRootGeneration} calls generation ``not yet proved''
(\S\ref{ch20:sec:generation}).  (iv) The docstring of \lean{AutVnat} calls it the group ``the campaign identifies with the
Monster'' (quoted); the identification is conditional on \lean{GmatEqAutO} (Chapter~\ref{ch:22}).  None affects a Lean statement.
\end{warning}

\begin{warning}
\textbf{The paper route was unrefereed.}  The written plan of the finite-order proof
in the authors' notes was not refereed and was not machine-checked; its normalisations (cocycle sign, $2^{-4}$ versus $2^{-3}$) were the main recorded risk.  The Lean proofs settle them:
the twisted coefficient $\tfrac s4\cdot2\cdot2^{-4}=\tfrac s{32}$ is the one in \lean{isingTwistedMode_hP}.  A referee should
trust the Lean statements as presented here, not the unrefereed plan.
\end{warning}

%======================================================================
\section{Referee's checklist}\label{ch20:sec:checklist}
%======================================================================

Namespaces: \lean{VOA.Moonshine.AutMonster} (A), \chtwentylean{VOA.Moonshine.ConwayNorton.M8} (M8), others as stated.  ``std.''
means \texttt{\#print axioms} $=$ \lean{[propext, Classical.choice, Quot.sound]} as recorded in the build records cited in
the status boxes.
\begin{center}\small
\begin{longtable}{p{0.27\textwidth}p{0.29\textwidth}p{0.2\textwidth}p{0.12\textwidth}}
\toprule
declaration & intended meaning & hypotheses; status & check\\
\midrule
\endhead
\chtwentylean{A.VnatVOA}, \chtwentylean{A.AutVnat} & $\Vnat$; its automorphisms (vacuum, $\omega$, all modes) & definitions &
  Def.~\ref{ch20:def:vnat}\\
\chtwentylean{GriessAut.griessAut}, \chtwentylean{A.vnatGriessRestrict} & automorphisms of $(\Vnat_2,\times,\langle,\rangle_3)$;
  restriction & definitions & Def.~\ref{ch20:def:griessaut}\\
\chtwentylean{A.vnatGriessRestrict_injective_of_gen} & faithfulness if $\Vnat_2$ generates & \textbf{conditional} on
  \chtwentylean{hgen} (unproved) & Thm.~\ref{ch20:thm:faithgen}\\
\chtwentylean{A.aut_eq_one_of_fix_weightTwo}, \chtwentylean{A.vnatGriessRestrict_injective} & $g|_{\Vnat_2}=1\Rightarrow g=1$;
  $\Aut(\Vnat)\hookrightarrow\mathrm{GriessAut}$ & none; std.\ (LiftAudit) & Thm.~\ref{ch20:thm:faithful}\\
\chtwentylean{LatGen.smul_quadState_one}, \chtwentylean{HeisenbergN.nprod_currentN_neg_two} & $o(q_{ab})=\gamma_a\gamma_b+N_{ab}$ & none;
  uncond. & Lem.~\ref{ch20:lem:zeromode}\\
\chtwentylean{HeisenbergN.quadOp_joint_eigen}, \chtwentylean{A.eigen_quadState_span}, \chtwentylean{A.eigen_plus} & joint
  eigenvectors are momentum states & displayed (rank $\ge2$, $\mu$); uncond. & Lem.~\ref{ch20:lem:jointeigen}\\
\chtwentylean{AutZeroTw.VnatVOA_smul_quad_one_snd}, \chtwentylean{AutZeroTw.VnatVOA_twisted_outer_eigen_eq_zero} & twisted zero
  modes; twisted exclusion & none; uncond. & Lem.~\ref{ch20:lem:twisted}\\
\chtwentylean{A.frameX0P_smul_frameX0P_of_neg}, \chtwentylean{A.VnatVOA_smul_xN_of_neg} & leading product & $\langle\alpha,\gamma\rangle<0$;
  uncond. & Lem.~\ref{ch20:lem:leading}\\
\chtwentylean{LeechSpan.leech_closure_minimal_eq_top}, \chtwentylean{LeechSpan.leech_induction} & $\Leech=\Z\Leech_4$ & none; uncond. &
  Thm.~\ref{ch20:thm:leechspan}\\
\chtwentylean{A.aut_frameX0P_eigen}, \chtwentylean{A.scalar_step} & scalars and their propagation & $g|_{\Vnat_2}=1$; uncond. &
  Lem.~\ref{ch20:lem:scalars}\\
\chtwentylean{A.sigmaAut} & $\sigma^\natural\in\Aut(\Vnat)$ & none; definition; used by std.\ theorems & Thm.~\ref{ch20:thm:sigma}\\
\chtwentylean{M8.pow_apply_eq_self_of_rigid} (and steps) & rigid idempotents are periodic & displayed (fin.\ dim., commutative,
  multiplicative $A$, rigid $e$); std. & Thm.~\ref{ch20:thm:periodic}\\
\chtwentylean{M8.isingVec}, \chtwentylean{M8.isingVec_one_self} & $e_{\alpha,s}$; $e\times e=2e$ & $\alpha$ minimal, $s^2\varepsilon(\alpha,-\alpha)=1$;
  std. & Def.~\ref{ch20:def:ising}\\
\chtwentylean{M8.leech_form_cases_of_normFour}, products, \chtwentylean{M8.quadState_one_quadState} & weight-two products &
  displayed; std. & Lem.~\ref{ch20:lem:products}\\
\chtwentylean{M8.leechVOAPlus_grading_two_eq_span}, \chtwentylean{M8.twistedPlus_grading_two_eq_span} & spanning sets of the
  weight-two blocks & none; std. & Lem.~\ref{ch20:lem:span}\\
\chtwentylean{M8.isingVec_zeroMode_plus_annihilated}, \chtwentylean{M8.isingVec_zeroMode_twisted_annihilated},
  \chtwentylean{M8.isingVec_rigid} & spectral lemma; rigidity & normalised $e$; std. & Thm.~\ref{ch20:thm:rigid}\\
\chtwentylean{M8.aut_fix_quad_minimal_of_fix_ising}, \chtwentylean{M8.aut_fix_untwisted_of_fix_quad_minimal},
  \chtwentylean{M8.aut_eq_one_of_fix_untwisted_twistedGen} & Ising vectors determine $g$ (with $\sigma$) & displayed; std. &
  Lem.~\ref{ch20:lem:fixising}\\
\chtwentylean{M8.autVnat_isOfFinOrder} & \textbf{every automorphism of $\Vnat$ has finite order} & \textbf{none}; std. &
  Thm.~\ref{ch20:thm:finorder}\\
\chtwentylean{M8.finite_setOf_rigid_idempotent}, \chtwentylean{M8.derivation_eq_zero_of_span_rigid} & finitely many rigid
  idempotents; derivations vanish if they span & displayed; std. & Thm.~\ref{ch20:thm:finrigid}\\
\chtwentylean{M8.vnatRigidSet}, \chtwentylean{M8.eq_one_of_fix_rigid}, \chtwentylean{M8.finite_autVnat} & \textbf{$\Aut(\Vnat)$
  is finite} & \textbf{none}; std. & Thm.~\ref{ch20:thm:finite}\\
\chtwentylean{M8.BorcherdsPreData}, \chtwentylean{M8.BorcherdsPreData.KNZ} & pre-data; the KNZ product identity & definitions &
  Def.~\ref{ch20:def:predata}\\
\chtwentylean{M8.BorcherdsPreData.lieSpan_eq_top_of_KNZ} & KNZ $\Rightarrow$ generation & \chtwentylean{hK}; std. &
  Thm.~\ref{ch20:thm:genKNZ}\\
\chtwentylean{M8.monsterLieSimpleRootGeneration} & \textbf{$\chtwentym$ is generated by its simple root spaces} & \textbf{none};
  std. & Thm.~\ref{ch20:thm:generation}\\
\chtwentylean{M8.conwayNorton_of_finOrder_of_completelyReplicable} & CN from finite order and replicability &
  \chtwentylean{hfin}, \chtwentylean{hrep}; std. & Thm.~\ref{ch20:thm:cnrep}\\
\chtwentylean{M8.conwayNorton_of_isOfFinOrder} & CN from finite order & \chtwentylean{hfin}; std. & Thm.~\ref{ch20:thm:final}\\
\chtwentylean{M8.conwayNortonConjecture_holds} ($=$ \chtwentylean{ConwayNorton.conwayNorton}) & \textbf{Conway--Norton
  (genus-zero form) for every automorphism of $\Vnat$} & \textbf{none}; std. & Thm.~\ref{ch20:thm:final}; Ch.~\ref{ch:16}\\
\bottomrule
\end{longtable}
\end{center}

%======================================================================
\section{Notes and references}\label{ch20:sec:notes}
%======================================================================

\paragraph{The moonshine module and its automorphisms.}  I.~B.~Frenkel, J.~Lepowsky and A.~Meurman, \emph{Vertex Operator
Algebras and the Monster}, Pure and Applied Mathematics 134, Academic Press, 1988: the construction of $\Vnat$, the triality
$\sigma$, the group $C=2^{1+24}.\mathrm{Co}_1$, and the proof that $\Aut(\Vnat)$ is the Monster (via Griess' construction
and Tits' theorem on $\Aut$ of the Griess algebra).  T.~Gannon, \emph{Moonshine beyond the Monster}, Cambridge University
Press, 2006, \S5.2 for the definition of automorphisms and the Griess algebra used in the library.  R.~L.~Griess, \emph{The
friendly giant}, Invent.\ Math.\ 69 (1982) 1--102.

\paragraph{Ising vectors.}  C.~Dong, G.~Mason and Y.~Zhu, \emph{Discrete series of the Virasoro algebra and the moonshine
module}, in Algebraic Groups and Their Generalizations, Proc.\ Sympos.\ Pure Math.\ 56, Part 2, AMS, 1994, 295--316: the
vectors $\tfrac1{16}\alpha(-1)^2\chtwentyvac\pm\tfrac14(e^\alpha+e^{-\alpha})$ and the decomposition of $\Vnat$ under $48$ commuting Ising
models.  M.~Miyamoto, \emph{Griess algebras and conformal vectors in vertex operator algebras}, J.\ Algebra 179 (1996)
523--548: conformal vectors of central charge $\tfrac12$, the eigenvalues $\{0,\tfrac12,\tfrac1{16}\}$, and the involutions
$\tau_e$.  C.~Dong and R.~L.~Griess, \emph{Rank one lattice type vertex operator algebras and their automorphism groups},
J.\ Algebra 208 (1998) 262--275 (the idea that $\Aut$ is algebraic with Lie algebra $\mathrm{Der}$).  C.~Dong, R.~L.~Griess and
C.~H.~Lam, \emph{Uniqueness results for the moonshine vertex operator algebra}, Amer.\ J.\ Math.\ 129 (2007) 583--609
(Ising frames).  These are design sources only; no theorem of theirs is an input of the library.

\paragraph{Borcherds' proof.}  R.~E.~Borcherds, \emph{Monstrous moonshine and monstrous Lie superalgebras}, Invent.\ Math.\ 109
(1992) 405--444: the Monster Lie algebra, its simple roots and the twisted denominator identities.  J.~H.~Conway and
S.~P.~Norton, \emph{Monstrous moonshine}, Bull.\ London Math.\ Soc.\ 11 (1979) 308--339.  The product formula for $j(p)-j(q)$
was found independently by M.~Koike, S.~P.~Norton and D.~Zagier in the 1980s; see Norton, \emph{More on moonshine}, in
Computational Group Theory (Durham, 1982), Academic Press, 1984, 185--193.  V.~G.~Kac, \emph{Infinite Dimensional Lie
Algebras}, 3rd ed., Cambridge University Press, 1990, for Kac--Moody and Borcherds algebras and their denominator identities.
P.~Goddard and C.~B.~Thorn, Phys.\ Lett.\ B 40 (1972) 235--238, for the no-ghost theorem.

\paragraph{Algebra used in the finite-dimensional lemmas.}  Baer's criterion and the injectivity of divisible abelian groups:
S.~Lang, \emph{Algebra}, 3rd ed., Springer GTM 211, 2002, Ch.~XX.  Minimal primes, Nakayama's lemma and the Krull intersection
theorem: M.~F.~Atiyah and I.~G.~Macdonald, \emph{Introduction to Commutative Algebra}, Addison--Wesley, 1969, Ch.~2, 4, 10.

\paragraph{The working notes.}  The development was recorded in the authors' working notes, which are not part of the
release.  The \texttt{Aut*} files were written by Fedor Popov's side; the \texttt{MoonshineLR*} files of this chapter by
Leonardo Rastelli's side, with machine assistance.  These notes are the authors' claims; where this
chapter states what the machine confirmed, it relies on the build and audit records cited.

\endgroup
